%=======================================================================
\chapter{Observers and Frames of Reference}
%=======================================================================

Every formula written so far belongs to one observer. Chapter 2 gave that
observer a point space and let him record where each material point is;
Chapters 3 and 4 let him measure how the body has been distorted and how
fast; Chapters 5 and 6 let him balance mass, momentum and energy. At no
stage was it asked whether a second observer, watching the same body from
a platform of his own, would write the same things down.

He would not. Two observers disagree about the velocity of a particle,
about its acceleration, about the spin of the material, and about the
components of every vector and tensor in sight. They nevertheless agree
about the distance between two material points, about the elapsed time
between two events, about the density, and about the stress power. The
business of this chapter is to sort the one list from the other, and to
do it by calculation rather than by intuition, because intuition gets it
wrong: the stretching tensor $\bD$ survives a change of observer and the
spin tensor $\bW$ does not, although the two were built side by side out
of the same $\bL$ in \eqr{4.7}.

The sorting matters for a reason that will only become visible in
Chapter 8. A constitutive equation is a claim about a material, and a
material cannot know who is looking at it. So a constitutive equation may
be built only out of quantities on which all observers agree, in the
precise sense set out in Section~\ref{sec:ch7kin} below. Every result of
this chapter is therefore a licence or a prohibition for the chapters
that follow: $\bD$ is admissible in a constitutive law and $\bW$ is not,
$\bC$ is admissible and $\bF$ is not, and the reason in each case is a
transformation law derived here.

One object generates the entire chapter, the relation
\[
   \bx^{+}=\bQ(t)\bx+\bc(t)
\]
between the place $\bx$ that one observer assigns to a material point and
the place $\bx^{+}$ that the other assigns to the same point at the same
instant. Section~\ref{sec:ch7obs} derives it from two physical
assumptions. Everything afterwards is obtained by differentiating it,
with respect to position to get $\bF$, $\bH$ and the polar factors, and
with respect to time to get $\bv$, $\ba$, $\bL$, $\bD$ and $\bW$. The
single term $\dot\bQ$, which appears whenever the second differentiation
is performed, is the whole difficulty and the whole content.

%=======================================================================
\section{Observers}\label{sec:ch7obs}
%=======================================================================

An observer $O$, perceiving a configuration $\kappa_{t}$ of a body, is
equipped for the purpose with a point space $\Et$ and its associated
translation space, also written $\Et$. These constitute the \emph{frame
of reference} adopted by $O$. We are concerned with the relationship
between this observer's perception and that of an alternative observer
$O^{+}$, who perceives the same configuration as $\kappa^{+}_{t^{+}}$
relative to a frame of reference consisting of the point space
$\mathbb{E}^{3+}$ and translation space $\mathbb{E}^{3+}$.

\begin{remark}[a frame of reference is not a coordinate
system]\label{rem:ch7frame}
This is the distinction that has to be fixed before anything else, and it
is the one most often lost.

A \emph{basis} is chosen inside a single frame. Changing it, as in
Section~\ref{ss:cob}, leaves every geometric object exactly where it is
and merely re-expresses it: the vector $\bu$ is one vector, and
$u_{i}=\ip\bu{\be_{i}}$ and $u_{i}'=\ip\bu{\be_{i}'}$ are two lists of
numbers naming it. The change-of-basis matrix $a_{ij}$ of
Theorem~\ref{thm:cob} is a fixed array of numbers; nothing in that
discussion had any time dependence, and nothing could, because the
question ``which basis?'' is not a question about when.

A \emph{frame} is a whole space, carried by an observer. Changing it
replaces the point $\bx\in\Et$ by a different point $\bx^{+}$ of a
different space $\mathbb{E}^{3+}$. The tensor $\bQ(t)$ that relates them
is a two-point tensor with one foot in each space, in exactly the sense
of Remark~\ref{rem:hatinch4}, and it depends on $t$ because the two
observers may be in relative motion.

The practical test: if $\dot\bQ$ can be non-zero, you are changing
observers; if it cannot, you are changing bases. Every difficulty in this
chapter, and every term in \eqr{7.9} and \eqr{7.63}, comes from
$\dot\bQ\neq\Ob$. A change of basis produces no such term and no such
difficulty, which is why Chapters 1 to 6 could ignore the whole subject.
\end{remark}

In classical non-relativistic physics, including classical continuum
mechanics, all observers are presumed to agree on two things: the time
intervals between successive events, and the distance between arbitrary
pairs of material points. Thus, if $t$ and $t^{+}$ are the times on the
watches worn by $O$ and $O^{+}$ respectively, then
\begin{equation}
   t^{+}=t+a ,
   \tag{7.1}\label{eq:7.1}
\end{equation}
where $a$ is a constant; and if $\bx,\by\in\Et$ and
$\bx^{+},\by^{+}\in\mathbb{E}^{3+}$ are the positions of material points
$p$ and $q$ as perceived by $O$ and $O^{+}$ (Figure~\ref{fig:ch7two}),
then
\begin{equation}
   \big|\bx^{+}-\by^{+}\big|=\big|\bx-\by\big| .
   \tag{7.2}\label{eq:7.2}
\end{equation}

\begin{remark}[what the two assumptions cost, and what they
buy]\label{rem:ch7twoassump}
Neither statement is a theorem. Both are the classical idealisation, and
both fail in relativity: \eqr{7.1} fails because simultaneity is not
absolute, and \eqr{7.2} fails because lengths contract. Writing them down
is therefore the point at which this book commits to Newtonian physics,
and it is worth knowing exactly how much is being committed.

Very little, and very much. Very little, because the two statements
concern only clocks and rulers and say nothing about matter. Very much,
because \eqr{7.2} alone will be shown to force the form of the change of
observer completely, leaving only the arbitrary $\bQ(t)$ and $\bc(t)$; and
the constant $a$ in \eqr{7.1} will be shown to be irrelevant, since only
$\dl t^{+}/\dl t=1$ is ever used. From two sentences about measurement
comes the entire kinematics of the chapter.
\end{remark}

\bookfig{1}\begin{figure}[htbp]\centering
\begin{tikzpicture}[scale=1.0,
   ar/.style={-{Stealth[length=2.3mm]},thick}]
% ================= as O sees it =================
\begin{scope}
  \blob{configline}
  \coordinate (oo) at (-2.35,-1.75);
  \coordinate (X) at (-0.40,0.30);
  \coordinate (Y) at (0.80,-0.42);
  \draw[thinvec,gray!70] (oo)--(X);
  \draw[thinvec,gray!70] (oo)--(Y);
  \draw[very thick,blue!55!black] (X)--(Y);
  \fill (X) circle (1.3pt); \fill (Y) circle (1.3pt);
  \fill (oo) circle (1.3pt);
  \node[lbl,anchor=south east] at (X) {$\bx$};
  \node[lbl,anchor=west] at (Y) {$\by$};
  \node[lbl,anchor=north east] at (oo) {$\bo$};
  \node[lbl,font=\small,text=blue!55!black] at (0.28,0.10) {$\ell$};
  \node[font=\small,anchor=north] at (0,-1.95) {$\kappa_{t}\subset\Et$};
  \node[font=\footnotesize,text=gray!85,anchor=north] at (0,-2.38)
     {observer $O$, clock reads $t$};
\end{scope}
% ================= the map =================
\draw[ar] (2.05,0.35) .. controls (3.05,0.95) and (4.05,0.95) .. (5.05,0.35);
\node[font=\small,anchor=south] at (3.55,0.92) {$\bQ(t),\ \bc(t)$};
% ================= as O+ sees it =================
\begin{scope}[xshift=7.2cm,rotate=-32]
  \blob{configline}
  \coordinate (X2) at (-0.40,0.30);
  \coordinate (Y2) at (0.80,-0.42);
  \draw[very thick,blue!55!black] (X2)--(Y2);
  \fill (X2) circle (1.3pt); \fill (Y2) circle (1.3pt);
  \node[lbl,anchor=south east] at (X2) {$\bx^{+}$};
  \node[lbl,anchor=west] at (Y2) {$\by^{+}$};
  \node[lbl,font=\small,text=blue!55!black] at (0.30,0.12) {$\ell$};
\end{scope}
\begin{scope}[xshift=7.2cm]
  \coordinate (oo2) at (-2.05,-1.95);
  \fill (oo2) circle (1.3pt);
  \node[lbl,anchor=north east] at (oo2) {$\bo^{+}$};
  \draw[thinvec,gray!70] (oo2)--(-0.05,0.72);
  \draw[thinvec,gray!70] (oo2)--(1.12,0.04);
  \node[font=\small,anchor=north] at (0,-1.95)
     {$\kappa^{+}_{t^{+}}\subset\mathbb{E}^{3+}$};
  \node[font=\footnotesize,text=gray!85,anchor=north] at (0,-2.38)
     {observer $O^{+}$, clock reads $t^{+}=t+a$};
\end{scope}
\end{tikzpicture}
\caption{A body as perceived by two observers. The same two material
points are seen at $\bx,\by$ by $O$ and at $\bx^{+},\by^{+}$ by $O^{+}$,
and the segment joining them has the same length $\ell$ in both pictures:
that is \eqr{7.2}, and it is the only geometric agreement assumed. The
two observers need not agree about which direction the segment points,
about where the origin is, or about anything else. Section~\ref{sec:ch7obs}
shows that this single agreement forces the relation between the pictures
to be \eqr{7.3}, a uniform orthogonal $\bQ(t)$ followed by a translation
$\bc(t)$, and that both may depend on time in any way whatever.}
\label{fig:ch7two}
\end{figure}

Assume the two point spaces are related by a differentiable mapping. Fix
$t$, and let $\phi$ denote the map that sends the place $\bx\in\Et$
assigned by $O$ to a material point to the place
$\bx^{+}=\phi(\bx)\in\mathbb{E}^{3+}$ assigned by $O^{+}$ to the same
material point at the same instant. Condition \eqr{7.2} reads
\[
   \big|\phi(\by)-\phi(\bx)\big|=\big|\by-\bx\big|
   \qquad\text{for all }\bx,\by ,
\]
which is word for word the hypothesis \eqr{3.113} of the analysis of
rigid-body motion in Section~\ref{sec:rigid}, with $\phi$ in place of
$\bchi_{\kappa}(\,\cdot\,,t)$. The argument given there therefore applies
verbatim. It is short enough to be worth repeating in the present
notation, because the conclusion is the pivot of the chapter.

Square both sides to remove the square root, so that each side is a
polynomial and differentiates cleanly:
\[
   \ip{\big[\phi(\by)-\phi(\bx)\big]}{\big[\phi(\by)-\phi(\bx)\big]}
   =\ip{(\by-\bx)}{(\by-\bx)} .
\]
Fix $\bx$ and differentiate with respect to $\by$. Both sides have the
form $\ip\ba\ba$, whose differential is
$\ip{\dl\ba}\ba+\ip\ba{\dl\ba}=2\ip\ba{\dl\ba}$, the two terms being
equal by the symmetry of the inner product; the factor $2$ appears on
both sides and cancels. Writing $\bG(\by)$ for the gradient of $\phi$ at
$\by$, a two-point tensor $\Et\to\mathbb{E}^{3+}$, so that
$\dl\phi=\bG\,\dl\by$,
\[
   \ip{\big[\phi(\by)-\phi(\bx)\big]}{\bG(\by)\,\dl\by}
   =\ip{(\by-\bx)}{\dl\by} .
\]
Move $\bG(\by)$ to the other slot by the definition of the transpose,
$\ip\ba{\bG\bb}=\ip{\bG^{t}\ba}\bb$, and collect:
\[
   \ip{\Big(\bG^{t}(\by)\big[\phi(\by)-\phi(\bx)\big]-(\by-\bx)\Big)}
      {\dl\by}=0\qquad\text{for all }\dl\by .
\]
A vector orthogonal to every vector is $\bze$ by Lemma~\ref{lem:cvec}, so
\[
   \bG^{t}(\by)\big[\phi(\by)-\phi(\bx)\big]=\by-\bx .
\]
Now reverse the roles: fix $\by$, so that $\bG^{t}(\by)$ and $\phi(\by)$
are both constant, and differentiate with respect to $\bx$. The only
$\bx$-dependent piece on the left is $-\phi(\bx)$, whose differential is
$-\bG(\bx)\,\dl\bx$; on the right it is $-\bx$. Hence
\[
   -\bG^{t}(\by)\bG(\bx)\,\dl\bx=-\dl\bx\qquad\text{for all }\dl\bx ,
\]
and since two tensors agreeing on every vector are equal,
\[
   \bG^{t}(\by)\bG(\bx)=\I\qquad\text{for all }\bx,\by\in\Et ,
\]
$\I$ being the identity of $\Et$. Setting $\by=\bx$ gives
$\bG^{t}(\bx)\bG(\bx)=\I$, so $\bG(\bx)$ is orthogonal at every $\bx$ and
$\bG^{t}=\bG^{-1}$ by Proposition~\ref{prop:orth}; substituting that back
gives $\bG^{-1}(\by)\bG(\bx)=\I$, i.e. $\bG(\by)=\bG(\bx)$ for every pair
of places. So $\bG$ does not depend on position at all. Call its common
value $\bQ(t)$, the dependence on $t$ surviving because $t$ was held
fixed throughout. Finally $\dl\bx^{+}=\bQ(t)\,\dl\bx=\dl(\bQ(t)\bx)$,
$\bQ$ being uniform, so $\dl(\bx^{+}-\bQ\bx)=\bze$ and the bracket is
independent of $\bx$; this is the integration step already used at
\eqr{3.99} and \eqr{4.30}. Therefore
\begin{equation}
   \bx^{+}=\bQ(t)\bx+\bc(t),
   \tag{7.3}\label{eq:7.3}
\end{equation}
where $\bc(t)\in\mathbb{E}^{3+}$ is arbitrary and $\bQ(t)$ is an
arbitrary spatially uniform two-point orthogonal tensor mapping $\Et$ to
$\mathbb{E}^{3+}$.

The converse is immediate and is what makes the condition sharp: if
\eqr{7.3} holds then
\[
   \big|\bx^{+}-\by^{+}\big|^{2}
   =\ip{\bQ(\bx-\by)}{\bQ(\bx-\by)}
   =\ip{(\bx-\by)}{\bQ^{t}\bQ(\bx-\by)}
   =\big|\bx-\by\big|^{2},
\]
so \eqr{7.2} is recovered. Thus \eqr{7.2} holds \emph{if and only if}
\eqr{7.3} does.

\begin{remark}[the shape of that proof, and what each differentiation
rules out]\label{rem:ch7isometry}
The derivation is four lines long but it is not obvious why it should
work at all, and the reason is worth extracting, because the same three
moves are used at \eqr{3.113} and in every rigidity argument in the book.

The hypothesis \eqr{7.2} is a statement about \emph{pairs} of points; the
conclusion \eqr{7.3} is a statement about \emph{one} point at a time. The
entire difficulty is the passage from the first kind of statement to the
second, and each step performs one part of it.

\emph{Squaring} is not cosmetic. The function $|\cdot|$ is not
differentiable where its argument vanishes, and $|\cdot|^{2}$ is a
polynomial in the components. Both sides being non-negative, nothing is
lost.

\emph{The first differentiation, in $\by$, makes the statement local.} It
replaces the map $\phi$ by its gradient and yields
$\bG^{t}(\by)\big[\phi(\by)-\phi(\bx)\big]=\by-\bx$. This is still a
statement about a pair of points, but it is now linear in the increment
at $\by$; and it says nothing yet about whether $\bG$ is the same at
different places, which is exactly the thing to be proved.

\emph{The second differentiation, in $\bx$, is what makes $\bG$
uniform.} Both $\phi(\by)$ and $\by$ are constants of that
differentiation and disappear, leaving $\bG^{t}(\by)\bG(\bx)=\I$ for
\emph{every pair} of places. That single relation carries two
conclusions: putting $\by=\bx$ gives orthogonality at each place, and
feeding that back gives $\bG(\by)=\bG(\bx)$ everywhere. The second is
the one that matters. A map that preserves all distances cannot bend,
because bending is precisely a $\bG$ that varies from place to place, and
two places whose $\bG$ differed would betray themselves in the distance
between a third pair.

\emph{The integration} turns a uniform gradient into an affine map, as
after \eqr{3.99} and \eqr{4.30}. It is here, and only here, that the
arbitrary $\bc(t)$ enters: an affine map is determined by its linear part
only up to a translation, which is the analytic form of the statement of
Remark~\ref{rem:ch7pointvec} that $\bc$ is arbitrary because the two
origins are.

Two footnotes on the hypotheses, since both are easy to overlook.
Differentiability of $\phi$ was assumed before the argument began; it is
a real assumption here, though a removable one, because a
distance-preserving map of one Euclidean space onto another is affine
with no smoothness assumed at all. Nothing in this book needs the sharper
statement. And what comes out is orthogonality, not rotation:
$\bQ^{t}\bQ=\I$ gives $\det\bQ=\pm1$ by
Proposition~\ref{prop:orth}, and no step of the argument prefers one sign
to the other. Remark~\ref{rem:ch7mirror} takes up what the other sign
would mean.
\end{remark}

Because tensors operate on vectors rather than on positions, \eqr{7.3}
should be interpreted to mean
\begin{equation}
   \bx^{+}-\by^{+}=\bQ(t)(\bx-\by).
   \tag{7.4}\label{eq:7.4}
\end{equation}
In fact, as suggested by Figure~\ref{fig:ch7two}, \eqr{7.3} also holds if
the positions therein are replaced by position vectors.

\begin{remark}[why \eqr{7.4} is the honest statement and \eqr{7.3} the
usable one]\label{rem:ch7pointvec}
A point is not a vector. The difference of two points is a vector, by
Definition~\ref{def:eps}, and it is on differences that $\bQ$ acts. That
is all \eqr{7.4} says, and it is the only form of the statement that is
correct before any origin has been named.

Once each observer names an origin, $\bo\in\Et$ for $O$ and
$\bo^{+}\in\mathbb{E}^{3+}$ for $O^{+}$, the deliberate abuse of
Remark~\ref{rem:abuse} lets the same letter denote a point and its
position vector $\bx-\bo$. Apply \eqr{7.4} to the pair $(\bx,\bo)$, which
gives $\phi(\bx)-\phi(\bo)=\bQ(t)(\bx-\bo)$, and insert the place
$\phi(\bo)$ between $\bx^{+}=\phi(\bx)$ and $\bo^{+}$ by Chasles'
relation, Definition~\ref{def:eps}(P2):
\[
\begin{aligned}
   \bx^{+}-\bo^{+}&=\big(\phi(\bx)-\phi(\bo)\big)
   +\big(\phi(\bo)-\bo^{+}\big)
   \\
   &=\bQ(t)(\bx-\bo)+\bc(t),\qquad
   \bc(t)=\phi(\bo)-\bo^{+},
\end{aligned}
\]
which is \eqr{7.3} verbatim for position vectors, with $\bc(t)$ now
identified as the position, in $O^{+}$'s frame, of the place $O^{+}$
assigns to $O$'s origin. So nothing is lost by reading \eqr{7.3} as an
equation between position vectors, and that is how it will be read from
here on. What is lost if the distinction is forgotten altogether is the
ability to see why $\bc$ is arbitrary: it is arbitrary precisely because
the choice of the two origins is.
\end{remark}

\begin{remark}[the two freedoms in \eqr{7.3}, and which one is
dangerous]\label{rem:ch7twofreedoms}
The change of observer carries two free functions of time, and they are
not on the same footing.

$\bc(t)$ is a relative translation. It disappears from every statement
about \emph{differences} of positions, and \eqr{7.4} is the proof: no
$\bc$ appears there. Consequently $\bc$ can never affect a deformation
gradient, a strain, or a stress, all of which are built from differences.
It survives only in the velocity, through $\dot\bc$, and in the
acceleration, through $\ddot\bc$.

$\bQ(t)$ is a relative rotation, and it does not disappear from
differences. It is the source of every non-trivial transformation law in
this chapter. The truly dangerous part of it is $\dot\bQ$: if $\bQ$ were
constant the whole chapter would collapse into the change-of-basis
formulas of Chapter 1, because a constant $\bQ$ commutes with $\dl/\dl t$
and every rate would transform exactly as the quantity itself does. It is
because $\bQ$ turns as time passes, and hence fails to commute with
differentiation, that $\bD$ and $\bW$ part company at \eqr{7.34}, that
$\ba$ acquires the extra terms of \eqr{7.9}, and that the equation of
motion acquires the inertial force of \eqr{7.63}.
\end{remark}

\begin{remark}[may one observer be the mirror image of
another?]\label{rem:ch7mirror}
Nothing said so far forbids it, and it is worth seeing that this is not
an oversight.

\emph{He is a legitimate observer.} A reflection preserves every
distance, so it satisfies \eqr{7.2} exactly, and the derivation of
\eqr{7.3} returned only $\bQ^{t}\bQ=\I$, hence $\det\bQ=\pm1$ by
Proposition~\ref{prop:orth}. Had the book wanted to exclude mirror
observers it would have had to say so, and it would have had to justify
saying so by something other than clocks and rulers, because clocks and
rulers cannot tell left from right.

\emph{What is unaffected.} Every transformation law derived in this
chapter. Each was obtained from $\bQ^{t}\bQ=\I$ alone, and none of them
used $\det\bQ=1$: $\bU^{+}=\bK\bU\bK^{t}$, $\bT^{+}=\bQ\bT\bQ^{t}$ and
the rest hold verbatim for a mirror observer. So does everything metric,
lengths and angles being exactly what \eqr{7.2} protects.

\emph{What he does not agree about.} Handedness. A right-handed helix is
left-handed in his picture (Figure~\ref{fig:ch7mirror}); a right-handed
basis of $\Et$ has a left-handed image in $\mathbb{E}^{3+}$; and every
cross product changes sign relative to what a rotation would have given,
because the box product carries the factor $\det\bQ$ by
Theorem~\ref{thm:det}. The consequence worth recording is that an axial
vector is not a spatial vector in the sense of
Definition~\ref{def:ch7obj}. If $\bA$ is skew with axial vector $\ba$ in
the sense of Proposition~\ref{prop:axial}, and if $\bA$ is objective,
$\bA^{+}=\bQ\bA\bQ^{t}$, then for every $\bs\in\mathbb{E}^{3+}$
\[
\begin{aligned}
   \bA^{+}\bs&=\bQ\bA\bQ^{t}\bs=\bQ\big(\ba\times\bQ^{t}\bs\big)\\
   &=(\det\bQ)\,(\bQ\ba)\times\big(\bQ\bQ^{t}\bs\big)
   =(\det\bQ)(\bQ\ba)\times\bs ,
\end{aligned}
\]
using $\bQ(\bu\times\bw)=(\det\bQ)(\bQ\bu)\times(\bQ\bw)$, which is
Theorem~\ref{thm:det} written for an orthogonal $\bQ$. Hence
\begin{equation*}
   \ba^{+}=(\det\bQ)\,\bQ\ba ,
\end{equation*}
which is $\bQ\ba$ for a rotation and $-\bQ\ba$ for a reflection. This is
why the whole of this chapter is written in the skew tensor $\bOm$ and
never in its axial vector $\bom$: the tensor transforms by one rule for
both signs and the vector does not. The same caution applies to $\bW$,
and it is the reason \eqr{7.34}$_{2}$ is stated as a tensor equation.

\emph{What \eqr{7.19} costs him.} Having chosen a mirrored frame,
$O^{+}$ is no longer free in his choice of reference configuration. The
requirement that his reference configuration be occupiable is $J^{+}>0$,
and that forced $\det\bQ\det\bK=1$; so $\det\bQ=-1$ obliges
$\det\bK=-1$ as well. A mirror observer must mirror his reference copy of
the body too. This is not a restriction on physics but a consistency
condition on bookkeeping, and it is the only place in the chapter where
the sign of a determinant does any work.

\emph{Why the question is not idle.} Whether a material can tell itself
from its mirror image is a genuine physical question, and for some
materials the answer is yes: a quartz crystal, a cholesteric liquid
crystal, a composite whose fibres wind one way. Chapter 8 will require
constitutive equations to be indifferent to a change of observer, and by
this remark that requirement may be posed over the full orthogonal group
$\Orth$ or over the proper orthogonal group alone. The two demands are
different, and the difference is exactly whether the material is
permitted to be chiral. Nothing in the present chapter decides between
them, because \eqr{7.2} does not.
\end{remark}

\begin{figure}[htbp]\centering
\renewcommand{\thefigure}{A}
\begin{tikzpicture}[scale=1.0,
   ar/.style={-{Stealth[length=2.3mm]},thick},
   sm/.style={-{Stealth[length=1.9mm]},semithick}]
% ---- the helix, drawn as front (solid) and back (pale) strands -------
% r(phi) = (a phi, r cos phi, r sin phi); depth = r sin phi
\def\coil{%
  \draw[guide] (-0.25,0)--(3.85,0);
  \draw[ar,gray!70] (3.85,0)--(4.30,0);
  \draw[thick,gray!45]
     plot[domain=180:360,samples=28,smooth]
     ({0.0042*\x},{0.42*cos(\x)});
  \draw[thick,gray!45]
     plot[domain=540:720,samples=28,smooth]
     ({0.0042*\x},{0.42*cos(\x)});
  \draw[thick,blue!55!black]
     plot[domain=0:180,samples=28,smooth]
     ({0.0042*\x},{0.42*cos(\x)});
  \draw[thick,blue!55!black]
     plot[domain=360:540,samples=28,smooth]
     ({0.0042*\x},{0.42*cos(\x)});
  \draw[thick,blue!55!black]
     plot[domain=720:900,samples=28,smooth]
     ({0.0042*\x},{0.42*cos(\x)});
  \draw[sm,blue!55!black] (0.30,0.29)--(0.44,-0.26);}
% ================= (a) as O sees it =================
\begin{scope}
  \coil
\end{scope}
\node[font=\small,align=center,anchor=north] at (2.00,-0.95)
   {(a) observer $O$\\[0.5mm] a right-handed helix};
% ================= the mirror =================
\draw[thick,gray!45,densely dashed] (5.15,-0.80)--(5.15,0.88);
\node[font=\footnotesize,text=gray!85,anchor=south] at (5.15,0.90)
   {mirror};
% ================= (b) as the mirror observer sees it =================
\begin{scope}[xshift=10.30cm,xscale=-1]
  \coil
\end{scope}
\node[font=\small,align=center,anchor=north] at (8.30,-0.95)
   {(b) observer $O^{+}$, $\det\bQ=-1$\\[0.5mm]
    the same helix, left-handed};
\end{tikzpicture}
\caption{A mirror observer. The solid strands are in front and the pale
ones behind, and the grey arrow is the direction of advance along the
axis, so the two panels are genuinely different curves in space and not
the same curve redrawn: no rotation carries one onto the other.
Yet every distance between two points of the coil is the same in both,
which is all \eqr{7.2} demands, so $\bQ$ with $\det\bQ=-1$ is an
admissible change of observer. What the two observers disagree about is
handedness, and with it the sign of every cross product and of every
axial vector, by the rule $\ba^{+}=(\det\bQ)\bQ\ba$ established in
Remark~\ref{rem:ch7mirror}. The disagreement is invisible in $\bU$, $\bC$
and $\bD$, which are built from lengths and angles alone; it becomes
visible the moment a formula contains $\bom$ rather than $\bOm$, or asks
whether a material is chiral.}
\label{fig:ch7mirror}
\end{figure}

We can now name the two properties that the rest of the chapter is
occupied with sorting.

\begin{definition}[objective and invariant]\label{def:ch7obj}
A field carried by the body is \emph{objective}, or \emph{frame
indifferent}, if the values recorded by the two observers are related as
the type of the field demands, that is, if
\[
\begin{aligned}
   \text{a scalar }\varphi:&\qquad \varphi^{+}=\varphi ,\\
   \text{a spatial vector }\bu:&\qquad \bu^{+}=\bQ\bu ,\\
   \text{a spatial tensor }\bA:&\qquad \bA^{+}=\bQ\bA\bQ^{t} .
\end{aligned}
\]
A quantity is \emph{invariant} under the change of observer if both
observers assign it the same value, $(\,\cdot\,)^{+}=(\,\cdot\,)$. An
\emph{equation} is invariant, or form invariant, if it has the same form
in the two frames, that is, if starring every symbol in it produces a
true statement.
\end{definition}

\begin{remark}[the distinction, and a counterexample on each
side]\label{rem:ch7objinv}
Objectivity is a property of a \emph{quantity}; invariance is a property
of a \emph{quantity's value} or of an \emph{equation}. Confusing the two
is the standard error in this subject, so it is worth having a
counterexample in each direction ready before the calculations start. All
four are proved in the sections that follow; they are collected here as
signposts.

\emph{Objective but not invariant.} The Cauchy stress. By \eqr{7.40},
$\bT^{+}=\bQ\bT\bQ^{t}$, which is exactly what
Definition~\ref{def:ch7obj} requires of a spatial tensor, so $\bT$ is
objective. But $\bQ\bT\bQ^{t}\neq\bT$ in general: the two observers write
down different tensors, with different components. Objectivity is not
sameness; it is agreement about \emph{how} to differ.

\emph{Invariant but not objective, in the sense that the test does not
apply.} The right Cauchy--Green tensor. By \eqr{7.30}$_{1}$,
$\bC^{+}=\bK\bC\bK^{t}$ with $\bK$ a \emph{fixed} tensor that has nothing
to do with the relative motion of the observers; if the two observers
happen to use a common reference configuration, $\bK=\hI$ and
$\bC^{+}=\bC$ literally. So $\bC$ carries no trace of $\bQ$ at all. It is
not objective in the sense of Definition~\ref{def:ch7obj} because it is
not a spatial tensor and the third line of that definition is not the
right test for it. Set beside $\bB^{+}=\bQ\bB\bQ^{t}$ of
\eqr{7.30}$_{2}$, this is instructive: $\bC$ and $\bB$ encode the same
deformation and are as close to being the same tensor as two tensors can
be, and yet one of them ignores the observer completely while the other
turns with him. The difference is only which space they live in.

\emph{An invariant scalar built from non-objective ingredients.} The
velocity is not objective, by \eqr{7.8}: $\bv^{+}\neq\bQ\bv$. Yet
$\dvg^{+}\bv^{+}=\dvg\bv$ by \eqr{7.56}. A non-objective field can have
an invariant scalar hidden inside it, and the extra terms can cancel.

\emph{An invariant scalar built from one objective and one
non-objective ingredient.} $\bT$ is objective and $\bL$ is not, by
\eqr{7.33}; nevertheless the stress power density is invariant,
$\bT^{+}\cdot\bL^{+}=\bT\cdot\bL$, by \eqr{7.58}. The offending part of
$\bL$ is annihilated by the symmetry of $\bT$.

The moral is that invariance must be checked, never assumed from the
ingredients, and that is what Section~\ref{sec:ch7balance} does law by
law. The one law that fails the check, linear momentum, fails it in a way
that is itself instructive.
\end{remark}

%=======================================================================
\section{Transformation of the motion}\label{sec:ch7motion}
%=======================================================================

Following the development of Section~\ref{sec:pva}, the motions of
$p\in B$, as perceived by $O$ and $O^{+}$, are expressible as
\begin{equation}
   \bx=\bchi(p,t)
   \qquad\text{and}\qquad
   \bx^{+}=\bchi^{+}(p,t^{+}),
   \tag{7.5}\label{eq:7.5}
\end{equation}
respectively. Accordingly, substituting both into \eqr{7.3},
\begin{equation}
   \bchi^{+}(p,t^{+})=\bQ(t)\bchi(p,t)+\bc(t).
   \tag{7.6}\label{eq:7.6}
\end{equation}

\begin{remark}[why the argument is $p$ and not $\bX$]\label{rem:ch7whyp}
Equations \eqr{7.5} take the material point itself as the independent
variable, which is the material description of
Definition~\ref{def:motion} and Remark~\ref{rem:decorations}, and they
have to. The body $B$ is the one thing the two observers share. It is not
a subset of $\Et$, nor of $\mathbb{E}^{3+}$; it is a set of material
points, and each observer supplies his own map of it, as
Remark~\ref{rem:fig123} insisted. Had we written $\bchi_{\kappa}(\bX,t)$
here, we would have had to say which observer chose $\kappa$, and the
statement \eqr{7.6} would have been circular. Section~\ref{sec:ch7kin}
introduces reference configurations only after \eqr{7.6} has been secured
without them, and pays for the convenience with the extra tensor $\bK$.
\end{remark}

The velocity of $p$ in the frame $\mathbb{E}^{3+}$ is
$\bv^{+}=\partial\bchi^{+}/\partial t^{+}$, by
Definition~\ref{def:va} written in that frame. The two time derivatives
are the same operation: $t^{+}=t+a$ with $a$ constant by \eqr{7.1}, so
$\dl t^{+}/\dl t=1$ and the chain rule gives
$\partial/\partial t^{+}=(\dl t/\dl t^{+})\,\partial/\partial t
=\partial/\partial t$, both taken at fixed $p$. Differentiating
\eqr{7.6} at fixed $p$, with the product rule applied to $\bQ(t)\bchi$
and remembering that $\bQ$ and $\bc$ depend on $t$ alone while $\bchi$
depends on $p$ and $t$,
\begin{equation}
   \bv^{+}=\frac{\partial}{\partial t^{+}}\bchi^{+}
   =\frac{\partial}{\partial t}\bchi^{+}
   =\dot\bQ\bx+\bQ\bv+\dot\bc ,
   \tag{7.7}\label{eq:7.7}
\end{equation}
where $\bv=\partial\bchi/\partial t$ is the velocity of $p$ in the frame
$\Et$ and $\bx=\bchi(p,t)$.

This is the first transformation law of the chapter that is not simply a
conjugation, and the reason is visible in the algebra: of the three terms
on the right, only the middle one came from differentiating the motion.
The other two came from differentiating the relation between the
observers. If the body stands still for $O$, so that $\bv=\bze$, then
$\bv^{+}=\dot\bQ\bx+\dot\bc$, which vanishes only if the observers are at
relative rest. \emph{The velocity is not objective.}

It is convenient to eliminate $\bx$, which belongs to $O$'s frame, in
favour of $\bx^{+}$. From \eqr{7.3}, $\bQ\bx=\bx^{+}-\bc$, so
\[
   \dot\bQ\bx=\dot\bQ\big(\bQ^{t}\bQ\big)\bx
   =\big(\dot\bQ\bQ^{t}\big)\big(\bQ\bx\big)
   =\bOm\,(\bx^{+}-\bc),
\]
where the identity $\bQ^{t}\bQ=\I$ of Proposition~\ref{prop:orth} was
inserted and the factors regrouped. Thus the velocities in the two frames
are related by
\begin{equation}
   \bv^{+}-\bQ\bv=\bOm(\bx^{+}-\bc)+\dot\bc ,
   \qquad\text{where}\qquad
   \bOm(t)=\dot\bQ\bQ^{t}
   \tag{7.8}\label{eq:7.8}
\end{equation}
is the spin of $\mathbb{E}^{3+}$.

Two properties of $\bOm$ are needed constantly and are established once
here.

\emph{$\bOm$ is skew.} Differentiate the orthogonality relation
$\bQ\bQ^{t}=\I^{+}$, in which $\I^{+}$ is the identity of
$\mathbb{E}^{3+}$ and is therefore constant. By the product rule and the
commutation of differentiation with transposition, Problem 1(a) and
Problem 1(b) of Chapter 4,
\[
   \dot\bQ\bQ^{t}+\bQ\dot{\bQ^{t}}
   =\dot\bQ\bQ^{t}+\big(\dot\bQ\bQ^{t}\big)^{t}
   =\Ob^{+},
\]
the middle step because $\bQ\dot{\bQ^{t}}=\bQ(\dot\bQ)^{t}
=(\dot\bQ\bQ^{t})^{t}$ by the rule $(\bA\bB)^{t}=\bB^{t}\bA^{t}$. Hence
$\bOm+\bOm^{t}=\Ob^{+}$, i.e. $\bOm^{t}=-\bOm$.

\emph{$\bOm$ is a spatial tensor of the frame of $O^{+}$.} The derivative
of a two-point tensor is a two-point tensor of the same type, so
$\dot\bQ\colon\Et\to\mathbb{E}^{3+}$, while
$\bQ^{t}\colon\mathbb{E}^{3+}\to\Et$. The composition
$\bOm=\dot\bQ\bQ^{t}$ therefore maps $\mathbb{E}^{3+}$ into itself. This
is why $\bOm$ may be added to $\bQ\bL\bQ^{t}$ in \eqr{7.33} without an
index mismatch, and it is worth checking by the alphabet test of
Remark~\ref{rem:hatinch4}: $\dot\bQ$ carries
$\be^{+}_{i}\otimes\be_{j}$, $\bQ^{t}$ carries
$\be_{j}\otimes\be^{+}_{k}$, the unstarred indices contract in pairs and
what is left is $\be^{+}_{i}\otimes\be^{+}_{k}$.

\begin{remark}[what $\dot\bQ\bQ^{t}$ \emph{is}]\label{rem:ch7omega}
The combination has been met before. Equation \eqr{4.29} says that in a
rigid motion $\bx=\bQ(t)\bX+\bc(t)$ the spatial velocity gradient is
$\bL=\dot\bQ\bQ^{t}$, that the stretching vanishes, $\bD=\Ob$, and that
the whole of $\bL$ is spin, $\bW=\dot\bQ\bQ^{t}$. Compare that motion
with \eqr{7.3} and they are the same formula. A change of observer is,
formally, a rigid motion: not a motion of the body, but of one frame
relative to the other.

So $\bOm$ is the velocity gradient of the rigid displacement that carries
$O$'s frame onto $O^{+}$'s, and it is a \emph{pure spin} for exactly the
reason \eqr{4.29} gives, namely that a rigid displacement stretches
nothing. By Proposition~\ref{prop:axial} the skew tensor $\bOm$ has a
unique axial vector $\bom$ with $\bOm\bs=\bom\times\bs$, and $\bom$ is
the angular velocity of $O$'s frame as recorded by $O^{+}$. If the
platform of $O^{+}$ turns with angular velocity $\bom_{p}$ relative to
$O$, then $\bom=-\bom_{p}$: what $O^{+}$ sees is the rest of the world
going round the other way. Problem 3 computes both for an observer
standing on a spinning planet.

Everything else in this chapter follows from this one identification. The
extra term in \eqr{7.33} is a pure spin, so it lands entirely in the skew
part of $\bL^{+}$ and leaves the symmetric part untouched: that is
\eqr{7.34}, and it is the whole reason $\bD$ is objective and $\bW$ is
not. The extra terms in \eqr{7.9} are what a pure spin does to an
acceleration, and they are the Coriolis and centrifugal terms.
\end{remark}

\begin{remark}[\eqr{7.8} read as a decomposition, and a one-line route to
\eqr{7.33}]\label{rem:ch7transport}
Equation \eqr{7.8} is usually met in elementary mechanics in words, and
the words are worth attaching to the symbols, because they make the rest
of Section~\ref{sec:ch7kin} predictable.

\emph{What the second and third terms are.} Put $\bv=\bze$ in \eqr{7.8}.
A particle at rest in the frame of $O$, sitting at the place $O^{+}$
calls $\bx^{+}$, is recorded by $O^{+}$ as having the velocity
\begin{equation*}
   \bv^{+}_{\text{tr}}(\bx^{+},t)=\bOm\big(\bx^{+}-\bc\big)+\dot\bc .
\end{equation*}
This is the \emph{transport} velocity: the velocity $O^{+}$ attributes to
the point of $O$'s frame that the particle happens to occupy. So
\eqr{7.8} says
\[
   \bv^{+}=\underbrace{\bQ\bv}_{\text{what }O\text{ records, rotated}}
   +\underbrace{\bv^{+}_{\text{tr}}}_{\text{what the frame contributes}},
\]
which is the familiar statement that a velocity relative to a moving
frame differs from a velocity relative to a fixed one by the velocity of
the frame at that place. Figure~\ref{fig:ch7transport}(a) is the
triangle.

\emph{The transport field is rigid.} As a field of $\bx^{+}$ it is a
constant $\dot\bc-\bOm\bc$ plus a skew tensor applied to $\bx^{+}$, which
is exactly the form \eqr{4.29} of a rigid velocity field, and
Remark~\ref{rem:ch7omega} said as much. Its gradient is
$\grad^{+}\bv^{+}_{\text{tr}}=\bOm$, a constant tensor and a pure spin;
it stretches nothing, anywhere.

\emph{Hence \eqr{7.33} in one line.} Take $\grad^{+}$ of \eqr{7.8}. The
first term gives $\bQ(\grad\bv)\bQ^{t}=\bQ\bL\bQ^{t}$, by the chain rule
with $\partial\bx/\partial\bx^{+}=\bQ^{t}$ from \eqr{7.3}, which is the
computation performed in general at \eqr{7.48}; the second gives $\bOm$;
and the left-hand side is $\bL^{+}$. So
\[
   \bL^{+}=\bQ\bL\bQ^{t}+\bOm ,
\]
which is \eqr{7.33} without passing through $\dot\bF$ at all. The route
through \eqr{7.31} is the book's and is worth having, because it exhibits
$\bOm$ as arising from $\dot\bQ$; but this route exhibits it as the
gradient of a rigid field, and that is the fact \eqr{7.34} needs. The
splitting into $\bD$ and $\bW$ can now be predicted before it is done:
the observer adds the gradient of a rigid field, a rigid field has no
stretching, so the observer cannot touch $\bD$.

\emph{Why the velocity itself is beyond repair.} Notice what the
decomposition does \emph{not} say. There is no way to subtract
$\bv^{+}_{\text{tr}}$ and be left with something both observers can
compute, because $\bv^{+}_{\text{tr}}$ is not a property of the body: it
depends on $\bOm$, $\bc$ and $\bx^{+}$, and $O$ has no access to any of
them. This is the difference between the velocity, which is irreparably
non-objective, and the velocity \emph{gradient}, whose offending part is
a single constant skew tensor that the splitting \eqr{4.7} isolates and
discards.
\end{remark}

\begin{figure}[htbp]\centering
\renewcommand{\thefigure}{B}
\begin{tikzpicture}[scale=1.0,
   ar/.style={-{Stealth[length=2.3mm]},thick},
   sm/.style={-{Stealth[length=1.8mm]},semithick},
   tn/.style={-{Stealth[length=1.5mm]},thin}]
% ================= (a) the triangle at one point =================
\begin{scope}
  \draw[thick,gray!55,fill=blue!4] (0,0) circle (2.15);
  \fill (0,0) circle (1.3pt);
  \coordinate (P) at (40:0.75);
  \draw[guide] (0,0)--(P);
  \node[font=\small,text=gray!85,anchor=south east] at (48:0.40)
     {$\bx^{+}$};
  \fill (P) circle (1.6pt);
  % Q v  and  the transport velocity  and their sum
  \draw[guide] ($(P)+(15:1.05)$)--($(P)+(-13.6:1.61)$);
  \draw[guide] ($(P)+(-50:0.85)$)--($(P)+(-13.6:1.61)$);
  \draw[ar,blue!55!black] (P)--++(15:1.05);
  \draw[ar,gray!70] (P)--++(-50:0.85);
  \draw[ar,red!65!black] (P)--++(-13.6:1.61);
  \node[font=\small,text=blue!55!black,anchor=south] at ($(P)+(17:1.02)$)
     {$\bQ\bv$};
  \node[font=\small,text=gray!85,anchor=north] at ($(P)+(-50:0.98)$)
     {$\bv^{+}_{\text{tr}}$};
  \node[font=\small,text=red!65!black,anchor=west] at ($(P)+(-12:1.66)$)
     {$\bv^{+}$};
  \draw[sm,gray!70] (128:2.32) arc (128:88:2.32);
  \node[font=\small,text=gray!85,anchor=south east] at (120:2.38)
     {$\bom_{p}$};
  \node[font=\small,align=center,anchor=north] at (0,-2.40)
     {(a) one particle:\\[0.5mm] $\bv^{+}=\bQ\bv+\bv^{+}_{\text{tr}}$};
\end{scope}
% ================= (b) the transport field =================
\begin{scope}[xshift=7.60cm]
  \draw[thick,gray!55,fill=blue!4] (0,0) circle (2.15);
  \fill (0,0) circle (1.3pt);
  \foreach \a in {0,60,...,300}{
     \fill[gray!60] (\a:0.80) circle (1.3pt);
     \draw[sm,gray!70] (\a:0.80)--++({\a-90}:0.48);
     \fill[gray!60] (\a:1.60) circle (1.3pt);
     \draw[ar,gray!70] (\a:1.60)--++({\a-90}:0.96);}
  \draw[sm,gray!70] (128:2.32) arc (128:88:2.32);
  \node[font=\small,text=gray!85,anchor=south east] at (120:2.38)
     {$\bom_{p}$};
  \node[font=\small,align=center,anchor=north] at (0,-2.40)
     {(b) the transport field:\\[0.5mm]
      rigid, so $\grad^{+}\bv^{+}_{\text{tr}}=\bOm$};
\end{scope}
\end{tikzpicture}
\caption{Equation \eqr{7.8} as a decomposition, drawn in the frame of
$O^{+}$ with $\bc=\bze$ for clarity. In (a) the velocity $O^{+}$ records
for one particle is the sum of the velocity $O$ records, carried over by
$\bQ$, and the transport velocity $\bv^{+}_{\text{tr}}$, which is what
$O^{+}$ would record if the particle were nailed to $O$'s frame. Since
the platform of $O^{+}$ turns at $\bom_{p}$, the transport velocity runs
backwards, $\bOm\bs=-\bom_{p}\times\bs$, as
Remark~\ref{rem:ch7omega} requires. In (b) the same field is drawn at
twelve places: its arrows grow linearly with the distance from the axis
and are everywhere tangential, which is to say that it is the velocity
field of a rigid rotation, \eqr{4.29}. That is the whole reason $\bD$
survives a change of observer and $\bW$ does not: the observer adds the
gradient of the field in (b), and the field in (b) stretches nothing.}
\label{fig:ch7transport}
\end{figure}

Using this result it is straightforward, though not short, to show that
the accelerations $\ba^{+}\in\mathbb{E}^{3+}$ and $\ba\in\Et$ of $p$ are
related by
\begin{equation}
   \ba^{+}-\bQ\ba=\ddot\bc+2\bOm\big(\bv^{+}-\dot\bc\big)
   +\big(\dot\bOm-\bOm^{2}\big)\big(\bx^{+}-\bc\big).
   \tag{7.9}\label{eq:7.9}
\end{equation}
This is Problem 1, where it is derived in full from \eqr{7.7} and
\eqr{7.8} with no step omitted. The reader who wants the algebra now
should turn to it; what is needed here is only the reading of the four
terms, and the fact that the left-hand side is not zero.

\begin{remark}[the four terms of \eqr{7.9}]\label{rem:ch7fourterms}
Each term answers to one way in which the frame of $O^{+}$ can fail to be
at rest relative to that of $O$, and each is familiar from elementary
mechanics under a different name.

$\ddot\bc$ is the acceleration of the origin of $O^{+}$'s frame. It is
the only term that survives when the frames do not rotate, and it is what
throws a passenger forward when a bus brakes.

$2\bOm(\bv^{+}-\dot\bc)$ is the \emph{Coriolis} term. It is proportional
to the velocity of the particle relative to the moving frame, and it
vanishes for a particle at rest in that frame; the factor $2$ is not a
convention but comes out of the calculation, as Problem 1 shows, because
$\dot\bQ$ is differentiated once in $\ba^{+}$ and once again through
$\bv$.

$-\bOm^{2}(\bx^{+}-\bc)$ is the \emph{centrifugal} term. With
$\bOm\bs=\bom\times\bs$, it is $-\bom\times(\bom\times\bs)$, which points
away from the axis of $\bom$ with magnitude $|\bom|^{2}$ times the
distance from that axis.

$\dot\bOm(\bx^{+}-\bc)$ is the \emph{Euler} or azimuthal term, present
only when the relative rotation is speeding up or slowing down.

Note that all four are proportional to nothing that the body is doing.
They depend on the relative motion of the observers and on the particle's
position and velocity, and on nothing else. That is exactly the property
that will let them be moved to the other side of the equation of motion
in \eqr{7.62} and called forces.
\end{remark}

\begin{figure}[htbp]\centering
\renewcommand{\thefigure}{C}
\begin{tikzpicture}[scale=1.0,
   ar/.style={-{Stealth[length=2.1mm]},thick},
   sm/.style={-{Stealth[length=1.7mm]},semithick}]
% ================= (a) the translational term =================
\begin{scope}
  \draw[thick,gray!55,fill=blue!4] (0,0) circle (1.15);
  \foreach \p in {(-0.55,0.45),(0.50,0.30),(-0.05,-0.60)}{
     \fill[gray!65] \p circle (1.3pt);
     \draw[sm,red!65!black] \p --++(25:0.62);}
  \node[font=\small,align=center,anchor=north] at (0,-1.35)
     {(a) translational\\[0.4mm] $\ddot\bc$\\[0.4mm]
      uniform};
\end{scope}
% ================= (b) the Coriolis term =================
\begin{scope}[xshift=4.05cm]
  \draw[thick,gray!55,fill=blue!4] (0,0) circle (1.15);
  \fill (0,0) circle (1.2pt);
  \coordinate (P) at (-0.20,0.15);
  \fill[gray!65] (P) circle (1.4pt);
  \draw[ar,blue!55!black] (P)--++(55:0.85);
  \node[font=\footnotesize,text=blue!55!black,anchor=south east]
     at ($(P)+(58:0.88)$) {$\bv^{+}$};
  \draw[ar,red!65!black] (P)--++(-35:0.70);
  \draw[sm,gray!70] (128:1.32) arc (128:92:1.32);
  \node[font=\footnotesize,text=gray!85,anchor=south east] at (124:1.38)
     {$\bom_{p}$};
  \node[font=\small,align=center,anchor=north] at (0,-1.35)
     {(b) Coriolis\\[0.4mm] $2\bOm(\bv^{+}-\dot\bc)$\\[0.4mm]
      $\perp\bv^{+}$, no work};
\end{scope}
% ================= (c) the centrifugal term =================
\begin{scope}[xshift=8.10cm]
  \draw[thick,gray!55,fill=blue!4] (0,0) circle (1.15);
  \fill (0,0) circle (1.2pt);
  \fill[gray!65] (152:0.40) circle (1.3pt);
  \draw[sm,red!65!black] (152:0.40)--(152:0.72);
  \fill[gray!65] (-28:0.80) circle (1.3pt);
  \draw[ar,red!65!black] (-28:0.80)--(-28:1.44);
  \draw[sm,gray!70] (128:1.32) arc (128:92:1.32);
  \node[font=\footnotesize,text=gray!85,anchor=south east] at (124:1.38)
     {$\bom_{p}$};
  \node[font=\small,align=center,anchor=north] at (0,-1.35)
     {(c) centrifugal\\[0.4mm] $-\bOm^{2}(\bx^{+}-\bc)$\\[0.4mm]
      outward, $\propto r$};
\end{scope}
% ================= (d) the Euler term =================
\begin{scope}[xshift=12.15cm]
  \draw[thick,gray!55,fill=blue!4] (0,0) circle (1.15);
  \fill (0,0) circle (1.2pt);
  \fill[gray!65] (152:0.40) circle (1.3pt);
  \draw[sm,red!65!black] (152:0.40)--++(62:0.32);
  \fill[gray!65] (-28:0.80) circle (1.3pt);
  \draw[ar,red!65!black] (-28:0.80)--++(-118:0.62);
  \draw[sm,gray!70] (128:1.32) arc (128:92:1.32);
  \node[font=\footnotesize,text=gray!85,anchor=south east] at (124:1.38)
     {$\dot\bom_{p}$};
  \node[font=\small,align=center,anchor=north] at (0,-1.35)
     {(d) Euler\\[0.4mm] $\dot\bOm(\bx^{+}-\bc)$\\[0.4mm]
      azimuthal, spin-up};
\end{scope}
\end{tikzpicture}
\caption{The four terms of \eqr{7.9}, drawn in the frame of $O^{+}$ as
the four contributions they make to the inertial force \eqr{7.63}. Each
answers to one way in which the frame can move, and each has a signature
that identifies it. (a) depends on neither the place nor the velocity of
the particle, so it is indistinguishable from a uniform gravitational
field, which is the content of Remark~\ref{rem:ch7inertial}. (b) is
perpendicular to $\bv^{+}$, because $\bOm$ is skew, and therefore does no
work: $\ip{\bOm\bv^{+}}{\bv^{+}}=0$ by Lemma~\ref{lem:symskw}. It
vanishes for a particle at rest in the frame. (c) points away from the
axis with magnitude proportional to the distance from it, and is the only
one of the four that a particle at rest in a steadily spinning frame
feels. (d) appears only while the spin rate is changing and is
azimuthal, in the sense opposite to the change. A frame may exhibit any
combination of the four; the spinning planet of Problem 3 has (b) and (c)
and nothing else.}
\label{fig:ch7fourterms}
\end{figure}

We observe that if $\ddot\bc=\bze$, i.e. if $\bc=\bV t+\bc_{0}$ for
constant $\bV$ and $\bc_{0}$, and if $\bQ$ is independent of $t$, then
$\dot\bQ=\Ob$, so $\bOm=\Ob$ and $\dot\bOm=\Ob$, and \eqr{7.9} reduces
term by term to $\ba^{+}-\bQ\ba=\bze$. Thus
\begin{equation}
   \bx^{+}=\bQ\bx+\bV t+\bc_{0}
   \qquad\text{and}\qquad
   \ba^{+}=\bQ\ba .
   \tag{7.10}\label{eq:7.10}
\end{equation}
This is called a \emph{Galilean transformation}. Its significance will be
discussed later, in connection with the equations of motion. Note that
$\bV$ here is a constant velocity and has nothing to do with the left
stretch tensor $\bV$ of \eqr{7.27}; the book uses the same letter for
both and no formula contains the two together.

\begin{remark}[what a Galilean transformation does and does not
change]\label{rem:ch7galilean}
Under \eqr{7.10} the acceleration is objective while the velocity is
still not: \eqr{7.8} with $\bOm=\Ob$ and $\dot\bc=\bV$ gives
$\bv^{+}=\bQ\bv+\bV$, so the two observers disagree about velocity by a
constant vector no matter what the body does. This is the familiar
statement that velocity is relative and acceleration is not, and
\eqr{7.9} shows precisely how much of it survives when the rotation is
allowed to depend on time: none of it.

The class of Galilean transformations is exactly the class that preserves
the form of the equation of motion, which is proved at \eqr{7.64}. It is
worth noticing already that the class is large: $\bQ$ is any fixed
orthogonal tensor, $\bV$ any constant vector, $\bc_{0}$ any fixed vector,
and $a$ in \eqr{7.1} any constant. That is the reason for the remark
after \eqr{7.64} that if one inertial frame exists then infinitely many
do.
\end{remark}


\begin{remark}[where the name comes from, and what the transformation is
for]\label{rem:ch7galileo}
Equation \eqr{7.10} takes three lines to write and is named after a man
who never wrote it. Both facts repay a paragraph, because the
transformation carries more of this book than its brevity suggests.

\emph{What it is, counted.} A Galilean transformation is the change of
observer \eqr{7.3} restricted by $\dot\bQ=\Ob$ and $\ddot\bc=\bze$,
together with the clock shift \eqr{7.1}:
\[
   \bx^{+}=\bQ\bx+\bV t+\bc_{0},
   \qquad
   t^{+}=t+a .
\]
It carries ten free constants: three in the fixed rotation $\bQ$, three
in the boost $\bV$, three in the translation $\bc_{0}$, and one in $a$.
These are the ten ways in which two observers may be set up differently
without either being able to find it out: pointing their axes
differently, drifting apart at constant speed, placing their origins
differently, and starting their clocks at different moments. (Strictly
$\bQ$ is orthogonal rather than a rotation, so the family has a second
piece consisting of the mirrored observers of
Remark~\ref{rem:ch7mirror}; the ten constants count the piece containing
the identity.)

\emph{Why they form a group, and why that matters.} Compose two of them.
By the calculation before \eqr{7.65}, $\bQ^{++}=\bQ^{+}\bQ$ and
$\bc^{++}=\bQ^{+}\bc+\bc^{+}$; a product of constant tensors is constant,
and $\bQ^{+}(\bV t+\bc_{0})+\bV^{+}t+\bc^{+}_{0}$ is again linear in $t$.
The identity and the inverses are immediate. So the Galilean
transformations form a group, and that is precisely what licenses the
sentence after \eqr{7.64}. Being inertial is not a property a frame has
by itself but a relation between two frames, and such a relation is
usable only if it is transitive. Were the composition of two Galilean
transformations not Galilean, a frame inertial relative to an inertial
frame need not have been inertial, and the whole notion would have
dissolved. Remark~\ref{rem:ch7addspin} is the same closure property
stated for the spins.

\emph{What it is for.} One thing, and \eqr{7.64} is it. The Galilean
transformations are exactly the changes of observer under which
\eqr{7.53}$_{3}$ keeps its bare form with no inertial force; by
Remark~\ref{rem:ch7whatinertial} the condition $\vek{i}^{+}=\bze$,
imposed as an identity in $\bx^{+}$ and $\bv^{+}$, is equivalent to
\eqr{7.10} and to nothing weaker. So \eqr{7.10} is not a computational
convenience but a definition in disguise: it names the class of frames in
which the balance laws of Chapters 5 and 6 may be read exactly as they
were written.

\emph{How Galileo came to it.} Not as a formula. The \emph{Dialogo} of
1632 is a defence of the claim that the Earth moves, and the standing
objection to that claim was mechanical rather than astronomical: if the
Earth turned, a stone dropped from a tower would be left behind by the
tower, and a cannon fired westward would outrange one fired eastward.
Galileo's answer, put into the mouth of Salviati, is the ship. Shut
yourself below decks with some flies, a bowl of water holding a fish, a
bottle dripping into a basin, and a friend to throw things to; watch
everything carefully while the ship stands still; then have the ship move
as fast as you please, provided only that the motion is uniform and does
not fluctuate this way and that. You will not find the least change in
any of those effects, nor will any of them tell you whether the ship is
moving or standing still. That is Figure~\ref{fig:ch7ship}, and it is
\eqr{7.10}$_{2}$ stated two centuries before anyone wrote
$\ba^{+}=\bQ\ba$.

Two qualifications keep the history honest. Galileo argued the
\emph{principle}, that a uniform motion shared by everything in a closed
room is undetectable inside it; the \emph{transformation}, and with it
the group, is a much later codification, and the name was attached to it
only when there was a Lorentz transformation to set beside it. And
Galileo's own notion of uniform motion still admitted uniform motion in a
circle, since what he needed was for a stone dropped from a mast to keep
pace with a turning Earth. The sharpening to uniform motion in a straight
line is Descartes' and Newton's, and \eqr{7.10}, with its \emph{constant}
$\bV$ and \emph{constant} $\bQ$, is the sharpened version. Newton states
the conclusion as Corollary V to the laws of motion in the
\emph{Principia} of 1687: the motions of bodies among themselves are the
same whether the space that contains them is at rest or moves uniformly
forwards in a straight line without circular motion. That last clause is
the requirement in \eqr{7.10} that $\bQ$ be constant, and Problem 3 is
what happens when it is dropped.

\emph{Where it stops.} The Galilean group is the symmetry of \eqr{7.1}
and \eqr{7.2} and of nothing else. Maxwell's equations do not share it,
and the attempt to make them share it is where the two assumptions of
Remark~\ref{rem:ch7twoassump} were given up in favour of the Lorentz
transformation. The relation between the two groups is not rivalry but
limit: the Galilean group is what the Poincar\'e group becomes as
$c\to\infty$. Inside a theory that has already committed itself to
\eqr{7.1} and \eqr{7.2}, as this one has, \eqr{7.10} is the last word.
\end{remark}

\begin{figure}[htbp]\centering
\renewcommand{\thefigure}{D}
\begin{tikzpicture}[scale=1.0,
   ar/.style={-{Stealth[length=2.3mm]},thick},
   sm/.style={-{Stealth[length=1.8mm]},semithick}]
% a cabin: \cab{x}{y}{halfwidth}{halfheight}{style}
\def\cab#1#2#3#4#5{\draw[#5] (#1-#3,#2-#4) rectangle (#1+#3,#2+#4);}
% ================= (a) below decks =================
\begin{scope}
  \cab{0}{0}{1.55}{1.10}{thick,gray!55,fill=blue!4}
  % the dripping bottle and its basin
  \fill[gray!60] (-0.68,0.72) rectangle (-0.44,0.98);
  \draw[guide] (-0.56,0.68)--(-0.56,-0.68);
  \foreach \y in {0.38,0.02,-0.36}
     {\fill[blue!55!black] (-0.56,\y) circle (1.6pt);}
  \draw[thick,gray!70] (-0.86,-0.70) arc (180:360:0.30 and 0.18);
  % a second object, simply let fall
  \draw[guide] (0.72,0.86)--(0.72,-0.62);
  \foreach \y in {0.56,0.20,-0.24}
     {\fill[red!65!black] (0.72,\y) circle (1.6pt);}
  \draw[thick,gray!70] (0.50,-0.64)--(0.94,-0.64);
  % water
  \draw[gray!55] plot[domain=0:720,samples=60]
     ({-1.90+\x/720*3.80},{-1.36+0.06*sin(\x)});
  \node[font=\small,align=center,anchor=north] at (0,-1.62)
     {(a) below decks, for $O^{+}$:\\[0.5mm]
      the fall is vertical};
\end{scope}
% ================= (b) from the shore =================
\begin{scope}[xshift=7.10cm]
  \cab{-1.85}{0}{0.95}{0.85}{thick,gray!45,densely dashed}
  \cab{1.05}{0}{0.95}{0.85}{thick,gray!55,fill=blue!4}
  % the bottle at t_0 and the basin at t
  \fill[gray!50] (-2.24,0.50) rectangle (-2.04,0.70);
  \draw[thick,gray!70] (0.48,-0.50) arc (180:360:0.24 and 0.15);
  % the track of one drop, as O records it
  \draw[guide] plot[smooth,tension=0.9] coordinates
     {(-2.14,0.46) (-0.90,0.22) (0.10,-0.16) (0.72,-0.52)};
  \foreach \p in {(-2.14,0.46),(-0.90,0.22),(0.10,-0.16),(0.72,-0.52)}
     {\fill[blue!55!black] \p circle (1.6pt);}
  % the ship's own displacement
  \draw[sm,gray!75] (-1.85,1.12)--(1.05,1.12);
  \node[font=\small,text=gray!85,anchor=south] at (-0.40,1.14)
     {$\bV(t-t_{0})$};
  \draw[gray!55] plot[domain=0:1080,samples=80]
     ({-3.00+\x/1080*5.90},{-1.36+0.06*sin(\x)});
  \node[font=\small,align=center,anchor=north] at (-0.40,-1.62)
     {(b) from the shore, for $O$:\\[0.5mm]
      the same fall, leaning forward};
\end{scope}
\end{tikzpicture}
\caption{Salviati's ship, the argument of the \emph{Dialogo} of 1632 and
the physical content of \eqr{7.10}. The cabin sails at the constant
velocity $\bV$. In (a) the observer $O^{+}$ shut inside it sees the drops
leave the bottle and fall vertically into the basin, and so does
everything else he tries; nothing he can do below decks tells him whether
the ship is moving. In (b) the observer $O$ on the shore sees the same
drop follow a track that leans forward, because to its fall he must add
the ship's displacement $\bV(t-t_{0})$. The two records differ, and by
exactly the amount \eqr{7.10}$_{1}$ predicts; but they differ by a term
linear in $t$, whose second derivative vanishes, so the two observers
agree about every acceleration, which is \eqr{7.10}$_{2}$. That agreement
is what lets $O^{+}$ use the equations of Chapter 6 unaltered, and it is
what fails, and is repaired by \eqr{7.62}, the moment the ship begins to
turn or to change speed.}
\label{fig:ch7ship}
\end{figure}
%=======================================================================
\section{Transformation of kinematic variables}\label{sec:ch7kin}
%=======================================================================

To analyze the motion, observers $O$ and $O^{+}$ select reference
configurations $\kappa$ and $\kappa^{+}$ respectively, in which the
positions of $p$ are $\bX=\kappa(p)$ and $\bX^{+}=\kappa^{+}(p)$. The
motions of $p$ in the two frames are then given by
\begin{equation}
   \bx=\bchi_{\kappa}(\bX,t)
   \qquad\text{and}\qquad
   \bx^{+}=\bchi^{+}_{\kappa^{+}}(\bX^{+},t^{+}),
   \tag{7.11}\label{eq:7.11}
\end{equation}
where $\bchi_{\kappa}(\bX,t)=\bchi(\kappa^{-1}(\bX),t)$ and
$\bchi^{+}_{\kappa^{+}}(\bX^{+},t^{+})
=\bchi^{+}((\kappa^{+})^{-1}(\bX^{+}),t^{+})$, exactly as in
Definition~\ref{def:ref}. Thus \eqr{7.6} becomes
\begin{equation}
   \bchi^{+}_{\kappa^{+}}(\bX^{+},t^{+})=\bQ(t)\bchi_{\kappa}(\bX,t)
   +\bc(t).
   \tag{7.12}\label{eq:7.12}
\end{equation}
The two sides are functions of different arguments, $\bX^{+}$ on the left
and $\bX$ on the right, and the equation is to be read as holding when
both label the same material point $p$. Finding the relation between
$\bX$ and $\bX^{+}$ is therefore the first order of business.

Suppose the observers adopt $\kappa_{t_{0}}$ and $\kappa^{+}_{t^{+}_{0}}$
as their respective reference configurations, the two being the placement
of the body at one and the same instant, since $t^{+}_{0}=t_{0}+a$ by
\eqr{7.1}. Following the development of Section~\ref{sec:shifter}, we
interpret this to mean
\begin{equation}
   \bchi_{\kappa}(\bX,t_{0})=\shift\bX
   \qquad\text{and}\qquad
   \bchi^{+}_{\kappa^{+}}(\bX^{+},t^{+}_{0})=\shift^{+}\bX^{+},
   \tag{7.13}\label{eq:7.13}
\end{equation}
with $\bX$ and $\bX^{+}$ regarded as position vectors in $\hEt$ and
$\hat{\mathbb{E}}^{3+}$ respectively, and where
$\shift\colon\hEt\to\Et$ and
$\shift^{+}\colon\hat{\mathbb{E}}^{3+}\to\mathbb{E}^{3+}$ are the
uniform, orthogonal shifters of Definition~\ref{def:shifter} in the two
frames. Equation \eqr{7.13} is \eqr{2.76} written twice, once for each
observer.

Now set $t=t_{0}$ in \eqr{7.12} and use \eqr{7.13} on both sides:
\[
   \shift^{+}\bX^{+}=\bQ(t_{0})\shift\bX+\bc(t_{0}).
\]
Multiply on the left by $(\shift^{+})^{t}$ and use
$(\shift^{+})^{t}\shift^{+}=\hI^{+}$, which is \eqr{2.72} read in the
frame of $O^{+}$, $\hI^{+}$ being the identity of
$\hat{\mathbb{E}}^{3+}$:
\begin{equation}
   \bX^{+}=\bK\bX+\bc^{+},
   \tag{7.14}\label{eq:7.14}
\end{equation}
where
\begin{equation}
   \bK=(\shift^{+})^{t}\bQ(t_{0})\shift
   \qquad\text{and}\qquad
   \bc^{+}=(\shift^{+})^{t}\bc(t_{0}),
   \tag{7.15}\label{eq:7.15}
\end{equation}
with the shifters evaluated at the times $t_{0}$ and $t^{+}_{0}$. Here
$\bK\colon\hEt\to\hat{\mathbb{E}}^{3+}$ is a fixed orthogonal tensor and
$\bc^{+}\in\hat{\mathbb{E}}^{3+}$ is a fixed vector.

Both claims about $\bK$ need checking, and one of them is the most
important single fact in the chapter.

\emph{$\bK$ is orthogonal.} It is a product of three orthogonal maps,
$\shift\colon\hEt\to\Et$, then $\bQ(t_{0})\colon\Et\to\mathbb{E}^{3+}$,
then $(\shift^{+})^{t}\colon\mathbb{E}^{3+}\to\hat{\mathbb{E}}^{3+}$, so
the composition maps $\hEt$ to $\hat{\mathbb{E}}^{3+}$ as claimed, and
\[
   \bK^{t}\bK=\shift^{t}\bQ^{t}(t_{0})\shift^{+}(\shift^{+})^{t}
   \bQ(t_{0})\shift
   =\shift^{t}\bQ^{t}(t_{0})\I^{+}\bQ(t_{0})\shift
   =\shift^{t}\I\shift=\shift^{t}\shift=\hI ,
\]
using $\shift^{+}(\shift^{+})^{t}=\I^{+}$ from \eqr{2.73},
$\bQ^{t}\bQ=\I$ from Proposition~\ref{prop:orth}, and
$\shift^{t}\shift=\hI$ from \eqr{2.72}.

\emph{$\bK$ is fixed.} Every one of the three factors in \eqr{7.15}$_{1}$
is evaluated at the single instant $t_{0}$. The tensor $\bQ(t)$ moves;
the tensor $\bQ(t_{0})$ is one particular tensor, chosen once, and does
not.

\begin{remark}[why the whole chapter splits in two at
\eqr{7.15}]\label{rem:ch7Kfixed}
The fixity of $\bK$ is the reason the transformation laws that follow
come in exactly two families, and it is worth stating before any of them
is derived, because it makes all of them predictable.

Referential objects, those living on $\hEt$, are related to their
counterparts on $\hat{\mathbb{E}}^{3+}$ by $\bK$ alone: $\bU$, $\bC$,
$\bE$, $\bS$, $\bq_{\kappa}$, $\bN$. Spatial objects, those living on
$\Et$, are related by $\bQ(t)$ alone: $\bB$, $\bV$, $\bT$, $\bq$, $\bn$,
$\bD$. Two-point objects get one of each: $\bF$, $\bR$, $\bP$, $\bH$,
$\shift$. The rule is not a coincidence but the alphabet test of
Remark~\ref{rem:hatinch4} applied to the transformation law itself, and
Remark~\ref{rem:2ptlaw} of Chapter 2 already made the general point that
the transformation law identifies the tensor's type.

The consequence is decisive. Since $\bK$ does not depend on time,
$\dot\bK=\hOb$, and differentiating a referential quantity produces no
extra term at all: whatever survives a change of observer before
differentiation survives it afterwards. Since $\bQ$ does depend on time,
differentiating a spatial quantity produces $\dot\bQ$, and that is where
every anomaly of this chapter comes from. Compare \eqr{7.30}$_{1}$, in
which $\bC^{+}=\bK\bC\bK^{t}$ is as clean after differentiation as
before, with \eqr{7.33}, in which $\bL$ picks up $\bOm$. Nothing else has
changed between the two; only which tensor does the conjugating.

This is also why constitutive theory is written in terms of referential
strain measures whenever possible. A referential measure cannot be
polluted by the observer's motion, because the observer's motion has been
evaluated at $t_{0}$ and frozen.
\end{remark}

Fixing $t$, and hence $t^{+}$, and combining \eqr{2.34} with \eqr{7.3}
and $\dl\bX^{+}=\bK\,\dl\bX$, which follows from \eqr{7.14} because
$\bc^{+}$ is constant, we find that the deformation gradients in the two
frames are related by
\begin{equation}
   \bF^{+}\dl\bX^{+}=\dl\bx^{+}=\bQ(t)\,\dl\bx=\bQ(t)\bF\,\dl\bX
   =\bQ(t)\bF\bK^{t}\,\dl\bX^{+},
   \tag{7.16}\label{eq:7.16}
\end{equation}
i.e., since $\dl\bX^{+}$ is arbitrary, that
\begin{equation}
   \bF^{+}=\bQ(t)\bF\bK^{t}.
   \tag{7.17}\label{eq:7.17}
\end{equation}
Every step of \eqr{7.16} deserves its justification. The first equality
is the definition \eqr{2.34} of the deformation gradient in the frame of
$O^{+}$. The second is \eqr{7.4} applied to a pair of neighbouring
places, which is legitimate at fixed $t$ because $\bQ$ is then a
constant tensor and a differential is taken at fixed $t$. The third is
\eqr{2.34} in the frame of $O$. The fourth substitutes
$\dl\bX=\bK^{t}\,\dl\bX^{+}$, obtained from $\dl\bX^{+}=\bK\,\dl\bX$ by
multiplying on the left by $\bK^{t}$ and using $\bK^{t}\bK=\hI$.

This is a linear map from $\hat{\mathbb{E}}^{3+}$ to $\mathbb{E}^{3+}$,
as the alphabet test confirms: $\bK^{t}$ carries
$\hbE_{A}\otimes\hbE^{+}_{B}$, $\bF$ carries $\be_{i}\otimes\hbE_{A}$,
$\bQ$ carries $\be^{+}_{j}\otimes\be_{i}$, and the contractions leave
$\be^{+}_{j}\otimes\hbE^{+}_{B}$. Because the shifter $\shift^{+}$ is
also such a map, it is natural to assume that
\begin{equation}
   \shift^{+}=\bQ(t)\shift\bK^{t}.
   \tag{7.18}\label{eq:7.18}
\end{equation}

This recovers \eqr{7.15}$_{1}$ at time $t_{0}$: substituting \eqr{7.18}
into $(\shift^{+})^{t}\bQ(t_{0})\shift$ and using
$\bQ^{t}(t_{0})\bQ(t_{0})=\I$ and $\shift^{t}\shift=\hI$,
\[
   (\shift^{+})^{t}\bQ(t_{0})\shift
   =\big(\bQ(t_{0})\shift\bK^{t}\big)^{t}\bQ(t_{0})\shift
   =\bK\shift^{t}\bQ^{t}(t_{0})\bQ(t_{0})\shift
   =\bK\shift^{t}\shift=\bK\hI=\bK .
\]
It also ensures that $\mathbb{E}^{3+}$ and $O^{+}$ evolve together as
time passes, and it yields $\bF^{+}(\bX^{+},t^{+}_{0})=\shift^{+}$ in
accordance with \eqr{2.79} and the chosen reference configurations:
by \eqr{7.17} and \eqr{2.79},
$\bF^{+}(\bX^{+},t^{+}_{0})=\bQ(t_{0})\bF(\bX,t_{0})\bK^{t}
=\bQ(t_{0})\shift\bK^{t}=\shift^{+}$.

\begin{remark}[why \eqr{7.18} is assumed rather than
derived]\label{rem:ch7whyshiftass}
The shifter is not determined by anything said so far. Its construction
in Definition~\ref{def:shifter} required a choice of fixed orthonormal
vectors $\bE_{A}$ in the current translation space, and $O^{+}$ must
likewise choose fixed orthonormal $\bE^{+}_{A}\in\mathbb{E}^{3+}$.
Nothing in \eqr{7.1} to \eqr{7.17} says which. Equation \eqr{7.18} is
that choice, and the two properties just verified are exactly the reasons
for making it: it agrees with \eqr{7.15} at the instant the reference
configurations were fixed, and it keeps $\shift^{+}$ playing for $O^{+}$
the role $\shift$ plays for $O$ at every later instant as well.

The alternative would be to let $O^{+}$ pick his $\bE^{+}_{A}$ once and
leave them alone in his own frame. That is a legitimate choice too, but
then $\shift^{+}$ would not satisfy \eqr{7.18}, the relation
$\bH=\bF-\shift$ would fail to transform correctly, and the two observers
would disagree about the displacement gradient. The discussion after
\eqr{7.26} is that failure worked out in detail for the sharper case in
which the shifter is dropped altogether.

An \emph{inertial frame} is a frame that neither accelerates nor spins.
In this book we always regard the frame $\Et$ of $O$ as being inertial.
Accordingly the shifter $\shift$ is fixed, independent of time. It then
follows from \eqr{7.18}, with $\dot\bQ=\Ob$ under a Galilean
transformation, that $(\shift^{+})^{\textstyle\cdot}
=\dot\bQ\shift\bK^{t}=\Ob$: the shifter is fixed in any frame related to
$O$ by a Galilean transformation, i.e. in all inertial frames. In a
frame that does spin, $\shift^{+}$ is not fixed, and Problem 2 computes
exactly how fast it turns. What an inertial frame is, what the two ways
of failing to be one are, and what it means for a shifter not to be
fixed, are taken apart in the remark that follows.
\end{remark}

\begin{remark}[what \emph{inertial} means, and which half of it the
shifter feels]\label{rem:ch7whatinertial}
The paragraph just ended packs a definition, a standing assumption and a
computation into four sentences. Each repays unpacking, the more so
because the standing assumption is what allowed Chapter 6 to be written
without mentioning observers at all, and because the last sentence is the
only place in the chapter where an object is allowed to depend on time
for a reason having nothing to do with the body.

\emph{Exactly two things can move.} A change of observer is \eqr{7.3},
$\bx^{+}=\bQ(t)\bx+\bc(t)$, and it has exactly two ingredients. Each may
depend on $t$, and each carries one of the two words in the definition.

The translation $\bc$ is not an abstract offset. Setting $\bx=\bze$ in
\eqr{7.3} gives $\bx^{+}=\bc(t)$: the vector $\bc(t)$ is the place at
which $O^{+}$ sees the origin of $O$'s frame. So $\dot\bc$ is the
velocity of that origin and $\ddot\bc$ its acceleration, both as recorded
by $O^{+}$, and $\ddot\bc\neq\bze$ is what \emph{accelerates} means. The
reading is worth keeping in mind at \eqr{7.63}, where it explains a fact
that surprises everyone who meets it: in Problem 3 the observer stands on
a spinning planet, his station is certainly being accelerated, and yet
$\ddot\bc=\bze$ there. It is $O$'s origin, the centre of the planet, that
$\bc$ tracks, and the observer sees it sitting still a distance $R$
straight down. The acceleration of his own station is not lost; it
reappears in the centrifugal term.

The rotation $\bQ(t)$ is the identification of directions. By \eqr{7.4} a
vector $\bu$ of $\Et$ is seen by $O^{+}$ as $\bQ(t)\bu$, so
$\bQ(t_{2})\bQ^{t}(t_{1})$ is the rotation carrying the direction $O^{+}$
assigned to $\bu$ at $t_{1}$ onto the one he assigns at $t_{2}$. If $\bQ$
is constant, that rotation is the identity for every pair of instants and
the two observers agree permanently about which way is which. If it is
not, the instantaneous rate is $\bOm=\dot\bQ\bQ^{t}$, with axial vector
$\bom$ by Proposition~\ref{prop:axial}, and Remark~\ref{rem:ch7omega} has
already identified it: $\bom$ is the angular velocity of $O$'s frame as
recorded by $O^{+}$, the negative of the rate $\bom_{p}$ at which
$O^{+}$'s own platform turns. That is what \emph{spins} means.

There is no third possibility. A frame is a Euclidean space and has no
shape to distort; \eqr{7.2} forbids it a change of scale; and \eqr{7.3},
which was derived and not assumed, exhausts what is left. Position and
orientation, two ingredients, two defects.
Figure~\ref{fig:ch7framekinds} is the whole of the classification.

\emph{Why acceleration and not motion.} The definition says
\emph{accelerates}, not \emph{moves}, and the asymmetry is not a
convenience. A frame drifting at constant velocity is inertial, because
by Remark~\ref{rem:ch7galilean} the two observers then disagree about
velocity by the fixed vector $\bV$ and about nothing else, and no
measurement returns $\bV$ without first nominating a frame in which to
make it. A frame whose drift is speeding up is a different matter, and
the passenger thrown forward when the bus brakes is the measurement.

\emph{The definition is the vanishing of the inertial force.} The two
conditions are not an extra postulate laid on top of the chapter. They
are \eqr{7.63} read as an identity. Suppose $\vek{i}^{+}=\bze$ for every
particle, that is, for every $\bx^{+}$ and every $\bv^{+}$. The only term
of \eqr{7.63} containing $\bv^{+}$ is $2\bOm\bv^{+}$, so subtracting the
statement for two particles at the same place with different velocities
leaves $\bOm(\bv^{+}_{1}-\bv^{+}_{2})=\bze$ for arbitrary velocities,
whence $\bOm=\Ob$. If that holds at every instant then $\dot\bOm=\Ob$ and
$\bOm^{2}=\Ob$ as well, and what remains of \eqr{7.63} is
$\ddot\bc=\bze$. Conversely all four terms vanish when $\bOm=\Ob$ and
$\ddot\bc=\bze$; and $\bOm=\Ob$ returns $\dot\bQ=\bOm\bQ=\Ob$. Hence
\[
\begin{aligned}
   \vek{i}^{+}=\bze\ \text{for every particle}
   &\iff \dot\bQ=\Ob\ \ \text{and}\ \ \ddot\bc=\bze\\[0.6mm]
   &\iff \bx^{+}=\bQ\bx+\bV t+\bc_{0},\ \text{which is \eqr{7.10}} .
\end{aligned}
\]
The first line says that a body free of tractions and of body force
satisfies $\ba^{+}=\bze$, i.e. that free particles move uniformly; the
second is the definition of an inertial frame; the third is a Galilean
transformation. The three are one statement wearing three costumes, and
\eqr{7.64} is the fourth.

\emph{Inertial relative to what.} Neither $\bQ$ nor $\bc$ exists until
two observers have been named, so ``neither accelerates nor spins'' is,
as it stands, a statement about a \emph{pair} of frames. The theory does
not pretend otherwise. What is postulated is that one inertial frame
exists, which is the content of Euler's postulates \eqr{6.34} and
\eqr{6.35}; the sentence after \eqr{7.64} then manufactures infinitely
many others, one for each Galilean transformation; and the standing
assumption that $\Et$ is inertial merely names one member of that class.
Every statement made afterwards is a statement about the class, not about
the member.

The postulate is nevertheless more than a definition, because a single
frame can be tested from inside it. Isolate a particle from all contact
and all body force and watch it: by \eqr{7.62} it moves uniformly if and
only if $\vek{i}^{+}=\bze$. Two instruments split the test in two. An
accelerometer reports $\ddot\bc$, subject to the caveat of
Remark~\ref{rem:ch7inertial} that no local measurement separates a
uniform $\ddot\bc$ from a uniform $\bb^{+}$. A gyroscope, or a Foucault
pendulum, reports $\bOm$, and that measurement carries no such caveat,
because the caveat rests on both quantities being uniform and the two
terms $\bOm$ contributes to \eqr{7.63} are not: one grows with the
particle's distance from the axis and the other with its velocity. No
body force does that. Uniform motion is undetectable from inside a closed
room; rotation is not.

\emph{Which half of it the shifter feels.} Return now to
$\shift^{+}=\bQ(t)\shift\bK^{t}$. Of the three factors $\shift$ is fixed,
being built in Definition~\ref{def:shifter} from vectors chosen once, and
$\bK$ is fixed by \eqr{7.15}, every factor in it having been evaluated at
the single instant $t_{0}$. Only $\bQ$ can move, so
\[
   \big(\shift^{+}\big)^{\textstyle\cdot}
   =\dot\bQ\shift\bK^{t}
   =\big(\dot\bQ\bQ^{t}\big)\big(\bQ\shift\bK^{t}\big)
   =\bOm\shift^{+},
\]
which is Problem 2. The formula deserves reading twice, once for what is
in it and once for what is not.

What is not in it is $\ddot\bc$. A frame that accelerates in a straight
line, however violently, has a fixed shifter. The reason is that the
shifter is a tensor, tensors act on vectors, and vectors are differences
of positions in which the translation cancels exactly as it does in
\eqr{7.4}. The whole translational half of non-inertiality is invisible
to it.

What is in it is $\bOm$, alone and undiluted. The shifter turns if and
only if the frame spins, and then at precisely the frame's own rate. So
of the two ways a frame can fail to be inertial the shifter feels one and
ignores the other, and knowing which is which saves a good deal of
worrying: whatever else an observer who merely accelerates disturbs, it
is not the shifter.

\emph{What a moving shifter actually looks like.} The statement can be
made completely concrete. Let $O^{+}$ label his reference space by
$\hbE^{+}_{A}=\bK\hbE_{A}$, which is legitimate because $\bK$ is
orthogonal and therefore carries an orthonormal basis of $\hEt$ onto one
of $\hat{\mathbb{E}}^{3+}$, and which is the choice that gives $\bX$ and
$\bX^{+}$ the same components in \eqr{7.14}. Then
$\bK=\bK\hI=(\bK\hbE_{A})\otimes\hbE_{A}=\hbE^{+}_{A}\otimes\hbE_{A}$, so
$\bK^{t}=\hbE_{A}\otimes\hbE^{+}_{A}$, and \eqr{7.18} with
$\shift=\bE_{B}\otimes\hbE_{B}$ and Lemma~\ref{lem:dyadprod} gives
\[
\begin{aligned}
   \shift^{+}
   &=\bQ\big(\bE_{B}\otimes\hbE_{B}\big)
    \big(\hbE_{A}\otimes\hbE^{+}_{A}\big)
   =\bQ\big(\bE_{A}\otimes\hbE^{+}_{A}\big)
   =\bE^{+}_{A}\otimes\hbE^{+}_{A},\\[0.6mm]
   \bE^{+}_{A}(t)&=\bQ(t)\bE_{A},
\end{aligned}
\]
and therefore $\dot\bE^{+}_{A}=\dot\bQ\bE_{A}=\bOm\bE^{+}_{A}$.

Read as an instruction this says something quite definite. The three
directions $O^{+}$ must use to build his shifter are the images of
$O$'s. They are \emph{not} the directions he calls north, east and up,
which are fixed in his own frame and which he has every reason to regard
as fixed; they are a triad that keeps pointing where $O$'s triad points.
For an observer on a turntable those are the directions of the fixed
stars, and his shifter triad is a gyroscope. He is quite entitled to say
that it turns, because in his frame it does, backwards at the rate of his
own platform. For an observer standing on the Earth it comes round once
per sidereal day, at $\omega\simeq7.29\times10^{-5}\,\mathrm{s}^{-1}$:
slow, but not zero. Figure~\ref{fig:ch7gyro} is this paragraph in two
pictures, and Problem 3 is this observer.

\emph{What is given up, and what is not.} Definition~\ref{def:shifter}
asked for \emph{fixed} orthonormal $\bE_{A}$, and
$\bE^{+}_{A}(t)=\bQ(t)\bE_{A}$ is not fixed, so \eqr{7.18} does relax
that definition for every observer but the inertial one. It is worth
being precise about how much is relaxed, because the answer is: the word
\emph{fixed}, and nothing else. At each instant $\{\bE^{+}_{A}(t)\}$ is
still orthonormal, $\bQ$ being orthogonal, so $\shift^{+}$ is still a
uniform orthogonal two-point tensor and every identity of
Proposition~\ref{prop:shifter} holds of it at every instant; in
particular $(\shift^{+})^{t}\shift^{+}=\hI^{+}$ and
$\shift^{+}(\shift^{+})^{t}=\I^{+}$, which were used at \eqr{7.14} and
again at \eqr{7.21}, remain available. Uniformity in $\bX$, which is what
makes $\shift^{+}\bX^{+}$ differentiate to $\shift^{+}$ and hence makes
\eqr{7.25} work, is untouched. Constancy in $t$ is the single property
surrendered, and it is surrendered only where the frame spins.

\emph{Why the surrender is the right way round.} The alternative offered
in the second paragraph of Remark~\ref{rem:ch7whyshiftass} can now be
priced. Suppose $O^{+}$ nails his $\bE^{+}_{A}$ to the walls of his
laboratory. His shifter is then constant,
$(\shift^{+})^{\textstyle\cdot}=\Ob$, but $\bF^{+}$ is not: for a body
sitting motionless in $O$'s frame, $\dot\bF^{+}=\dot\bQ\bF\bK^{t}
=\bOm\bF^{+}$. The difference $\bH^{+}=\bF^{+}-\shift^{+}$ would then
change in time while absolutely nothing happened to the body, and
$O^{+}$ would read his own rotation as a growing displacement gradient.
With \eqr{7.18} instead, both terms obey the same evolution law and
$\dot\bH^{+}=\bOm\bH^{+}$: the displacement gradient rides round with the
observer and reports the body alone. Problem 2 works this out; it is the
dynamic form of the objection raised after \eqr{7.26}. The shifter of a
spinning observer must be gyroscopic, and if it is not, his displacement
gradient measures him.

\emph{Nothing physical is being smuggled in.} It is worth closing on why
none of this should be troubling. The shifter is not a field on the body
and carries no mechanical information whatever; by
Remark~\ref{rem:shifteruse} it exists only so that
$\bchi_{\kappa}(\bX,t)-\shift\bX$ has both of its terms in the same
space. Requiring it to turn therefore costs nothing, because there was
nothing in it to conserve. What the requirement buys is that $\bH$ means
the same thing to every observer, which is precisely the property
Remark~\ref{rem:ch7shiftertrap} shows cannot be arranged by any choice of
bases.
\end{remark}

\begin{figure}[htbp]\centering
\renewcommand{\thefigure}{E}
\begin{tikzpicture}[scale=1.0,
   ar/.style={-{Stealth[length=2.2mm]},thick},
   sm/.style={-{Stealth[length=1.9mm]},semithick}]
% a small pair of axes standing for the frame of O+ at one instant:
%   \frm{x}{y}{angle}{colour}{line width}
\def\frm#1#2#3#4#5{%
  \begin{scope}[shift={(#1,#2)},rotate=#3]
    \draw[-{Stealth[length=1.9mm]},#4,#5] (0,0)--(0.80,0);
    \draw[-{Stealth[length=1.9mm]},#4,#5] (0,0)--(0,0.80);
    \fill[#4] (0,0) circle (1.2pt);
  \end{scope}}
% ================= (a) inertial =================
\begin{scope}
  \draw[ar] (-2.15,1.80)--(-1.45,1.80);
  \node[font=\small,anchor=west] at (-1.38,1.80) {$\bE_{1}$};
  \draw[guide] (-2.25,0)--(2.25,0);
  \frm{-1.85}{0}{0}{gray!55}{thin}
  \frm{-0.35}{0}{0}{gray!55}{thin}
  \frm{1.15}{0}{0}{blue!55!black}{semithick}
  \node[font=\small,text=gray!85,anchor=north] at (-1.85,-0.12) {$t_{0}$};
  \node[font=\small,text=gray!85,anchor=north] at (-0.35,-0.12) {$t_{1}$};
  \node[font=\small,text=blue!55!black,anchor=north] at (1.15,-0.12)
     {$t_{2}$};
  \draw[sm,gray!75] (-1.85,-0.92)--(-0.35,-0.92);
  \draw[sm,gray!75] (-0.35,-0.92)--(1.15,-0.92);
  \node[font=\small,align=center,anchor=north] at (-0.35,-1.30)
     {(a) inertial\\[0.4mm]
      $\dot\bQ=\Ob$, $\ddot\bc=\bze$\\[0.4mm]
      $\vek{i}^{+}=\bze$\\[0.4mm]
      $\shift^{+}$ fixed};
\end{scope}
% ================= (b) accelerating =================
\begin{scope}[xshift=5.80cm]
  \draw[ar] (-2.15,1.80)--(-1.45,1.80);
  \node[font=\small,anchor=west] at (-1.38,1.80) {$\bE_{1}$};
  \draw[guide] (-2.25,0)--(2.25,0);
  \frm{-1.95}{0}{0}{gray!55}{thin}
  \frm{-1.25}{0}{0}{gray!55}{thin}
  \frm{0.85}{0}{0}{red!65!black}{semithick}
  \node[font=\small,text=gray!85,anchor=north east] at (-1.90,-0.12)
     {$t_{0}$};
  \node[font=\small,text=gray!85,anchor=north west] at (-1.20,-0.12)
     {$t_{1}$};
  \node[font=\small,text=red!65!black,anchor=north] at (0.85,-0.12)
     {$t_{2}$};
  \draw[sm,gray!75] (-1.95,-0.92)--(-1.25,-0.92);
  \draw[sm,red!65!black] (-1.25,-0.92)--(0.85,-0.92);
  \node[font=\small,align=center,anchor=north] at (-0.35,-1.30)
     {(b) accelerating\\[0.4mm]
      $\dot\bQ=\Ob$, $\ddot\bc\neq\bze$\\[0.4mm]
      $\vek{i}^{+}=\ddot\bc$\\[0.4mm]
      $\shift^{+}$ still fixed};
\end{scope}
% ================= (c) spinning =================
\begin{scope}[xshift=11.60cm]
  \draw[ar] (-2.15,1.80)--(-1.45,1.80);
  \node[font=\small,anchor=west] at (-1.38,1.80) {$\bE_{1}$};
  \begin{scope}[shift={(-0.35,0)}]
    \draw[guide] (0,0) circle (0.80);
    \frm{0}{0}{0}{gray!55}{thin}
    \frm{0}{0}{40}{gray!55}{thin}
    \frm{0}{0}{80}{red!65!black}{semithick}
    \node[font=\small,text=gray!85,anchor=north west] at (-10:0.90)
       {$t_{0}$};
    \node[font=\small,text=red!65!black,anchor=south east] at (94:0.90)
       {$t_{2}$};
    \draw[sm,red!65!black] (10:1.25) arc (10:70:1.25);
    \node[font=\small,text=red!65!black,anchor=south west] at (40:1.38)
       {$\bom_{p}$};
  \end{scope}
  \node[font=\small,align=center,anchor=north] at (-0.35,-1.30)
     {(c) spinning\\[0.4mm]
      $\dot\bQ\neq\Ob$\\[0.4mm]
      Coriolis, centrifugal\\[0.4mm]
      $(\shift^{+})^{\textstyle\cdot}=\bOm\shift^{+}$};
\end{scope}
\end{tikzpicture}
\caption{The two ways a frame can fail to be inertial, and which of them
the shifter feels. Each panel is drawn in the frame of $O$ and shows the
frame of $O^{+}$ --- a pair of axes, that is, the directions $O^{+}$
calls fixed --- at three successive instants, beside one direction
$\bE_{1}$ that really is fixed in $\Et$. In (a) the frame drifts at
constant velocity, so that the steps taken in equal times are equal:
$\bQ$ is constant, $\bc$ is linear in $t$, the transformation is the
Galilean one \eqr{7.10}, and $O^{+}$ is inertial. In (b) the steps grow,
so $\ddot\bc\neq\bze$, while the axes nevertheless stay parallel to
$\bE_{1}$ for ever. In (c) the origin does not move at all and the axes
turn, $\dot\bQ\neq\Ob$, and the velocity- and position-dependent terms
of \eqr{7.63} appear. The last line of each panel is the point of the
figure: by \eqr{7.18} the shifter is carried by $\bQ$ alone, so it is
blind to (b) and sensitive only to (c). The shifter records the
orientation of the frame and nothing else.}
\label{fig:ch7framekinds}
\end{figure}

\begin{figure}[htbp]\centering
\renewcommand{\thefigure}{F}
\begin{tikzpicture}[scale=1.0,
   ar/.style={-{Stealth[length=2.4mm]},thick},
   sm/.style={-{Stealth[length=2.0mm]},semithick}]
% ================= (a) as O records it =================
\begin{scope}
  \draw[thick,gray!55,fill=blue!4] (0,0) circle (1.90);
  \fill (0,0) circle (1.4pt);
  % the mark painted on the rim, at t_0 and at t
  \draw[line width=1.5pt,gray!70] (0:1.70)--(0:2.12);
  \node[font=\small,text=gray!85,anchor=west] at (0:2.20) {$t_{0}$};
  \draw[line width=1.5pt,red!65!black] (55:1.70)--(55:2.12);
  \node[font=\small,text=red!65!black,anchor=south west] at (55:2.20)
     {$t$};
  \draw[sm,red!65!black] (12:2.58) arc (12:46:2.58);
  \node[font=\small,text=red!65!black,anchor=south west] at (50:2.60)
     {$\vartheta$};
  % the shifter triad: it does not move
  \draw[ar,blue!55!black] (0,0)--(20:1.55);
  \draw[ar,blue!55!black] (0,0)--(110:1.55);
  \node[font=\small,text=blue!55!black,anchor=north] at (9:1.28)
     {$\bE_{1}$};
  \node[font=\small,text=blue!55!black,anchor=east] at (126:1.24)
     {$\bE_{2}$};
  \node[font=\small,align=center,anchor=north] at (0,-2.35)
     {(a) as $O$ records it:\\[0.5mm]
      the platform turns, $\bE_{A}$ does not};
\end{scope}
% ================= (b) as O+ records it =================
\begin{scope}[xshift=8.60cm]
  \draw[thick,gray!55,fill=blue!4] (0,0) circle (1.90);
  \fill (0,0) circle (1.4pt);
  % his platform: at rest, by definition
  \draw[line width=1.5pt,gray!70] (0:1.70)--(0:2.12);
  \node[font=\small,text=gray!85,anchor=south west] at (0:2.20)
     {$t_{0}$ and $t$};
  % the shifter triad, turning backwards
  \draw[ar,blue!55!black,densely dashed] (0,0)--(20:1.55);
  \draw[ar,blue!55!black,densely dashed] (0,0)--(110:1.55);
  \node[font=\small,text=blue!55!black,anchor=south east] at (118:1.62)
     {$\bE^{+}_{A}(t_{0})$};
  \draw[ar,blue!55!black] (0,0)--(-35:1.55);
  \draw[ar,blue!55!black] (0,0)--(55:1.55);
  \node[font=\small,text=blue!55!black,anchor=south west] at (57:1.62)
     {$\bE^{+}_{A}(t)$};
  \draw[sm,blue!55!black] (-8:2.45) arc (-8:-40:2.45);
  \node[font=\small,text=blue!55!black,anchor=west] at (-38:2.62)
     {$\bom=-\bom_{p}$};
  \node[font=\small,align=center,anchor=north] at (0,-2.35)
     {(b) as $O^{+}$ records it:\\[0.5mm]
      the platform is at rest, $\bE^{+}_{A}(t)$ turns};
\end{scope}
\end{tikzpicture}
\caption{What it means for the shifter not to be fixed. An observer
$O^{+}$ stands on a turntable spinning at $\bom_{p}$, and the mark on the
rim records how far it has gone round, through the angle
$\vartheta=|\bom_{p}|(t-t_{0})$. The two panels are the same physical
situation, and the blue triad, dashed at $t_{0}$ and solid at $t$, is the
same pair of directions in both. In (a) the vectors $\bE_{A}$ of
Definition~\ref{def:shifter} do not move; the platform does. In (b) the
platform is at rest, by definition, and it
is the triad $\bE^{+}_{A}(t)=\bQ(t)\bE_{A}$ carrying the shifter
$\shift^{+}=\bE^{+}_{A}\otimes\hbE^{+}_{A}$ that turns, backwards at the
observer's own rate: $\dot\bE^{+}_{A}=\bOm\bE^{+}_{A}$, which is
$(\shift^{+})^{\textstyle\cdot}=\bOm\shift^{+}$ of Problem 2. The
practical content is in panel (b). Observer $O^{+}$ cannot build his
shifter out of the walls of his own laboratory; he must sight three fixed
stars, or mount a gyroscope. Had he used the walls, his shifter would be
constant, \eqr{7.18} would fail, and $\bH^{+}=\bF^{+}-\shift^{+}$ would
change in time for a body on which nothing whatever was happening.}
\label{fig:ch7gyro}
\end{figure}

Because both observers have selected occupiable reference
configurations, we have that
$0<J^{+}=\det\bF^{+}=J\det\bQ\det\bK$, in which $J=\det\bF>0$ by
\eqr{2.50}. The determinant of a two-point tensor is computed from its
component matrix relative to a right-handed orthonormal basis in each of
the two spaces, and $\det(\bA\bB)=\det\bA\det\bB$ holds for the
corresponding matrices by Theorem~\ref{thm:det}; that is all that is
being used, together with $\det\bK^{t}=\det\bK$. Then, because $\bQ$ and
$\bK$ are orthogonal, each of $\det\bQ$ and $\det\bK$ is $\pm1$ by
Proposition~\ref{prop:orth}, so their product is $\pm1$; and since
$J>0$ and $J^{+}>0$, that product cannot be $-1$. Hence
$\det\bQ\det\bK=1$, and
\begin{equation}
   J^{+}=J.
   \tag{7.19}\label{eq:7.19}
\end{equation}
It is then a simple matter to show that the cofactor of $\bF$ transforms
to
\begin{equation}
   (\bF^{+})^{*}=\bQ\bF^{*}\bK^{t}.
   \tag{7.20}\label{eq:7.20}
\end{equation}
Here is the computation. By Proposition~\ref{prop:I2cof},
$\bA^{*}=(\det\bA)\bA^{-t}$, so $(\bF^{+})^{*}=J^{+}(\bF^{+})^{-t}
=J(\bF^{+})^{-t}$ by \eqr{7.19}. Invert \eqr{7.17}, reversing the order
of the factors as in Proposition~\ref{prop:prodinv} and using
$\bQ^{-1}=\bQ^{t}$, $(\bK^{t})^{-1}=\bK$:
\[
   (\bF^{+})^{-1}=\big(\bQ\bF\bK^{t}\big)^{-1}
   =(\bK^{t})^{-1}\bF^{-1}\bQ^{-1}=\bK\bF^{-1}\bQ^{t},
\]
and transposing,
$(\bF^{+})^{-t}=\big(\bK\bF^{-1}\bQ^{t}\big)^{t}
=\bQ\bF^{-t}\bK^{t}$. Therefore
\[
   (\bF^{+})^{*}=J\bQ\bF^{-t}\bK^{t}=\bQ\big(J\bF^{-t}\big)\bK^{t}
   =\bQ\bF^{*}\bK^{t},
\]
the scalar $J$ passing through $\bQ$ because scalars commute with
everything. So the cofactor transforms exactly as $\bF$ does, which had
to be the case: by Remark~\ref{rem:cofgeom} the cofactor carries area
elements, and an area element is built from a pair of line elements, each
of which transforms with $\bQ$ and $\bK$.

With reference to \eqr{2.77}, the displacement of $p$ in the frame of
$O^{+}$ is
\begin{equation}
\begin{aligned}
   \bu^{+}(\bX^{+},t^{+})
   &=\bchi^{+}_{\kappa^{+}}(\bX^{+},t^{+})-\shift^{+}\bX^{+}\\
   &=\bQ(t)\big[\bchi_{\kappa}(\bX,t)
     -\shift\bK^{t}\big(\bK\bX+\bc^{+}\big)\big]+\bc(t)\\
   &=\bQ(t)\big[\bchi_{\kappa}(\bX,t)-\shift\bX\big]+\bc(t)
     -\bQ(t)\shift\bK^{t}\bc^{+}\\
   &=\bQ(t)\bu(\bX,t)+\bd(t),
\end{aligned}
   \tag{7.21}\label{eq:7.21}
\end{equation}
where $\bu$ is the displacement in the frame of $O$ and
$\bd(t)=\bc(t)-\bQ(t)\shift\bK^{t}\bc^{+}$ vanishes at time $t_{0}$.
Here we have used the fact that $\bK^{t}\bK=\hI$, the identity for
$\hEt$: it is what turns $\shift\bK^{t}\bK\bX$ into $\shift\bX$ between
the second and third lines. The second line substituted \eqr{7.12} for
the first term and \eqr{7.18} together with \eqr{7.14} for the second.

That $\bd(t_{0})=\bze$ is worth verifying, since it is the statement that
the two displacements are both zero at the instant from which they are
measured. At $t=t_{0}$, \eqr{7.18} gives
$\bQ(t_{0})\shift\bK^{t}=\shift^{+}\bK\bK^{t}=\shift^{+}\hI^{+}
=\shift^{+}$, so
\[
   \bd(t_{0})=\bc(t_{0})-\shift^{+}\bc^{+}
   =\bc(t_{0})-\shift^{+}(\shift^{+})^{t}\bc(t_{0})
   =\bc(t_{0})-\I^{+}\bc(t_{0})=\bze ,
\]
using \eqr{7.15}$_{2}$ and then $\shift^{+}(\shift^{+})^{t}=\I^{+}$.

\begin{remark}[the displacement is not objective, but its gradient
is]\label{rem:ch7disp}
Equation \eqr{7.21} is $\bu^{+}=\bQ\bu+\bd$, not $\bu^{+}=\bQ\bu$, so
the displacement fails the test of Definition~\ref{def:ch7obj}. It could
not pass it. A displacement is the difference between a position now and
a position at $t_{0}$, and the two observers were in different relative
positions at those two instants; the mismatch is exactly $\bd(t)$.

The gradient does pass, because $\bd$ is uniform in $\bX$ and a uniform
field has zero gradient. This is the same phenomenon as in
Remark~\ref{rem:ch7twofreedoms}: translations disappear on taking
differences. It is also the reason strain measures may be built from
$\bH$ but not from $\bu$.
\end{remark}

Let $\bH$ be the gradient of $\bu$ with respect to $\bX$, i.e.
$\bH\,\dl\bX=\dl\bu$ at fixed $t$. It follows from \eqr{2.77} by taking
the referential gradient of $\bu=\bchi_{\kappa}-\shift\bX$, the shifter
being uniform so that its gradient is itself, that
\begin{equation}
   \bH=\bF-\shift .
   \tag{7.22}\label{eq:7.22}
\end{equation}
In terms of Cartesian components, cf. \eqr{2.78} and \eqr{2.70},
\begin{equation}
   \bH=H_{iA}\,\be_{i}\otimes\hbE_{A},\qquad\text{where}\qquad
   H_{iA}=\hat u_{i,A}=F_{iA}-\dd_{iA}.
   \tag{7.23}\label{eq:7.23}
\end{equation}
This is a two-point tensor, mapping $\hEt$ to $\Et$.

Let $\bH^{+}$ be the gradient of $\bu^{+}$ with respect to $\bX^{+}$.
Then, from \eqr{7.21}, differentiating at fixed $t$ so that $\bd(t)$
contributes nothing,
\begin{equation}
   \bH^{+}\dl\bX^{+}=\dl\bu^{+}=\bQ(t)\,\dl\bu=\bQ(t)\bH\,\dl\bX
   =\bQ(t)\bH\bK^{t}\,\dl\bX^{+},
   \tag{7.24}\label{eq:7.24}
\end{equation}
yielding, since $\dl\bX^{+}$ is arbitrary,
\begin{equation}
   \bH^{+}=\bQ(t)\bH\bK^{t}=\bF^{+}-\shift^{+}.
   \tag{7.25}\label{eq:7.25}
\end{equation}
The second equality is \eqr{7.22} transformed: subtracting \eqr{7.18}
from \eqr{7.17},
$\bF^{+}-\shift^{+}=\bQ\bF\bK^{t}-\bQ\shift\bK^{t}
=\bQ(\bF-\shift)\bK^{t}=\bQ\bH\bK^{t}$, the middle step by the
distributivity of tensor multiplication. This, of course, is also a
two-point tensor, mapping $\hat{\mathbb{E}}^{3+}$ to $\mathbb{E}^{3+}$.

It is at this stage that the significance of the shifter becomes
apparent. In particular, it ensures that every tensor in \eqr{7.22} is of
the same type, and hence that the relation remains invariant under basis
transformations in a given frame. It also ensures that the relationship
between the deformation and displacement gradients is the same in all
frames: \eqr{7.22} holds for $O$ and \eqr{7.25}$_{2}$, which is the same
sentence with every symbol starred, holds for $O^{+}$.

In contrast, in the extant literature the shifter in \eqr{7.22} is
typically replaced by the referential identity $\hI$. That tensor
transforms differently. From \eqr{7.14},
\[
   \hI^{+}\dl\bX^{+}=\dl\bX^{+}=\bK\,\dl\bX=\bK\hI\,\dl\bX
   =\bK\hI\bK^{t}\,\dl\bX^{+},
\]
whence, $\dl\bX^{+}$ being arbitrary,
\begin{equation}
   \hI^{+}=\bK\hI\bK^{t}.
   \tag{7.26}\label{eq:7.26}
\end{equation}

Now, if \eqr{7.22} is replaced by $\bH=\bF-\hI$, then
\eqr{7.25}$_{1}$ yields
\[
\begin{aligned}
   \bH^{+}&=\bQ(t)\bH\bK^{t}=\bQ(t)\big(\bF-\hI\big)\bK^{t}
   \\
   &=\bQ(t)\bF\bK^{t}-\bQ(t)\hI\bK^{t}
   =\bF^{+}-\bQ(t)\hI\bK^{t}\neq\bF^{+}-\hI^{+},
\end{aligned}
\]
since $\bQ(t)\hI\bK^{t}$ and $\bK\hI\bK^{t}$ differ whenever
$\bQ(t)$ differs from $\bK$, which it does at every instant other than
$t_{0}$ and, in general, even there. So the proposed relationship, if
deemed valid by $O$, must be regarded by $O^{+}$ as being invalid. This
would imply that $O$ enjoys a privileged status, not shared by other
observers. Such a state of affairs runs counter to the philosophy
underlying modern physics. We will consider the displacement gradient
again later, in the context of linear elasticity theory.

\begin{remark}[the same objection, stated so that it cannot be
evaded]\label{rem:ch7shiftertrap}
The temptation is to answer that $\hI$ and $\shift$ are the same thing
when the two reference bases are aligned with the spatial one, so that
$\dd_{iA}$ really are Kronecker deltas, and that the distinction is
therefore pedantry. Remark~\ref{rem:shifteruse} already answered that for
one observer: $\bF-\hI$ subtracts a map $\hEt\to\hEt$ from a map
$\hEt\to\Et$ and is meaningless whatever the numerical values of the
components. The present objection is stronger and does not depend on
type checking at all.

Grant, for the sake of argument, that $O$ has aligned his bases so that
the components of $\shift$ and of $\hI$ agree, and that he therefore sees
no difference between the two formulas. The tensor $\bQ(t)$ still turns.
At any instant $t\neq t_{0}$ the observer $O^{+}$, aligning \emph{his}
bases in the same way, computes $\hI^{+}$ by \eqr{7.26} and
$\shift^{+}$ by \eqr{7.18}, and these are not the same tensor, because
one is built from the frozen $\bK$ and the other from the moving
$\bQ(t)$. So a convention that costs $O$ nothing costs $O^{+}$ the
correctness of his formula. There is no way to align bases once and for
all that repairs this, because $\bQ(t)$ is a function of time and a
choice of bases is not.

The general lesson is worth extracting, because it recurs in Chapter 8.
An identity that holds in one frame and is written with objects that
transform by different rules is not an identity at all; it is a statement
about that frame. Testing a formula against a change of observer is the
cheapest available check on whether it is a law or a convention.
\end{remark}

Continuing, we have the polar decompositions \eqr{3.32} of $\bF$ and the
polar decompositions
\begin{equation}
   \bF^{+}=\bR^{+}\bU^{+}=\bV^{+}\bR^{+}
   \tag{7.27}\label{eq:7.27}
\end{equation}
of $\bF^{+}$ in terms of the symmetric, positive-definite stretch tensors
$\bU^{+}$ and $\bV^{+}$ and the two-point rotation $\bR^{+}$, all
guaranteed to exist and to be unique by Theorem~\ref{thm:polar} applied
in the frame of $O^{+}$. Then, substituting \eqr{7.17} and inserting an
identity in the one place where it does the required work,
\begin{equation}
\begin{aligned}
   \bR^{+}\bU^{+}&=\bQ\bR\bU\bK^{t}=\bQ\bR\,\hI\,\bU\bK^{t}
     =\big(\bQ\bR\bK^{t}\big)\big(\bK\bU\bK^{t}\big)
   \qquad\text{and}\\[1mm]
   \bV^{+}\bR^{+}&=\bQ\bV\bR\bK^{t}=\bQ\bV\,\I\,\bR\bK^{t}
     =\big(\bQ\bV\bQ^{t}\big)\big(\bQ\bR\bK^{t}\big),
\end{aligned}
   \tag{7.28}\label{eq:7.28}
\end{equation}
the inserted identities being $\hI=\bK^{t}\bK$ in the first line and
$\I=\bQ^{t}\bQ$ in the second. The device is the one used throughout
Chapter 3: to convert a product into a product of factors of prescribed
type, insert an identity written as a product of the tensor that has the
right domain and its transpose.

The uniqueness of the polar decompositions then furnishes
\begin{equation}
   \bU^{+}=\bK\bU\bK^{t},\qquad
   \bV^{+}=\bQ\bV\bQ^{t}
   \qquad\text{and}\qquad
   \bR^{+}=\bQ\bR\bK^{t},
   \tag{7.29}\label{eq:7.29}
\end{equation}
but only once it has been checked that the two brackets in each line of
\eqr{7.28} really are of the kind that Theorem~\ref{thm:polar} declares
unique. Three checks are needed.

\emph{$\bQ\bR\bK^{t}$ is a two-point rotation
$\hat{\mathbb{E}}^{3+}\to\mathbb{E}^{3+}$.} It is a product of orthogonal
maps, hence orthogonal, by the computation already performed for $\bK$
after \eqr{7.15}; and its determinant is
$(\det\bQ)(\det\bR)(\det\bK)=(\det\bQ)(\det\bK)=1$ by \eqr{7.19} and
$\det\bR=1$.

\emph{$\bK\bU\bK^{t}$ is symmetric and positive definite on
$\hat{\mathbb{E}}^{3+}$.} Symmetry:
$(\bK\bU\bK^{t})^{t}=\bK\bU^{t}\bK^{t}=\bK\bU\bK^{t}$, using
$(\bA\bB\bC)^{t}=\bC^{t}\bB^{t}\bA^{t}$ and the symmetry of $\bU$.
Positive definiteness: for $\ba\neq\bze$ in $\hat{\mathbb{E}}^{3+}$,
\[
   \ip\ba{\bK\bU\bK^{t}\ba}=\ip{\bK^{t}\ba}{\bU\bK^{t}\ba}>0 ,
\]
by Definition~\ref{def:pd} applied to $\bU$, since $\bK^{t}\ba\neq\bze$:
$\bK^{t}$ is invertible.

\emph{$\bQ\bV\bQ^{t}$ is symmetric and positive definite on
$\mathbb{E}^{3+}$.} The same two computations with $\bQ$ for $\bK$ and
$\bV$ for $\bU$.

So each line of \eqr{7.28} exhibits $\bF^{+}$ as a rotation times a
symmetric positive definite tensor, in the two possible orders, and
uniqueness identifies the factors term by term, giving \eqr{7.29}.

The right- and left-Cauchy--Green tensors in the frame of $O^{+}$ are
thus given by
\begin{equation}
\begin{aligned}
   \bC^{+}&=(\bU^{+})^{2}=\bK\bU\bK^{t}\bK\bU\bK^{t}
    =\bK\bU\,\hI\,\bU\bK^{t}=\bK\bU^{2}\bK^{t}=\bK\bC\bK^{t}
   \qquad\text{and}\\[1mm]
   \bB^{+}&=(\bV^{+})^{2}=\bQ\bV\bQ^{t}\bQ\bV\bQ^{t}
    =\bQ\bV\,\I\,\bV\bQ^{t}=\bQ\bV^{2}\bQ^{t}=\bQ\bB\bQ^{t},
\end{aligned}
   \tag{7.30}\label{eq:7.30}
\end{equation}
respectively. These also follow directly from \eqr{7.17} with
$\bC^{+}=(\bF^{+})^{t}\bF^{+}$ and $\bB^{+}=\bF^{+}(\bF^{+})^{t}$:
\[
\begin{aligned}
   \bC^{+}&=\big(\bQ\bF\bK^{t}\big)^{t}\bQ\bF\bK^{t}
    =\bK\bF^{t}\bQ^{t}\bQ\bF\bK^{t}=\bK\bF^{t}\,\I\,\bF\bK^{t}
    =\bK\bC\bK^{t},\\
   \bB^{+}&=\bQ\bF\bK^{t}\big(\bQ\bF\bK^{t}\big)^{t}
    =\bQ\bF\bK^{t}\bK\bF^{t}\bQ^{t}=\bQ\bF\,\hI\,\bF^{t}\bQ^{t}
    =\bQ\bB\bQ^{t}.
\end{aligned}
\]

\begin{remark}[which measures of deformation see the observer, and
why]\label{rem:ch7whosees}
Equations \eqr{7.29} and \eqr{7.30} are the complete answer to the
question the chapter was written for, at least as far as deformation is
concerned, and it is worth reading them as a list.

$\bU$ and $\bC$ do not see $\bQ$ at all. They are referential, and by
Remark~\ref{rem:ch7Kfixed} they can only be touched by the fixed $\bK$.
If both observers use the same reference configuration, so that
$\bK=\hI$, they write down literally the same tensors. This is the
algebraic form of the statement, made intuitively in
Remark~\ref{rem:CBintuition}, that $\bC$ measures lengths and angles of
material fibres: lengths and angles are what \eqr{7.2} guarantees the two
observers agree about.

$\bV$ and $\bB$ see only $\bQ$, and see it in the objective way of
Definition~\ref{def:ch7obj}. They measure the same lengths and angles,
but record them in the current space, which the two observers hold
differently.

$\bR$ and $\bF$ see both, one on each side. They are two-point tensors
and cannot be otherwise.

The practical consequence, which Chapter 8 will use and Chapter 10 will
use again, is that a scalar function of $\bC$ is automatically agreed on
by all observers, whereas a scalar function of $\bF$ is not. That is why
constitutive equations reduce to functions of $\bC$, and why
Chapter 10 will be able to say that a kinematic constraint may not see
$\bR$: by \eqr{7.29}$_{3}$ the rotation is exactly the part of $\bF$
that a change of observer alters. The reduction is not a modelling
preference; it is forced by \eqr{7.29}$_{3}$.
\end{remark}

Equation \eqr{7.17} also furnishes, on differentiating with respect to
$t$ and remembering that $\bK$ is fixed so that $\dot\bK=\hOb$,
\begin{equation}
   \dot\bF^{+}=\bQ\dot\bF\bK^{t}+\dot\bQ\bF\bK^{t}.
   \tag{7.31}\label{eq:7.31}
\end{equation}
Using $\dot\bF^{+}=\bL^{+}\bF^{+}$ and $\dot\bF=\bL\bF$, both instances
of \eqr{4.3}, where $\bL^{+}$ and $\bL$ respectively are the spatial
velocity gradients as perceived by $O^{+}$ and $O$, the left side of
\eqr{7.31} becomes $\bL^{+}\bF^{+}$ and the first term on the right
becomes $\bQ\bL\bF\bK^{t}$. Both terms on the right now contain the
factor $\bF\bK^{t}$, and \eqr{7.17} identifies it:
$\bF\bK^{t}=\bQ^{t}\bQ\bF\bK^{t}=\bQ^{t}\bF^{+}$. Substituting,
\[
   \bQ\bL\bF\bK^{t}=\bQ\bL\bQ^{t}\bF^{+},
   \qquad
   \dot\bQ\bF\bK^{t}=\dot\bQ\bQ^{t}\bF^{+}=\bOm\bF^{+},
\]
the second using the definition \eqr{7.8}$_{2}$ of $\bOm$. We conclude
that
\begin{equation}
   \bL^{+}\bF^{+}=\bQ\bL\bQ^{t}\bF^{+}+\dot\bQ\bQ^{t}\bF^{+},
   \tag{7.32}\label{eq:7.32}
\end{equation}
and right-multiplication by $(\bF^{+})^{-1}$, which exists because
$\det\bF^{+}=J>0$ by \eqr{7.19}, gives
\begin{equation}
   \bL^{+}=\bQ\bL\bQ^{t}+\bOm ,
   \tag{7.33}\label{eq:7.33}
\end{equation}
where $\bOm$ is the skew tensor defined in \eqr{7.8}$_{2}$.

The stretching and spin tensors thus transform as
\begin{equation}
   \bD^{+}=\bQ\bD\bQ^{t}
   \qquad\text{and}\qquad
   \bW^{+}=\bQ\bW\bQ^{t}+\bOm .
   \tag{7.34}\label{eq:7.34}
\end{equation}
The passage from \eqr{7.33} to \eqr{7.34} is the splitting \eqr{4.7}
applied in the frame of $O^{+}$, and it goes through in three steps.
First, conjugation by $\bQ$ commutes with transposition:
\[
   \big(\bQ\bL\bQ^{t}\big)^{t}=\bQ\bL^{t}\bQ^{t},
\]
so that
\[
   \Sym\big(\bQ\bL\bQ^{t}\big)
   =\tfrac12\bQ\big(\bL+\bL^{t}\big)\bQ^{t}=\bQ\bD\bQ^{t},
   \qquad
   \Skw\big(\bQ\bL\bQ^{t}\big)=\bQ\bW\bQ^{t},
\]
using the linearity of $\bA\mapsto\bQ\bA\bQ^{t}$ to take the factor
$\tfrac12$ and the sum through. Second, $\bOm$ is skew, so
$\Sym\bOm=\Ob^{+}$ and $\Skw\bOm=\bOm$: the extra term contributes
nothing at all to the symmetric part and all of itself to the skew part.
Third, the splitting of a tensor into symmetric and skew parts is unique,
Problem 14 of Chapter 1 and Definition~\ref{def:splits}, so the
symmetric parts of the two sides of \eqr{7.33} may be equated and
likewise the skew parts. That is \eqr{7.34}.

\begin{remark}[why $\bD$ survives and $\bW$ does not]\label{rem:ch7DW}
This is the central result of the chapter and it should not be left as
two lines of algebra.

\emph{The mechanism.} By Remark~\ref{rem:ch7omega} the contamination
$\bOm$ is the velocity gradient of a rigid displacement, and \eqr{4.29}
established that a rigid displacement has zero stretching and pure spin.
So the pollution introduced by the observer is, by its very nature, spin
and nothing but spin. The splitting \eqr{4.7} sorts $\bL$ into a
stretching part and a spinning part; the observer can only add to the
second bin. That is the whole reason, and it explains why the result
could have been predicted without any calculation.

\emph{Why it is not an accident of the algebra.} $\bD$ measures rates of
change of lengths and angles. Equation \eqr{4.15} makes this exact:
$\ip\bm{\bD\bm}=(\ln\mu)^{\textstyle\cdot}$, the rate of logarithmic
stretch of the material fibre in the direction $\bm$. Both observers
agree about $\mu$, because \eqr{7.2} says they agree about lengths, so
they must agree about $\dot\mu/\mu$, and hence about $\bD$ up to the
rotation that relates their pictures. $\bW$ measures no length and no
angle; by Remark~\ref{rem:Wtrap} it is not even the spin of a material
fibre. There is therefore nothing in \eqr{7.2} to protect it, and nothing
does.

\emph{The counterexample, which is the point.} Let the body be at rest
for $O$: $\bv=\bze$, hence $\bL=\Ob$, $\bD=\Ob$, $\bW=\Ob$. Let $O^{+}$
spin steadily. Then \eqr{7.33} gives $\bL^{+}=\bOm$ and \eqr{7.34} gives
\[
   \bD^{+}=\Ob^{+},\qquad \bW^{+}=\bOm\neq\Ob^{+} .
\]
The two observers agree that nothing is being deformed, and they disagree
completely about the spin. $O^{+}$ records a non-zero $\bW^{+}$ although
no material fibre changes its length or its angle to any other. The spin
was manufactured by the act of looking. Figure~\ref{fig:ch7DW} is that
example drawn.

\emph{What it forbids.} Suppose someone proposed the constitutive
equation $\bT=2\mu\bW$ for some material. In the example above $O$ would
compute $\bT=\Ob$ and $O^{+}$ would compute
$\bT^{+}=2\mu\bOm\neq\Ob^{+}$. But by \eqr{7.40} the two are obliged to
satisfy $\bT^{+}=\bQ\bT\bQ^{t}=\Ob^{+}$. The proposed law is therefore
not merely inaccurate but self-contradictory: it asserts that a body
sitting still on a table develops stress if the person looking at it
starts to turn round. The same argument disqualifies $\bT=f(\bL)$ for any
$f$ that actually depends on the skew part of $\bL$, and it is why the
viscous fluid of Chapter 9 is written with $\bD$ and never with $\bL$.
That prohibition is worth the whole chapter.
\end{remark}

\begin{figure}[htbp]\centering
\renewcommand{\thefigure}{G}
\begin{tikzpicture}[scale=1.0,
   ar/.style={-{Stealth[length=2.2mm]},thick},
   sm/.style={-{Stealth[length=1.7mm]},semithick}]
% ================= what O records =================
\begin{scope}
  \draw[thick,fill=blue!6] (0,0) circle (1.15);
  \foreach \a in {0,45,...,315}
     {\fill (\a:1.15) circle (1.1pt);}
  \fill (0,0) circle (1.2pt);
  % "no motion" marks
  \foreach \a in {20,110,200,290}
     {\node[font=\footnotesize,text=gray!80] at (\a:0.62) {$\cdot$};}
  \node[font=\small,anchor=north] at (0,-1.45) {observer $O$};
  \node[font=\small,align=center,anchor=north] at (0,-1.90)
     {$\bv=\bze$\\[0.6mm] $\bD=\Ob,\quad \bW=\Ob$};
  \node[font=\footnotesize,text=gray!85,anchor=south] at (0,1.32)
     {the body sits still};
\end{scope}
% ================= the change of observer =================
\draw[ar] (2.10,0.25) .. controls (3.10,0.80) and (4.10,0.80) .. (5.10,0.25);
\node[font=\small,anchor=south] at (3.60,0.78)
   {$\bQ(t)$,\ \ $\bOm=\dot\bQ\bQ^{t}\neq\Ob$};
% ================= what O+ records =================
\begin{scope}[xshift=7.2cm]
  \draw[thick,fill=blue!6] (0,0) circle (1.15);
  \foreach \a in {0,45,...,315}
     {\fill (\a:1.15) circle (1.1pt);}
  \fill (0,0) circle (1.2pt);
  % tangential velocity arrows: the body appears to turn
  \foreach \a in {0,45,...,315}
     {\draw[sm,red!65!black] (\a:1.15) -- ++(\a+90:0.55);}
  \draw[ar,red!65!black] (1.55,-1.05) arc (-34:34:1.9);
  \node[font=\small,anchor=north] at (0,-1.45) {observer $O^{+}$};
  \node[font=\small,align=center,anchor=north] at (0,-1.90)
     {$\bv^{+}=\bOm\bx^{+}+\dot\bc\neq\bze$\\[0.6mm]
      $\bD^{+}=\Ob^{+},\quad \bW^{+}=\bOm$};
  \node[font=\footnotesize,text=gray!85,anchor=south] at (0,1.32)
     {the body appears to turn};
\end{scope}
\end{tikzpicture}
\caption{One motion, two records: the counterexample of
Remark~\ref{rem:ch7DW}. The body is at rest for $O$ and the observer
$O^{+}$ spins, so $O^{+}$ sees every material point moving on a circle.
The circle of material drawn in each panel is a circle in both, and every
marked fibre keeps its length and its angle to every other; that is what
\eqr{7.2} guarantees, and it is why both observers record
$\bD=\Ob$. The velocity fields differ, and so do the spins: $O$ records
$\bW=\Ob$ and $O^{+}$ records $\bW^{+}=\bOm$, the entire spin having been
manufactured by the observer's own rotation. A constitutive law written
in $\bW$ would give these two panels different stresses, which is
impossible by \eqr{7.40}.}
\label{fig:ch7DW}
\end{figure}

To any unit vector $\bn\in\Et$ we can associate two positions
$\bx,\by\in\Et$ by $\bn=(\by-\bx)/|\by-\bx|$
(Figure~\ref{fig:ch7unit}). Its counterpart in $\mathbb{E}^{3+}$ is
\begin{equation}
   \bn^{+}=\frac{\by^{+}-\bx^{+}}{\big|\by^{+}-\bx^{+}\big|}
   =\frac{\by^{+}-\bx^{+}}{|\by-\bx|}
   =\frac{\bQ(\by-\bx)}{|\by-\bx|}
   =\bQ\bn ,
   \tag{7.35}\label{eq:7.35}
\end{equation}
the second equality by \eqr{7.2} and the third by \eqr{7.4}; the last
step takes the scalar $|\by-\bx|^{-1}$ inside the linear map $\bQ$.
Similarly, for a unit vector $\bN\in\hEt$ we have its counterpart in
$\hat{\mathbb{E}}^{3+}$:
\begin{equation}
   \bN^{+}=\frac{\bY^{+}-\bX^{+}}{\big|\bY^{+}-\bX^{+}\big|}
   =\frac{\bY^{+}-\bX^{+}}{|\bY-\bX|}
   =\frac{\bK(\bY-\bX)}{|\bY-\bX|}
   =\bK\bN ,
   \tag{7.36}\label{eq:7.36}
\end{equation}
now using \eqr{7.14}, from which $\bY^{+}-\bX^{+}=\bK(\bY-\bX)$ because
the constant $\bc^{+}$ cancels in the difference, and
$|\bK(\bY-\bX)|=|\bY-\bX|$ because $\bK$ is orthogonal.

\begin{remark}[why \eqr{7.35} is a theorem and \eqr{7.38} will be an
assumption]\label{rem:ch7unitwhy}
It would have been possible to declare that unit vectors transform by
$\bn^{+}=\bQ\bn$ and be done with it. The book does not, and the reason
is instructive.

A unit vector attached to a configuration is not a primitive object. It
is built from two positions, and positions are the primitives about which
\eqr{7.2} and \eqr{7.3} already say everything. So its transformation law
is a consequence, and \eqr{7.35} derives it in three lines with nothing
assumed that was not assumed in Section~\ref{sec:ch7obs}.

Contrast \eqr{7.38} below. A traction is not built out of positions; it
is a primitive of the mechanical theory, introduced in Chapter 6 by a
postulate of its own. Nothing in the geometry of Section~\ref{sec:ch7obs}
can determine how it transforms, and the statement $\bt^{+}=\bQ\bt$ has
to be adopted as a further physical hypothesis. Keeping the two kinds of
statement apart is what makes it possible to say later exactly which
physical assumptions the theory rests on.
\end{remark}

\bookfig{2}\begin{figure}[htbp]\centering
\begin{tikzpicture}[scale=1.0,
   ar/.style={-{Stealth[length=2.3mm]},thick}]
% ================= as O sees it =================
\begin{scope}
  \blob{configline}
  \coordinate (X) at (-0.55,-0.15);
  \coordinate (Y) at (0.72,0.42);
  \draw[guide] (X)--(Y);
  \fill (X) circle (1.3pt); \fill (Y) circle (1.3pt);
  \node[lbl,anchor=north east] at (X) {$\bx$};
  \node[lbl,anchor=south west] at (Y) {$\by$};
  \draw[ar,red!65!black] (X) -- ($(X)!0.72cm!(Y)$);
  \node[lbl,font=\small,text=red!65!black,anchor=north west]
     at (-0.30,-0.10) {$\bn$};
  \node[font=\small,anchor=north] at (0,-1.55) {$\Et$};
\end{scope}
% ================= the map =================
\draw[ar] (2.05,0.20) .. controls (2.95,0.72) and (3.85,0.72) .. (4.75,0.20);
\node[font=\small,anchor=south] at (3.40,0.70) {$\bQ(t)$};
% ================= as O+ sees it =================
\begin{scope}[xshift=6.9cm,rotate=-38]
  \blob{configline}
  \coordinate (X2) at (-0.55,-0.15);
  \coordinate (Y2) at (0.72,0.42);
  \draw[guide] (X2)--(Y2);
  \fill (X2) circle (1.3pt); \fill (Y2) circle (1.3pt);
  \node[lbl,anchor=north east] at (X2) {$\bx^{+}$};
  \node[lbl,anchor=south west] at (Y2) {$\by^{+}$};
  \draw[ar,red!65!black] (X2) -- ($(X2)!0.72cm!(Y2)$);
  \node[lbl,font=\small,text=red!65!black,anchor=north west]
     at (-0.32,-0.12) {$\bn^{+}$};
\end{scope}
\node[font=\small,anchor=north] at (6.9,-1.75) {$\mathbb{E}^{3+}$};
\end{tikzpicture}
\caption{A unit vector as perceived by two observers. The vector is not
posited but built: $\bn$ is the direction from $\bx$ to $\by$, normalised
by the distance between them. Both the numerator and the denominator are
already governed by Section~\ref{sec:ch7obs}, the first by \eqr{7.4} and
the second by \eqr{7.2}, so the transformation law $\bn^{+}=\bQ\bn$ of
\eqr{7.35} is derived rather than assumed. The same construction in the
reference configuration, with \eqr{7.14} in place of \eqr{7.4}, gives
$\bN^{+}=\bK\bN$, and the change from $\bQ$ to $\bK$ is the whole
difference between a spatial and a referential vector.}
\label{fig:ch7unit}
\end{figure}

%=======================================================================
\section{Further transformations}\label{sec:ch7further}
%=======================================================================

We assume the density, energy, heating supply, temperature and entropy to
be absolute scalars, and hence the same for all observers. Thus
\begin{equation}
   \rho^{+}=\rho,\quad \eps^{+}=\eps,\quad r^{+}=r,\quad
   h^{+}=h,\quad \theta^{+}=\theta
   \quad\text{and}\quad \eta^{+}=\eta .
   \tag{7.37}\label{eq:7.37}
\end{equation}
The first of these is a natural consequence of the fact that neither the
mass of a body, nor its volume, are affected by a change from one
observer to another. Mass is assigned to the body itself in
Section~\ref{sec:mass}, before any observer appears; and volume is
unaffected because, by \eqr{7.4},
\[
\begin{aligned}
   \dl v^{+}&=\big|\tbox{\dl\bx^{+}}{\dl\by^{+}}{\dl\bz^{+}}\big|
   \\
   &=\big|\tbox{\bQ\dl\bx}{\bQ\dl\by}{\bQ\dl\bz}\big|
   =|\det\bQ|\,\big|\tbox{\dl\bx}{\dl\by}{\dl\bz}\big|=\dl v ,
\end{aligned}
\]
by Theorem~\ref{thm:det} and $|\det\bQ|=1$. Equation \eqr{7.37}$_{1}$
combines with \eqr{7.19} to yield the equality of the referential mass
densities: since $\rho_{\kappa}=\rho J$ by the mass conservation of
Section~\ref{sec:mass}, and both factors are unchanged,
$\rho^{+}_{\kappa^{+}}=\rho^{+}J^{+}=\rho J=\rho_{\kappa}$.

\begin{remark}[why \eqr{7.37} is an assumption, and how one would notice
if it were wrong]\label{rem:ch7absolute}
Nothing derived so far bears on any of the six quantities in \eqr{7.37}.
Sections~\ref{sec:ch7obs} to \ref{sec:ch7kin} concerned only clocks,
rulers and the geometry that follows from them; mass, energy and
temperature were never mentioned. So \eqr{7.37} is new physical input,
and it is worth being explicit about that rather than letting it pass as
obvious.

It is also worth seeing that it has content. Suppose temperature were not
absolute, and that a moving observer read a different value. Then the
Clausius--Duhem inequality \eqr{7.53}$_{5}$, which carries $\theta$ in
two places, would hold for one observer and could fail for another, and
the second law of thermodynamics would become a statement about who is
looking. The assumption \eqr{7.37} is precisely what forbids that, and
\eqr{7.59}$_{3}$ is the pay-off.

The one entry that is not simply a postulate is $\rho^{+}=\rho$, since it
follows from mass and volume separately, and the volume computation above
is the proof.
\end{remark}

Further, we assume that the Cauchy and Piola tractions $\bt,\bp\in\Et$
and the body force density $\bb\in\Et$ transform as spatial vectors, i.e.
\begin{equation}
   \bt^{+}=\bQ\bt,\qquad \bp^{+}=\bQ\bp
   \qquad\text{and}\qquad \bb^{+}=\bQ\bb .
   \tag{7.38}\label{eq:7.38}
\end{equation}
The Cauchy stress $\bT^{+}$ in the frame of $O^{+}$ is then such that
\begin{equation}
   \bT^{+}\bn^{+}=\bt^{+}=\bQ\bt=\bQ\bT\bn=\bQ\bT\bQ^{t}\bn^{+}
   \tag{7.39}\label{eq:7.39}
\end{equation}
for all unit vectors $\bn^{+}$; thus
\begin{equation}
   \bT^{+}=\bQ\bT\bQ^{t}.
   \tag{7.40}\label{eq:7.40}
\end{equation}
The chain \eqr{7.39} uses Cauchy's relation $\bt=\bT\bn$ of Chapter 6 in
each frame, \eqr{7.38}$_{1}$ in the middle, and finally
$\bn=\bQ^{t}\bn^{+}$, which is \eqr{7.35} solved for $\bn$ by
multiplying on the left by $\bQ^{t}$ and using $\bQ^{t}\bQ=\I$.

Two small points make the passage from \eqr{7.39} to \eqr{7.40} legal.
First, every unit vector $\bn^{+}\in\mathbb{E}^{3+}$ arises as $\bQ\bn$
for a unit $\bn\in\Et$, namely $\bn=\bQ^{t}\bn^{+}$, which is a unit
vector because $\bQ^{t}$ is orthogonal; so \eqr{7.39} really does hold
for every unit vector of the starred frame and not merely for those that
happen to be images. Second, two tensors that agree on every \emph{unit}
vector agree on every vector: given $\bs\neq\bze$, write
$\bs=|\bs|\,\bn$ with $\bn=\bs/|\bs|$ a unit vector, and linearity gives
$\bA\bs=|\bs|\bA\bn=|\bs|\bB\bn=\bB\bs$, while $\bA\bze=\bze=\bB\bze$.
Hence $\bT^{+}$ and $\bQ\bT\bQ^{t}$ are the same tensor.

\begin{remark}[what kind of statement \eqr{7.38} is, and how much of it
is independent]\label{rem:ch7tractionass}
Three statements were written down in \eqr{7.38}, and it matters that
they are not three.

\emph{Two are assumptions and the third is a theorem.} Equation
\eqr{7.42} will derive $\bp^{+}=\bQ\bp$ from \eqr{7.38}$_{1}$ and the
kinematics of Section~\ref{sec:ch7kin} alone, so \eqr{7.38}$_{2}$ adds
nothing. What is genuinely assumed is that the Cauchy traction and the
body force density are objective spatial vectors. Nothing proves either,
and Remark~\ref{rem:ch7unitwhy} said why: a traction is a primitive of
Chapter 6 rather than a construction out of positions, and the geometry
of Section~\ref{sec:ch7obs} has no hold on it. This is the point at which
the mechanics of the chapter, as opposed to its geometry, enters.

\emph{What the assumption says physically.} That the force transmitted
across a surface is a property of the two pieces of material on either
side of it, and not of the person watching. Cut the body and put a spring
balance in the cut: both observers read the same needle, and they
disagree only about which direction in their own space the reading points
along, by exactly the $\bQ$ that relates their spaces. Stated that way it
is nearly a tautology about what the word \emph{force} means, and that is
why it usually passes without comment. It is nevertheless an assumption,
and it is a strong one: it is what forbids a force whose magnitude
depends on the observer's motion.

\emph{Where it fails, and why the failure is the whole story of
Section~\ref{sec:ch7balance}.} The one force in this chapter that is
\emph{not} objective is the inertial force $\vek{i}^{+}$ of \eqr{7.63},
and it could not be, since by construction it is built out of $\bQ$ and
$\bc$ themselves. That is the precise content of the word
\emph{fictitious}. A centrifugal force has the units of a body force
density, enters the equations of motion exactly like one, and is
distinguished from a real one by a single property: it fails the test of
Definition~\ref{def:ch7obj}. Chapter 6's postulates concerned real body
forces, \eqr{7.38}$_{3}$ is the statement that they are objective, and
\eqr{7.62} is what happens when the two kinds are added together.

\emph{What rests on it.} Because \eqr{7.38}$_{1}$ is an assumption, so is
\eqr{7.40}, and with it the whole of Chapter 8. Had the traction carried
an extra term, in the manner of \eqr{7.8}, Cauchy's relation
$\bt=\bT\bn$ would have produced a $\bT^{+}$ that is not a conjugate of
$\bT$; the stress would not have been objective; the stress power would
not have been invariant, so \eqr{7.59}$_{2}$ and \eqr{7.59}$_{3}$ would
have failed as well; and there would have been nothing left on the stress
side of a constitutive equation for both observers to agree about. One
assumed line, \eqr{7.38}$_{1}$, carries all of that.
\end{remark}

The Piola stress $\bP^{+}$ follows from $\bP=\bT\bF^{*}$, the definition
of Chapter 6, written in each frame and combined with \eqr{7.40} and
\eqr{7.20}:
\begin{equation}
   \bP^{+}=\bT^{+}(\bF^{+})^{*}=\bQ\bT\bQ^{t}\bQ\bF^{*}\bK^{t}
   =\bQ\bT\,\I\,\bF^{*}\bK^{t}=\bQ\bP\bK^{t}.
   \tag{7.41}\label{eq:7.41}
\end{equation}
This furnishes
\begin{equation}
   \bp^{+}=\bP^{+}\bN^{+}=\bQ\bP\bK^{t}\bK\bN=\bQ\bP\,\hI\,\bN
   =\bQ\bP\bN=\bQ\bp ,
   \tag{7.42}\label{eq:7.42}
\end{equation}
in accordance with \eqr{7.38}$_{2}$, using \eqr{7.36} for $\bN^{+}$ and
$\bK^{t}\bK=\hI$. So \eqr{7.38}$_{2}$ was not an independent assumption:
it is forced by \eqr{7.38}$_{1}$ together with the kinematics of
Section~\ref{sec:ch7kin}. Lastly, the Piola--Kirchhoff stress $\bS^{+}$
follows from $\bP^{+}=\bF^{+}\bS^{+}$. Solving,
\[
   \bS^{+}=(\bF^{+})^{-1}\bP^{+}
   =\big(\bK\bF^{-1}\bQ^{t}\big)\big(\bQ\bP\bK^{t}\big)
   =\bK\bF^{-1}\,\I\,\bP\bK^{t}=\bK\big(\bF^{-1}\bP\big)\bK^{t},
\]
using the inverse of \eqr{7.17} computed after \eqr{7.20}; and
$\bF^{-1}\bP=\bS$ by the same relation $\bP=\bF\bS$ in the frame of $O$.
We obtain
\begin{equation}
   \bS^{+}=\bK\bS\bK^{t}.
   \tag{7.43}\label{eq:7.43}
\end{equation}

\begin{remark}[three stresses, three transformation
laws]\label{rem:ch7threestress}
Put \eqr{7.40}, \eqr{7.41} and \eqr{7.43} side by side and they say what
Remark~\ref{rem:2ptlaw} of Chapter 2 promised: the transformation law
announces the type.
\[
\begin{aligned}
   \bT^{+}&=\bQ\bT\bQ^{t}
     &&\text{spatial: }\Et\to\Et ,\\
   \bP^{+}&=\bQ\bP\bK^{t}
     &&\text{two-point: }\hEt\to\Et ,\\
   \bS^{+}&=\bK\bS\bK^{t}
     &&\text{referential: }\hEt\to\hEt .
\end{aligned}
\]
The pattern is exactly that of $\bB$, $\bF$, $\bC$ in \eqr{7.30} and
\eqr{7.17}, and for the same reason: each stress measures force per unit
area, and the three differ only in which configuration supplies the
force and which supplies the area. $\bT$ takes both from the current
configuration, $\bS$ takes both from the reference one, and $\bP$ takes
one from each.

The practically important entry is the last. Because $\bK$ is fixed, a
constitutive equation written for $\bS$ in terms of $\bC$ relates two
objects neither of which sees $\bQ$, and observer agreement is then
automatic. Written for $\bT$ in terms of $\bB$ it is equally automatic,
since both sides carry one $\bQ$ on each side. Written for $\bT$ in terms
of $\bC$ it is not, and that mismatch is what Chapter 8 has to repair.
\end{remark}

The relationship between the heat flux $h$ through a surface with unit
normal $\bn$ and the heat flux vector $\bq$ is $h=-\ip\bq\bn$, cf.
\eqr{6.166}. The corresponding relationship in the frame of $O^{+}$ is
$h^{+}=-\ip{\bq^{+}}{\bn^{+}}$. We thus have, from \eqr{7.35} and
\eqr{7.37}$_{4}$, that
\[
   \ip{\bq^{+}}{\bn^{+}}=\ip\bq\bn=\ip\bq{\bQ^{t}\bn^{+}}
   =\ip{\bQ\bq}{\bn^{+}}
\]
for all unit vectors $\bn^{+}$, the last step being the definition of the
transpose. Hence $\ip{(\bq^{+}-\bQ\bq)}{\bn^{+}}=0$ for every unit
$\bn^{+}$, and therefore for every vector of $\mathbb{E}^{3+}$ by the
scaling argument used after \eqr{7.39}; Lemma~\ref{lem:cvec} then gives
\begin{equation}
   \bq^{+}=\bQ\bq ,
   \tag{7.44}\label{eq:7.44}
\end{equation}
and the referential heat flux vectors in the two frames, defined by
$\bq_{\kappa}=(\bF^{*})^{t}\bq$ in Chapter 6, are related by
\begin{equation}
\begin{aligned}
   \bq^{+}_{\kappa^{+}}&=\big[(\bF^{+})^{*}\big]^{t}\bq^{+}
    =\big(\bQ\bF^{*}\bK^{t}\big)^{t}\bQ\bq\\
   &=\bK(\bF^{*})^{t}\bQ^{t}\bQ\bq=\bK(\bF^{*})^{t}\,\I\,\bq
    =\bK\bq_{\kappa}.
\end{aligned}
   \tag{7.45}\label{eq:7.45}
\end{equation}


\begin{remark}[the whole of it is one commutative square]%
\label{rem:ch7square}
Everything derived in Sections~\ref{sec:ch7kin} and
\ref{sec:ch7further} is a statement about the diagram of
Figure~\ref{fig:ch7square}, and once that diagram is in front of one, the
dictionary below can be read off instead of memorised.

\emph{Four spaces, two vertical maps.} There are four spaces in play: the
reference spaces $\hEt$ and $\hat{\mathbb{E}}^{3+}$ and the current
spaces $\Et$ and $\mathbb{E}^{3+}$. Changing observer means going down,
by $\bK$ on the reference side and by $\bQ(t)$ on the current side, and
the whole content of Remark~\ref{rem:ch7Kfixed} is the difference between
those two verticals: the left one was frozen at $t_{0}$ by \eqr{7.15} and
the right one moves.

\emph{Every two-point law says that the square commutes.} Equation
\eqr{7.17} is $\bF^{+}=\bQ\bF\bK^{t}$, and multiplying on the right by
$\bK$ and using $\bK^{t}\bK=\hI$ turns it into
\[
   \bF^{+}\bK=\bQ\bF ,
\]
that is, into the statement that the two routes from $\hEt$ to
$\mathbb{E}^{3+}$ agree: deform and then change observer, or change
observer and then deform. Read that way, \eqr{7.20}, \eqr{7.25},
\eqr{7.29}$_{3}$ and \eqr{7.41} are the same sentence for $\bF^{*}$,
$\bH$, $\bR$ and $\bP$. And \eqr{7.18} is the demand that the shifter be
a horizontal arrow of the same square. Seen from here it stops looking
like a convention: it is what puts $\shift$ into the diagram at all, and
Remark~\ref{rem:ch7whyshiftass} is the observation that nothing else
would.

\emph{Every one-space law is a conjugation.} A tensor mapping a space to
itself is a loop at one corner, and carrying it down the diagram means
going down, round the loop, and back up: $\bA^{+}=\bQ\bA\bQ^{t}$ on the
right, $\bA^{+}=\bK\bA\bK^{t}$ on the left. That is the entire difference
between the second and third groups of the dictionary, and the polar
decomposition displays it: by \eqr{7.29} and
Figure~\ref{fig:ch7square}(b), $\bU$ is a loop on the reference side,
$\bV$ a loop on the current side, and $\bR$ the horizontal arrow they
share. The pairs $(\bC,\bB)$ and $(\bS,\bT)$ are loops in the same two
places. A vector is a third case: it lives \emph{at} a corner rather than
on an arrow, and it is carried down by a single application of the
vertical, which is $\bN^{+}=\bK\bN$ and $\bq^{+}_{\kappa^{+}}
=\bK\bq_{\kappa}$ on the left and $\bn^{+}=\bQ\bn$, $\bt^{+}=\bQ\bt$,
$\bq^{+}=\bQ\bq$ on the right.

\emph{What the diagram cannot show.} One thing only, and it is the
subject of the rest of the chapter: the right-hand vertical depends on
$t$. Every law free of time derivatives is a statement about the square
as drawn. Every law containing one must differentiate that vertical too,
and $\dot\bQ$ is what comes out. That is why $\bv$, $\ba$, $\bL$ and
$\bW$ are absent from the figure and present in the last group of the
dictionary: they are statements about how the square changes, not about
the square.

\emph{Using it as a check.} The alphabet test of
Remark~\ref{rem:hatinch4} is precisely the instruction to read a
formula's corners off this square, and it catches nearly every error
available in this chapter. A term $\bK\bT\bK^{t}$ is wrong because $\bT$
is a loop on the right while $\bK$ runs down the left; $\bQ\bC\bQ^{t}$ is
wrong for the mirror-image reason; the sum $\bF-\shift$ is legitimate and
$\bF-\hI$ is not, because the first subtracts one horizontal arrow from
another and the second subtracts a loop from a horizontal arrow. That
last is the objection of Remark~\ref{rem:ch7shiftertrap}, drawn.
\end{remark}

\begin{figure}[htbp]\centering
\renewcommand{\thefigure}{H}
\begin{tikzpicture}[scale=1.0,
   ar/.style={-{Stealth[length=2.4mm]},thick},
   sm/.style={-{Stealth[length=2.0mm]},semithick},
   sp/.style={inner sep=2.5pt}]
% ================= (a) the square =================
\begin{scope}
  \node[sp] (TL) at (0,2.45)   {$\hEt$};
  \node[sp] (TR) at (4.10,2.45) {$\Et$};
  \node[sp] (BL) at (0,0)      {$\hat{\mathbb{E}}^{3+}$};
  \node[sp] (BR) at (4.10,0)    {$\mathbb{E}^{3+}$};
  \draw[ar] (TL)--(TR) node[midway,above,font=\small] {$\bF$};
  \draw[ar] (BL)--(BR) node[midway,below,font=\small] {$\bF^{+}$};
  \draw[ar] (TL)--(BL) node[midway,left,font=\small] {$\bK$};
  \draw[ar] (TR)--(BR) node[midway,right,font=\small] {$\bQ(t)$};
  \node[font=\small,anchor=north] at (2.05,1.62) {$\bF^{+}\bK=\bQ\bF$};
  \node[font=\small,align=center,anchor=north] at (2.05,-1.66)
     {(a) the square of \eqr{7.17}:\\[0.5mm]
      $\bK$ fixed, $\bQ(t)$ moving};
\end{scope}
% ================= (b) the polar factors on it =================
\begin{scope}[xshift=8.20cm]
  \node[sp] (tl) at (0,2.45)   {$\hEt$};
  \node[sp] (tr) at (4.10,2.45) {$\Et$};
  \node[sp] (bl) at (0,0)      {$\hat{\mathbb{E}}^{3+}$};
  \node[sp] (br) at (4.10,0)    {$\mathbb{E}^{3+}$};
  \draw[ar] (tl)--(tr) node[midway,above,font=\small] {$\bR$};
  \draw[ar] (bl)--(br) node[midway,below,font=\small] {$\bR^{+}$};
  \draw[ar] (tl)--(bl) node[midway,left,font=\small] {$\bK$};
  \draw[ar] (tr)--(br) node[midway,right,font=\small] {$\bQ(t)$};
  \draw[sm] (tl) to[out=115,in=65,looseness=9] (tl);
  \node[font=\small,anchor=south] at (0,3.30) {$\bU$};
  \draw[sm] (tr) to[out=115,in=65,looseness=9] (tr);
  \node[font=\small,anchor=south] at (4.10,3.30) {$\bV$};
  \draw[sm] (bl) to[out=245,in=295,looseness=9] (bl);
  \node[font=\small,anchor=north] at (0,-0.92) {$\bU^{+}$};
  \draw[sm] (br) to[out=245,in=295,looseness=9] (br);
  \node[font=\small,anchor=north] at (4.10,-0.92) {$\bV^{+}$};
  \node[font=\small,align=center,anchor=north] at (2.05,-1.66)
     {(b) the polar factors of \eqr{7.29}:\\[0.5mm]
      $\bU$, $\bV$ are loops, $\bR$ the arrow};
\end{scope}
\end{tikzpicture}
\caption{The diagram the chapter is about. Going right is deforming;
going down is changing observer. (a) Equation \eqr{7.17} says that the
square commutes, and the same statement with $\bF$ replaced by
$\bF^{*}$, $\bR$, $\bH$, $\bP$ or $\shift$ is \eqr{7.20}, \eqr{7.29}$_{3}$,
\eqr{7.25}, \eqr{7.41} and \eqr{7.18} respectively. The two verticals are
not alike: $\bK$ is one fixed tensor, every factor of it evaluated at
$t_{0}$ by \eqr{7.15}, while $\bQ(t)$ turns, and every anomaly of this
chapter is $\dot\bQ$. (b) A tensor that maps a space to itself is a loop
at a corner, and a loop is carried down the diagram by conjugation:
$\bU^{+}=\bK\bU\bK^{t}$ on the left and $\bV^{+}=\bQ\bV\bQ^{t}$ on the
right, which is \eqr{7.29}. The pairs $(\bC,\bB)$ and $(\bS,\bT)$ are
loops in the same two places, and $\bD$ is a loop on the right that
happens also to be objective; a vector such as $\bN$ or $\bn$ sits at a
corner instead and is carried down by one application of the vertical.
The dictionary below is this figure sorted into a list.}
\label{fig:ch7square}
\end{figure}
\subsection*{The dictionary}

It is convenient to have the results of Sections~\ref{sec:ch7kin} and
\ref{sec:ch7further} in one place, since the remaining sections and the
whole of Chapters 8 to 10 consult them constantly. Sorted by what each
quantity sees:
\[
\begin{aligned}
   \text{nothing:}&\quad
     \rho,\ \eps,\ r,\ h,\ \theta,\ \eta,\ J,\ \rho_{\kappa}
     &&\eqr{7.19},\ \eqr{7.37}\\[0.7mm]
   \bK\text{ only:}&\quad
     \bN^{+}=\bK\bN,\quad \bU^{+}=\bK\bU\bK^{t},
     &&\eqr{7.29}\\
   &\quad \bC^{+}=\bK\bC\bK^{t},
     &&\eqr{7.30}\\
   &\quad \bS^{+}=\bK\bS\bK^{t},\quad
     \bq^{+}_{\kappa^{+}}=\bK\bq_{\kappa},
     &&\eqr{7.43},\ \eqr{7.45}\\
   &\quad \hI^{+}=\bK\hI\bK^{t}
     &&\eqr{7.26}\\[0.7mm]
   \bQ\text{ only:}&\quad
     \bn^{+}=\bQ\bn,\quad \bt^{+}=\bQ\bt,
     &&\eqr{7.35},\ \eqr{7.38}\\
   &\quad \bb^{+}=\bQ\bb,\quad \bq^{+}=\bQ\bq,
     &&\eqr{7.38},\ \eqr{7.44}\\
   &\quad \bV^{+}=\bQ\bV\bQ^{t},\quad \bB^{+}=\bQ\bB\bQ^{t},
     &&\eqr{7.29},\ \eqr{7.30}\\
   &\quad \bT^{+}=\bQ\bT\bQ^{t},
     &&\eqr{7.40}\\
   &\quad \bD^{+}=\bQ\bD\bQ^{t}
     &&\eqr{7.34}\\[0.7mm]
   \text{both:}&\quad
     \bF^{+}=\bQ\bF\bK^{t},\quad \bR^{+}=\bQ\bR\bK^{t},
     &&\eqr{7.17},\ \eqr{7.29}\\
   &\quad \bH^{+}=\bQ\bH\bK^{t},
     &&\eqr{7.25}\\
   &\quad \bP^{+}=\bQ\bP\bK^{t},\quad
     (\bF^{+})^{*}=\bQ\bF^{*}\bK^{t},
     &&\eqr{7.41},\ \eqr{7.20}\\
   &\quad \shift^{+}=\bQ\shift\bK^{t}
     &&\eqr{7.18}\\[0.7mm]
   \dot\bQ\text{ as well:}&&&\\[-1mm]
   &\quad \bv^{+}=\bQ\bv+\bOm(\bx^{+}-\bc)+\dot\bc ,
     &&\eqr{7.8}\\
   &\quad \bL^{+}=\bQ\bL\bQ^{t}+\bOm,
     &&\eqr{7.33}\\
   &\quad \bW^{+}=\bQ\bW\bQ^{t}+\bOm,
     &&\eqr{7.34}\\
   &\quad \ba^{+}=\bQ\ba+\vek{i}^{+}
     &&\eqr{7.9}
\end{aligned}
\]
The first four groups are the objective and the observer-blind
quantities, in the sense of Definition~\ref{def:ch7obj}; they are the raw
material of constitutive theory. The last group is the list of everything
that goes wrong, and every entry in it contains $\dot\bQ$, whether
written as $\dot\bQ$, as $\bOm$, or as $\vek{i}^{+}$.

%=======================================================================
\section{Transformation of the balance laws}\label{sec:ch7balance}
%=======================================================================

Let $\bw(\bx,t)$ be an arbitrary spatial vector in the frame of $O$, and
suppose $\bw^{+}(\bx^{+},t^{+})$ is its counterpart in the frame of
$O^{+}$, i.e. that $\bw$ is objective in the sense of
Definition~\ref{def:ch7obj}:
\begin{equation}
   \bw^{+}(\bx^{+},t^{+})=\bQ(t)\bw(\bx,t).
   \tag{7.46}\label{eq:7.46}
\end{equation}
Then, differentiating at fixed $t$, so that $\bQ$ is a constant tensor
throughout and passes through the differential,
\begin{equation}
   \bQ(\grad\bw)\dl\bx=\bQ\,\dl\bw=\dl\bw^{+}
   =(\grad^{+}\bw^{+})\dl\bx^{+}=(\grad^{+}\bw^{+})\bQ\,\dl\bx ,
   \tag{7.47}\label{eq:7.47}
\end{equation}
using the definition of the gradient of a vector field in each frame,
Definition~\ref{def:diffv}, and $\dl\bx^{+}=\bQ\,\dl\bx$ from \eqr{7.4}.
As $\dl\bx$ is arbitrary the two ends are equal as tensors,
$\bQ(\grad\bw)=(\grad^{+}\bw^{+})\bQ$, and multiplying on the right by
$\bQ^{t}$ therefore gives
\begin{equation}
   \grad^{+}\bw^{+}=\bQ(\grad\bw)\bQ^{t}.
   \tag{7.48}\label{eq:7.48}
\end{equation}
The gradient of an objective vector field is an objective tensor field.

This implies that the divergence is the same in both frames:
\begin{equation}
\begin{aligned}
   \dvg^{+}\bw^{+}&=\tr(\grad^{+}\bw^{+})=\tr\big[\bQ(\grad\bw)\bQ^{t}\big]
   \\
   &=\tr\big[\bQ^{t}\bQ(\grad\bw)\big]=\tr(\grad\bw)=\dvg\bw ,
\end{aligned}
   \tag{7.49}\label{eq:7.49}
\end{equation}
by Proposition~\ref{prop:divv} at the two ends and the cyclic property
$\tr(\bA\bB)=\tr(\bB\bA)$ in the middle.

\begin{remark}[two traces, on two different spaces]\label{rem:ch7twotraces}
The middle step of \eqr{7.49} needs the care that
Remark~\ref{rem:twopointlegal} taught. The first trace is taken on
$\mathbb{E}^{3+}$, since $\bQ(\grad\bw)\bQ^{t}$ maps that space to
itself; the second is taken on $\Et$, since $\bQ^{t}\bQ(\grad\bw)$ maps
$\Et$ to itself. They are different operations on different spaces, and
the assertion is that they produce the same number.

They do, and the reason is the component identity behind
Remark~\ref{rem:twopointlegal}: with
$\bQ=Q_{\alpha i}\be^{+}_{\alpha}\otimes\be_{i}$ and
$\grad\bw=w_{i,j}\be_{i}\otimes\be_{j}$,
\[
   \tr\big[\bQ(\grad\bw)\bQ^{t}\big]=Q_{\alpha i}w_{i,j}Q_{\alpha j},
   \qquad
   \tr\big[\bQ^{t}\bQ(\grad\bw)\big]=Q_{\alpha j}Q_{\alpha i}w_{i,j},
\]
the same sum of products of the same numbers, written in a different
order. Both traces are legitimate, each on its own space, and they are
equal. Every use of the cyclic property in this chapter is of that kind
and none of them is an abuse.
\end{remark}

From \eqr{7.40},
\begin{equation}
   \bT^{+}\bw^{+}=\bQ\bT\bQ^{t}\bQ\bw=\bQ\bT\,\I\,\bw=\bQ(\bT\bw).
   \tag{7.50}\label{eq:7.50}
\end{equation}
If $\bw$ is independent of $\bx$, i.e. $\grad\bw=\Ob$, then $\bw^{+}$ is
independent of $\bx^{+}$, since $\bQ$ is spatially uniform and
\eqr{7.48} gives $\grad^{+}\bw^{+}=\Ob^{+}$; and the definition
\eqr{1.148} of the divergence of a tensor, Definition~\ref{def:divA},
yields
\begin{equation}
\begin{aligned}
   \bQ^{t}(\dvg^{+}\bT^{+})\cdot\bw
   &=\dvg^{+}\bT^{+}\cdot\bw^{+}
    =\dvg^{+}\big[(\bT^{+})^{t}\bw^{+}\big]\\
   &=\dvg(\bT^{t}\bw)
    =\dvg\bT\cdot\bw ,
\end{aligned}
   \tag{7.51}\label{eq:7.51}
\end{equation}
leading, as $\bw$ is arbitrary, to the conclusion that
\begin{equation}
   \dvg^{+}\bT^{+}=\bQ(\dvg\bT).
   \tag{7.52}\label{eq:7.52}
\end{equation}
Each equality of \eqr{7.51} carries its own justification. The first
moves $\bQ^{t}$ across the inner product by the definition of the
transpose and then uses \eqr{7.46}. The second is
Definition~\ref{def:divA} in the frame of $O^{+}$, applicable because
$\bw^{+}$ is uniform there. The third is \eqr{7.49} applied to the
vector field $\bs=\bT^{t}\bw$, which is objective:
\[
   (\bT^{+})^{t}\bw^{+}=\big(\bQ\bT\bQ^{t}\big)^{t}\bQ\bw
   =\bQ\bT^{t}\bQ^{t}\bQ\bw=\bQ\bT^{t}\bw=\bQ\bs=\bs^{+} .
\]
The fourth is Definition~\ref{def:divA} in the frame of $O$. The passage
from \eqr{7.51} to \eqr{7.52} is Lemma~\ref{lem:cvec}: the vector
$\bQ^{t}\dvg^{+}\bT^{+}-\dvg\bT$ has zero inner product with every
uniform $\bw$, hence vanishes, and multiplying by $\bQ$ gives \eqr{7.52}.

Suppose now that $O$ is an inertial observer. The laws holding in the
frame of this observer are the conservation of mass, the balance of
energy, the balances of linear and rotational momenta, and the
Clausius--Duhem inequality:
\begin{equation}
\begin{aligned}
   &\dot\rho+\rho\,\dvg\bv=0,\\
   &\rho\dot\eps+\dvg\bq-\bT\cdot\bL=\rho r,\\
   &\rho\ba-\dvg\bT=\rho\bb,\qquad \bT^{t}=\bT,\\
   &\quad\text{and}\quad
   -\rho\big(\dot\psi+\eta\dot\theta\big)+\bT\cdot\bL
   -\theta^{-1}\ip\bq{\grad\theta}\geq0 ,
\end{aligned}
   \tag{7.53}\label{eq:7.53}
\end{equation}
where $\ba$ is the material acceleration. We seek the corresponding laws
pertaining to an arbitrary observer $O^{+}$. To this end we use
\eqr{7.44}, \eqr{7.46} and \eqr{7.49} to infer that
\begin{equation}
   \dvg^{+}\bq^{+}=\dvg\bq ,
   \tag{7.54}\label{eq:7.54}
\end{equation}
the heat flux vector being objective by \eqr{7.44} and hence an
admissible $\bw$ in \eqr{7.49}.

Also, from $\ip{\grad^{+}\theta^{+}}{\dl\bx^{+}}=\dl\theta^{+}
=\dl\theta=\ip{\grad\theta}{\dl\bx}=\ip{\grad\theta}{\bQ^{t}\dl\bx^{+}}
=\ip{\bQ\grad\theta}{\dl\bx^{+}}$, valid for every $\dl\bx^{+}$ because
$\theta$ is an absolute scalar by \eqr{7.37}$_{5}$ and hence has the same
increment for both observers, Lemma~\ref{lem:cvec} gives
$\grad^{+}\theta^{+}=\bQ(\grad\theta)$, and therefore
\begin{equation}
\begin{aligned}
   \grad^{+}\theta^{+}&=\bQ(\grad\theta),
   \qquad
   \ip{\bq^{+}}{\grad^{+}\theta^{+}}\\
   &=\ip{\bQ\bq}{\bQ\grad\theta}
   =\ip\bq{\bQ^{t}\bQ\grad\theta}=\ip\bq{\grad\theta}.
\end{aligned}
   \tag{7.55}\label{eq:7.55}
\end{equation}

The transformation formula \eqr{7.8} for the material velocity is not of
the form \eqr{7.46}. Nevertheless
\begin{equation}
   \dvg^{+}\bv^{+}=\dvg\bv .
   \tag{7.56}\label{eq:7.56}
\end{equation}
This follows from \eqr{7.33}, i.e.
\begin{equation}
   \dvg^{+}\bv^{+}=\tr\bL^{+}=\tr\big(\bQ\bL\bQ^{t}+\bOm\big)
   =\tr\big(\bQ\bL\bQ^{t}\big)=\tr\big(\bQ^{t}\bQ\bL\big)
   =\tr\bL=\dvg\bv ,
   \tag{7.57}\label{eq:7.57}
\end{equation}
the third equality because $\tr\bOm=0$: a skew tensor has
$\Omega_{\alpha\alpha}=-\Omega_{\alpha\alpha}$ for each $\alpha$, with no
sum, so every diagonal component vanishes.

Further, the symmetry of $\bT$, and hence that of $\bQ\bT\bQ^{t}$, yields
\begin{equation}
\begin{aligned}
   \bT^{+}\cdot\bL^{+}
   &=\bQ\bT\bQ^{t}\cdot\big(\bQ\bL\bQ^{t}+\bOm\big)
    =\bQ\bT\bQ^{t}\cdot\bQ\bL\bQ^{t}\\
   &=\tr\big[\bQ\bT\bQ^{t}\big(\bQ\bL^{t}\bQ^{t}\big)\big]
    =\tr\big(\bQ\bT\bL^{t}\bQ^{t}\big)\\
   &=\tr\big(\bQ^{t}\bQ\bT\bL^{t}\big)=\tr\big(\bT\bL^{t}\big)
    =\bT\cdot\bL .
\end{aligned}
   \tag{7.58}\label{eq:7.58}
\end{equation}
The second equality dropped $\bQ\bT\bQ^{t}\cdot\bOm$, which vanishes by
Lemma~\ref{lem:symskw} because $\bQ\bT\bQ^{t}$ is symmetric,
$(\bQ\bT\bQ^{t})^{t}=\bQ\bT^{t}\bQ^{t}=\bQ\bT\bQ^{t}$, and $\bOm$ is
skew. The third used Definition~\ref{def:tip},
$\bA\cdot\bB=\tr(\bA\bB^{t})$, with
$(\bQ\bL\bQ^{t})^{t}=\bQ\bL^{t}\bQ^{t}$; the fourth cancelled
$\bQ^{t}\bQ=\I$ in the middle; the fifth is the cyclic property once
more.

\begin{remark}[the two invariant scalars, and what they teach about
Definition~\ref{def:ch7obj}]\label{rem:ch7twoscalars}
Equations \eqr{7.56} and \eqr{7.58} are the counterexamples promised in
Remark~\ref{rem:ch7objinv}, and they are worth pausing over because they
are the sharpest available warning against reasoning from ingredients.

In \eqr{7.56} both sides are built from a field, $\bv$, that is not
objective: \eqr{7.8} carries two extra terms. The extra terms survive
differentiation, since $\bL^{+}$ carries $\bOm$; they are annihilated
only at the very last step, by the trace. So a non-objective field can
have a perfectly invariant scalar buried inside it, and no inspection of
the field alone would reveal it.

In \eqr{7.58} the ingredients are mixed: $\bT$ is objective, $\bL$ is
not. What kills the extra term is not a property of $\bL$ at all but the
symmetry of $\bT$, which is itself a balance law, \eqr{7.53}$_{4}$. Had
the material admitted couple stresses, so that $\bT$ were not symmetric,
the stress power would not have been invariant and the energy balance
\eqr{7.59}$_{2}$ would have failed to have the same form for both
observers. The invariance of the stress power is therefore not a
kinematic fact but a consequence of the balance of rotational momentum.
That is worth remembering, because Chapters 8 and 9 use
$\bT\cdot\bL=\bT\cdot\bD$ constantly, and the deletion of $\bW$ there is
the same deletion as here.
\end{remark}

With these results and \eqr{7.37}, together with the invariance of the
Helmholtz free energy, cf. \eqr{6.186}, and of the material derivative,
we conclude that the balances of mass and energy, and the
Clausius--Duhem inequality, are the same in all frames:
\begin{equation}
\begin{aligned}
   &\dot\rho^{+}+\rho^{+}\dvg^{+}\bv^{+}=0,\\
   &\rho^{+}\dot\eps^{+}+\dvg^{+}\bq^{+}-\bT^{+}\cdot\bL^{+}
     =\rho^{+}r^{+},\\
   &\quad\text{and}\quad
   -\rho^{+}\big(\dot\psi^{+}+\eta^{+}\dot\theta^{+}\big)
   +\bT^{+}\cdot\bL^{+}
   -(\theta^{+})^{-1}\ip{\bq^{+}}{\grad^{+}\theta^{+}}\geq0 ,
\end{aligned}
   \tag{7.59}\label{eq:7.59}
\end{equation}
and, of course, \eqr{7.40} and \eqr{7.53}$_{4}$ ensure that
\begin{equation}
   (\bT^{+})^{t}=\bT^{+}.
   \tag{7.60}\label{eq:7.60}
\end{equation}
The verification is term by term, and it is entirely a matter of
substituting the results already obtained. In \eqr{7.59}$_{1}$,
$\rho^{+}=\rho$ by \eqr{7.37}$_{1}$ and $\dvg^{+}\bv^{+}=\dvg\bv$ by
\eqr{7.56}; the material derivative is unchanged because it follows the
material point $p$ and, by \eqr{7.1}, $\partial/\partial t^{+}
=\partial/\partial t$, so $\dot\rho^{+}=\dot\rho$. Every term of
\eqr{7.59}$_{1}$ therefore equals the corresponding term of
\eqr{7.53}$_{1}$, and the equation holds because that one does. In
\eqr{7.59}$_{2}$, use \eqr{7.37}$_{1,2,3}$, \eqr{7.54} and \eqr{7.58}. In
\eqr{7.59}$_{3}$, use \eqr{7.37} again, \eqr{7.58} and \eqr{7.55}$_{2}$.
And \eqr{7.60} is $(\bQ\bT\bQ^{t})^{t}=\bQ\bT^{t}\bQ^{t}
=\bQ\bT\bQ^{t}$.

\begin{remark}[why the symmetry of $\bT$ survives a rotating
observer]\label{rem:ch7symT}
Equation \eqr{7.60} was obtained by conjugation in one line, and the ease
of it hides a question worth asking. The symmetry of $\bT$ is not an
algebraic identity. It is the local form of the balance of moment of
momentum, derived in Chapter 6 for an inertial observer. In the frame of
$O^{+}$ that balance picks up the moment of the inertial force, which
does not vanish and which varies from place to place. Why is the
conclusion unchanged?

\emph{Because a body force never enters the symmetry argument.} Taking
moments about a point of $\mathbb{E}^{3+}$ adds the term
$\int\rho^{+}\,\bs^{+}\times\vek{i}^{+}\,\dl v$, a volume integral of
exactly the same form as the moment
$\int\rho^{+}\,\bs^{+}\times\bb^{+}\,\dl v$ of the true body force. In
the localisation of Chapter 6 the volume terms are one order higher in
the diameter of the shrinking region than the surface terms, and it is
the surface terms alone that survive to give $\eps_{ijk}T_{jk}=0$. Adding
$\vek{i}^{+}$ to $\bb^{+}$ therefore changes the balance of moment of
momentum in precisely the way that changes nothing in its local
consequence.

\emph{The observation generalises, and explains why \eqr{7.62} is the
only casualty.} By \eqr{7.63} a change of observer contributes to the
mechanics through a term of body-force type and through nothing else. Any
law whose local content is insensitive to the body force is therefore
untouched: the balance of moment of momentum for the reason just given,
and the mass balance and the Clausius--Duhem inequality because no force
appears in them at all. Only the balance of linear momentum reads the
body force pointwise, and only it changes.

\emph{The energy balance is the interesting case.} No body force appears
in \eqr{7.53}$_{2}$, but one does appear in the statement of Chapter 6
from which it came, through the working $\int\rho\ip\bb\bv\,\dl v$; it
cancels against the same term in the kinetic energy theorem, and
\eqr{7.59}$_{2}$ is what is left after that cancellation. Adding
$\vek{i}^{+}$ to $\bb^{+}$ adds
$\int\rho^{+}\ip{\vek{i}^{+}}{\bv^{+}}\,\dl v$ to both sides of the
cancellation, so it survives in neither. An inertial force does work on a
body in a rotating frame, and none of that work reaches the internal
energy.

Exactly one of the four terms of \eqr{7.63} does no work at all, and it
is worth saying which. The Coriolis term is orthogonal to $\bv^{+}$
whenever $\dot\bc=\bze$, since $\ip{\bOm\bv^{+}}{\bv^{+}}=0$ for a skew
$\bOm$ by Lemma~\ref{lem:symskw}: a Coriolis force deflects a particle
and never speeds it up or slows it down. The Euler term is orthogonal to
$\bs^{+}$, by the same lemma applied to the skew $\dot\bOm$, which is
what makes it azimuthal; but it is not orthogonal to $\bv^{+}$, and it
does work. The other two do work as well, and they are the two that
possess potentials:
\[
   -\bOm^{2}\bs^{+}=\grad^{+}\psi ,\quad
   \psi=\tfrac12\big|\bOm\bs^{+}\big|^{2},
   \qquad
   \ddot\bc=\grad^{+}\ip{\ddot\bc}{\bx^{+}},
\]
the first because $\ip{\bOm^{2}\bs}{\bs}=-|\bOm\bs|^{2}$ and $\bOm^{2}$
is symmetric. With $\bOm\bs=\bom\times\bs$ the potential $\psi$ is the
familiar $\tfrac12|\bom|^{2}r^{2}$, $r$ being the distance from the axis.
That is why a steadily rotating frame admits an effective potential
energy, and why the Coriolis force cannot be absorbed into one.
\end{remark}

The situation for the linear momentum balance is different. To
investigate this we combine \eqr{7.38}$_{3}$ and \eqr{7.53}$_{3}$,
obtaining
\begin{equation}
\begin{aligned}
   \rho^{+}\bb^{+}&=\rho^{+}\bQ\bb=\rho\bQ\bb
    =\bQ\big(\rho\ba-\dvg\bT\big)\\
   &=\rho\bQ\ba-\bQ(\dvg\bT)=\rho^{+}\bQ\ba-\dvg^{+}\bT^{+},
\end{aligned}
   \tag{7.61}\label{eq:7.61}
\end{equation}
the last step by \eqr{7.52} and \eqr{7.37}$_{1}$ read backwards.

Every term of \eqr{7.53}$_{3}$ has now been carried into the frame of
$O^{+}$ intact except one. The divergence of the stress went over by
\eqr{7.52}, the body force by \eqr{7.38}$_{3}$, and the density is
invariant by \eqr{7.37}$_{1}$; each of the three appears in the starred
frame as the image under $\bQ$ of what $O$ wrote, which is what
objectivity means. The acceleration is the single exception, and
\eqr{7.9} says by exactly how much it fails. Give that failure a name:
\begin{equation*}
   \vek{i}^{+}:=\ba^{+}-\bQ\ba ,
   \qquad\text{equivalently}\qquad
   \bQ\ba=\ba^{+}-\vek{i}^{+}.
\end{equation*}
This introduces nothing. The quantity $\vek{i}^{+}$ is the right-hand
side of \eqr{7.9}, already computed in Problem 1 out of $\bQ(t)$ and
$\bc(t)$ alone; the definition merely gives it a symbol so that it can be
moved about. Substituting it into \eqr{7.61} we arrive at
\[
   \rho^{+}\bb^{+}=\rho^{+}\big(\ba^{+}-\vek{i}^{+}\big)
   -\dvg^{+}\bT^{+},
\]
i.e., moving $\dvg^{+}\bT^{+}$ and $\rho^{+}\vek{i}^{+}$ to the left, the
linear momentum balance
\begin{equation}
   \dvg^{+}\bT^{+}+\rho^{+}\big(\bb^{+}+\vek{i}^{+}\big)=\rho^{+}\ba^{+},
   \tag{7.62}\label{eq:7.62}
\end{equation}
where
\begin{equation}
   \vek{i}^{+}=\ddot\bc+2\bOm\big(\bv^{+}-\dot\bc\big)
   +\big(\dot\bOm-\bOm^{2}\big)\big(\bx^{+}-\bc\big)
   \tag{7.63}\label{eq:7.63}
\end{equation}
is the contribution to the net body force arising from the spin and
translation of the frame of $O^{+}$. This is typically referred to as the
\emph{inertial force}.

Three features of the symbol are worth fixing before it is used. It is an
\emph{acceleration}, of the same dimensions as $\ba^{+}$, which is why it
may stand beside $\bb^{+}$ inside the bracket of \eqr{7.62} and why that
bracket must be multiplied by $\rho^{+}$ before it becomes a force per
unit volume. Its sign is not a convention but is forced by the
substitution: $\bQ\ba$ was replaced by $\ba^{+}-\vek{i}^{+}$, and
carrying the resulting $-\rho^{+}\vek{i}^{+}$ from the right-hand side to
the left turned it into $+\rho^{+}\vek{i}^{+}$. And it has no star-free
counterpart. For $O$ himself the change of observer is the identity,
$\bQ=\I$ and $\bc=\bze$, so that $\bOm=\Ob$ and \eqr{7.63} returns
$\vek{i}=\bze$; \eqr{7.62} then collapses to \eqr{7.53}$_{3}$. The symbol
exists only because a second observer was introduced, and it measures
that observer and nothing else.

Notice, finally, what \eqr{7.63} does \emph{not} contain. There is no
$\bT$, no $\bb$, no material constant, nothing about what is being done
to the body: only $\bc$ and $\bOm$, which describe the relation between
the two frames, and $\bx^{+}$ and $\bv^{+}$, which say where the particle
is and how fast it is moving as $O^{+}$ records it. That is precisely the
property that licenses moving $\vek{i}^{+}$ to the force side of the
equation and treating it as though it were a body force, in the manner of
Remark~\ref{rem:ch7inertial}: $O^{+}$ can evaluate it once, as a field
over his whole frame, before he knows anything whatever about the body
that is going to occupy it. Figure~\ref{fig:ch7fourterms} is that field,
term by term.

However, if the frame of $O^{+}$ is related to
that of $O$ by a Galilean transformation, then $\bOm=\Ob$, $\dot\bOm=\Ob$
and $\ddot\bc=\bze$, so every term of \eqr{7.63} vanishes,
$\vek{i}^{+}=\bze$, and the balance of linear momentum is identical in
both frames, i.e.
\begin{equation}
   \dvg^{+}\bT^{+}+\rho^{+}\bb^{+}=\rho^{+}\ba^{+}.
   \tag{7.64}\label{eq:7.64}
\end{equation}
Thus, if an inertial frame exists, cf. Euler's postulates \eqr{6.34} and
\eqr{6.35}, then there are infinitely many, simply because there are
infinitely many Galilean transformations.

\begin{remark}[what an inertial force is, and what it is
not]\label{rem:ch7inertial}
Three things about \eqr{7.62} repay attention.

\emph{It is not a repair.} Nothing was added to the physics. Equation
\eqr{7.62} is \eqr{7.53}$_{3}$ rewritten, and every symbol in
$\vek{i}^{+}$ was computed in Section~\ref{sec:ch7motion} from the
relative motion of the two observers. The non-inertial observer is not
doing worse mechanics; he is doing the same mechanics with a longer
right-hand side, and he can compute the extra term himself once he knows
$\bQ(t)$ and $\bc(t)$.

\emph{Why it looks like a body force.} In \eqr{7.62} the inertial force
enters multiplied by $\rho^{+}$, exactly as $\bb^{+}$ does, and that is
not a notational accident: $\vek{i}^{+}$ is an acceleration, and
multiplying an acceleration by a density produces a force per unit
volume. The consequence is that no measurement made inside a small region
can distinguish a uniform $\ddot\bc$ from a uniform gravitational field,
since both act on every particle in proportion to its mass. That is the
observation from which general relativity starts. Within the present
theory it means only that the split of the right-hand side of \eqr{7.62}
into $\bb^{+}$ and $\vek{i}^{+}$ is a matter of bookkeeping unless the
frames are specified.

\emph{Why only this law fails.} Look back at \eqr{7.59}. The mass
balance contains no acceleration; the energy balance contains $\ba$
nowhere, only $\dot\eps$ and the stress power; the Clausius--Duhem
inequality likewise. Linear momentum is the only one of the five in which
the material acceleration appears, and $\ba$ is the only quantity in the
whole dictionary of Section~\ref{sec:ch7further} whose transformation law
carries $\ddot\bc$ and $\dot\bOm$. One non-objective ingredient, one
non-invariant law. The correspondence is exact, and it is the clearest
illustration in the book of Remark~\ref{rem:ch7objinv}.
\end{remark}

\begin{figure}[htbp]\centering
\renewcommand{\thefigure}{I}
\begin{tikzpicture}[scale=1.0,
   ar/.style={-{Stealth[length=2.2mm]},thick},
   sm/.style={-{Stealth[length=1.8mm]},semithick}]
% ================= panel (a): the inertial observer =================
\begin{scope}
  \draw[thick,gray!55] (0,0) circle (1.55);
  \fill (0,0) circle (1.3pt);
  \node[font=\footnotesize,text=gray!85,anchor=north east] at (-0.05,-0.06)
     {axis};
  % straight, uniform path of a free particle
  \draw[ar,blue!55!black] (-1.85,-0.95) -- (1.95,0.72);
  \foreach \s in {0,0.25,0.5,0.75,1}
     {\fill[blue!55!black]
        ($(-1.85,-0.95)!\s!(1.95,0.72)$) circle (1.4pt);}
  \node[font=\small,text=blue!55!black,anchor=south east] at (1.95,0.72)
     {$\bv$ constant};
  % the turntable's own rotation
  \draw[sm,gray!70] (1.10,-1.30) arc (-50:-14:1.70);
  \node[font=\footnotesize,text=gray!85,anchor=north west] at (1.12,-1.32)
     {$\bom_{p}$};
  \node[font=\small,anchor=north,align=center] at (0,-2.15)
     {(a) observer $O$, inertial\\[0.6mm]
      $\bb=\bze$ and $\ba=\bze$: no force, no bending};
\end{scope}
% ================= panel (b): the rotating observer =================
\begin{scope}[xshift=7.0cm]
  \draw[thick,gray!55] (0,0) circle (1.55);
  \fill (0,0) circle (1.3pt);
  % the same journey, recorded in the rotating frame: a curved track
  \draw[ar,red!65!black] plot[smooth,tension=0.85] coordinates
     {(-1.85,-0.95) (-0.95,-0.30) (-0.10,0.10) (0.80,0.10) (1.60,-0.35)};
  \foreach \p in {(-1.85,-0.95),(-0.95,-0.30),(-0.10,0.10),(0.80,0.10),
                  (1.60,-0.35)}
     {\fill[red!65!black] \p circle (1.4pt);}
  % Coriolis and centrifugal at one instant
  \coordinate (P) at (-0.10,0.10);
  \draw[sm,blue!55!black] (P) -- ++(0.86,0.10)
     node[right=0pt,font=\small,text=blue!55!black] {$\bv^{+}$};
  \draw[sm,red!65!black] (P) -- ++(-0.10,-0.80)
     node[below=0pt,font=\small,text=red!65!black]
     {$2\bOm\bv^{+}$};
  \draw[sm,gray!75] (P) -- ++(-0.62,0.34)
     node[above left=-3pt,font=\small,text=gray!85]
     {$-\bOm^{2}\bx^{+}$};
  \node[font=\small,anchor=north,align=center] at (0,-2.15)
     {(b) observer $O^{+}$, spinning with the disc\\[0.6mm]
      $\bb^{+}=\bze$ but $\ba^{+}\neq\bze$: the deficit is
      $\vek{i}^{+}$};
\end{scope}
\end{tikzpicture}
\caption{Why a rotating observer manufactures forces. A free particle
crosses a spinning disc. In (a) the inertial observer records a straight
path at constant speed, consistent with \eqr{7.64} and no force at all.
In (b) the observer turning with the disc records the same journey as a
curved track, so $\ba^{+}\neq\bze$; since the tractions and the true body
force are unchanged, \eqr{7.64} cannot hold in his frame. Equation
\eqr{7.62} restores the balance by crediting him with the inertial force
$\vek{i}^{+}$ of \eqr{7.63}, whose Coriolis part $2\bOm\bv^{+}$ bends the
track and whose centrifugal part $-\bOm^{2}\bx^{+}$ pushes outward.
Nothing touches the particle in either panel; the difference is entirely
in $\dot\bQ$.}
\label{fig:ch7apparent}
\end{figure}

It is possible to show that
\begin{equation}
   \ba^{++}-\vek{i}^{++}=\bQ^{+}(t^{+})\big(\ba^{+}-\vek{i}^{+}\big)
   \tag{7.65}\label{eq:7.65}
\end{equation}
under a further transformation
\begin{equation}
   \bx^{++}=\bQ^{+}(t^{+})\bx^{+}+\bc^{+}(t^{+}),
   \qquad
   t^{++}=t^{+}+a^{+}
   \tag{7.66}\label{eq:7.66}
\end{equation}
from the frame of $O^{+}$ to that of another observer $O^{++}$, where
$\vek{i}^{++}$ is the inertial force in the latter frame, and hence that
\begin{equation}
   \dvg^{++}\bT^{++}+\rho^{++}\big(\bb^{++}+\vek{i}^{++}\big)
   =\rho^{++}\ba^{++}.
   \tag{7.67}\label{eq:7.67}
\end{equation}
Thus the linear momentum balance, with the inertial forces included, is
valid in all frames. Moreover, the jump conditions derived in Chapter 6
are identical in all frames.

The proof of \eqr{7.65} is two observations and one line, and it is worth
giving, because it is what makes the theory consistent rather than merely
serviceable.

\emph{First, changes of observer compose.} Substituting \eqr{7.3} into
\eqr{7.66},
\[
   \bx^{++}=\bQ^{+}\big(\bQ\bx+\bc\big)+\bc^{+}
   =\big(\bQ^{+}\bQ\big)\bx+\big(\bQ^{+}\bc+\bc^{+}\big)
   =:\bQ^{++}\bx+\bc^{++},
\]
and $\bQ^{++}=\bQ^{+}\bQ$ is orthogonal as a product of orthogonal
tensors, mapping $\Et$ to $\mathbb{E}^{3++}$. Likewise
$t^{++}=t+(a+a^{+})$, again of the form \eqr{7.1}. So the passage from
$O$ directly to $O^{++}$ is itself a change of observer of exactly the
kind studied in Section~\ref{sec:ch7obs}, and everything proved for a
single change applies to it.

\emph{Second, \eqr{7.9} is the definition of the inertial force.}
Comparing \eqr{7.9} with \eqr{7.63}, the content of the two together is
\[
   \vek{i}^{+}=\ba^{+}-\bQ\ba ,
\]
the acceleration deficit of the frame of $O^{+}$ measured against the
inertial frame of $O$. Applying the same statement to the composite
change, with $\bQ^{++}$ in place of $\bQ$,
$\vek{i}^{++}=\ba^{++}-\bQ^{++}\ba$. Therefore
\[
   \ba^{++}-\vek{i}^{++}=\bQ^{++}\ba=\bQ^{+}\bQ\ba
   =\bQ^{+}\big(\ba^{+}-\vek{i}^{+}\big),
\]
which is \eqr{7.65}. And \eqr{7.67} is \eqr{7.62} written for the
composite change of observer, which is legitimate precisely because the
first observation showed the composite to be a change of observer.

\begin{remark}[spins add]\label{rem:ch7addspin}
The composition rule has a corollary worth recording, since it is the
statement that the $\bOm$ of \eqr{7.8} behaves like an angular velocity.
Differentiating $\bQ^{++}=\bQ^{+}\bQ$, all time derivatives being with
respect to the common time by \eqr{7.1} and \eqr{7.66}$_{2}$,
\[
   \bOm^{++}=\dot{\bQ^{++}}\big(\bQ^{++}\big)^{t}
   =\big(\dot\bQ^{+}\bQ+\bQ^{+}\dot\bQ\big)\bQ^{t}(\bQ^{+})^{t}
   =\dot\bQ^{+}(\bQ^{+})^{t}
     +\bQ^{+}\big(\dot\bQ\bQ^{t}\big)(\bQ^{+})^{t},
\]
that is, $\bOm^{++}=\bOm^{+}_{\text{rel}}+\bQ^{+}\bOm(\bQ^{+})^{t}$,
where $\bOm^{+}_{\text{rel}}=\dot\bQ^{+}(\bQ^{+})^{t}$ is the spin of the
second change taken by itself. The spin of the composite is the spin of
the second plus the spin of the first, the latter carried over by
$\bQ^{+}$ into the final frame. Written with axial vectors, by
Proposition~\ref{prop:axial}, this reads
$\bom^{++}=\bom^{+}_{\text{rel}}+\bQ^{+}\bom$: angular velocities add.
Nothing in this chapter would work if they did not, since \eqr{7.65}
would then depend on the route taken from $O$ to $O^{++}$.
\end{remark}

\begin{remark}[what the chapter licenses, and what it
forbids]\label{rem:ch7mfi}
It is worth setting down in one place the principle that the whole
chapter exists to make precise, because Chapter 8 will invoke it on every
page and will not repeat the derivation.

\emph{The principle.} A constitutive equation is a claim about a
material, and a material has no access to the identity of the observer
looking at it. A constitutive equation must therefore survive the
substitution of \eqr{7.3} for arbitrary $\bQ(t)$ and $\bc(t)$: if it
holds between the quantities recorded by $O$, then the same equation with
every symbol starred must hold between the quantities recorded by
$O^{+}$. This is the \emph{principle of material frame indifference}. It
is a restriction on the constitutive equation and never on the motion:
the motion is whatever it is, and it is the form of the law that is being
constrained.

\emph{Why it is not vacuous.} The dictionary of
Section~\ref{sec:ch7further} is the list of what each variable does under
that substitution, and its entries are not alike. A law written in
variables drawn from the first four groups satisfies the principle
automatically, because the $\bQ$'s cancel in pairs or never appear at
all. A law containing any entry from the last group does not, because
$\dot\bQ$ enters with nothing to cancel against and $\dot\bQ$ is
arbitrary. The principle is therefore a genuine sieve, and the dictionary
is that sieve written out.

\emph{The three prohibitions established here.} Each is a mistake that
has been made.
\begin{itemize}\itemsep2pt
\item No constitutive law may contain $\bW$, nor $\bL$ otherwise than
      through $\bD$. Remark~\ref{rem:ch7DW}: the observer manufactures
      spin out of nothing, so such a law would predict stress in a body
      lying still on a table.
\item No constitutive law may contain $\bv$ or $\ba$. By \eqr{7.8} and
      \eqr{7.9} the two observers disagree about both even when the body
      is doing nothing whatever.
\item No constitutive law may contain $\bF$ except through $\bC$, or
      equivalently none may see $\bR$. By \eqr{7.29}$_{3}$ the rotation
      is exactly the part of $\bF$ that a change of observer alters.
\end{itemize}
The corresponding licences are the first four groups of the dictionary,
and the safest of them is the second: a law relating $\bS$ to $\bC$
involves nothing that sees $\bQ$ at all, which is why constitutive theory
is written there.

\emph{One thing the principle does not govern.} It says nothing about
the balance laws, which are not constitutive statements.
Section~\ref{sec:ch7balance} showed that four of the five are form
invariant on their own and that the fifth is repaired by \eqr{7.62}. A
frequent confusion is to read the failure of \eqr{7.64} in a rotating
frame as a failure of frame indifference. It is not. Frame indifference
is a demand made on constitutive equations; \eqr{7.62} is a theorem about
a balance law; and the theorem tells the engineer in the rotating frame
which extra terms to add, not that he must stop.

\emph{How strong the principle is.} What has been established in this
chapter is that the theory is consistent with the principle and what the
principle sieves out. Whether the principle is exactly true is a separate
question, and it has been argued over since M\"uller and Truesdell: for a
gas whose molecular mean free path is not negligible against the scale of
the flow, kinetic theory does predict a weak dependence of the
constitutive response on $\bOm$, and the principle is then an excellent
approximation rather than an axiom. Nothing in this book is at that
scale, and the principle is used at full strength throughout; but it is
an assumption about materials, not a theorem about geometry, and it is
worth knowing which of the two one is holding. Compare
Remark~\ref{rem:ch7tractionass}, where the same point was made about
\eqr{7.38}$_{1}$: the geometry of Section~\ref{sec:ch7obs} is proved, and
everything in this chapter that concerns matter rather than space is
assumed.
\end{remark}

%=======================================================================
\setcounter{problemenv}{0}
\section{Problems}\label{sec:problems:c7}
%=======================================================================

%-----------------------------------------------------------------------
\begin{problem}[The transformation of the acceleration]
Derive \eqr{7.9}.
\end{problem}

\begin{solution}
The acceleration in the frame of $O^{+}$ is
$\ba^{+}=\partial\bv^{+}/\partial t^{+}$ at fixed $p$, by
Definition~\ref{def:va}, and $\partial/\partial t^{+}=\partial/\partial
t$ by the argument given before \eqr{7.7}. So the whole problem is to
differentiate \eqr{7.7} once more and then to eliminate the unstarred
quantities $\bx$ and $\bv$ in favour of the starred ones. Neither
operation is difficult; what makes the result look complicated is that
both $\bx$ and $\bv$ have to be eliminated, and the second elimination
feeds on the first. The elimination is therefore split into two steps
below, giving four in all.

\textbf{Step 1: differentiate.} From \eqr{7.7},
$\bv^{+}=\dot\bQ\bx+\bQ\bv+\dot\bc$. Differentiate at fixed $p$, using
the product rule of Problem 1(a) of Chapter 4 on the first two terms and
remembering that $\dot\bx=\bv$ and $\dot\bv=\ba$ because $\bx$ and $\bv$
are being followed along the motion of the material point:
\[
   \ba^{+}=\ddot\bQ\bx+\dot\bQ\bv+\dot\bQ\bv+\bQ\ba+\ddot\bc
   =\ddot\bQ\bx+2\dot\bQ\bv+\bQ\ba+\ddot\bc .
\]
The two middle terms are equal and are the origin of the factor $2$ in
the Coriolis term: one copy comes from differentiating $\dot\bQ\bx$ and
the other from differentiating $\bQ\bv$. Rearranged,
\begin{equation*}
   \ba^{+}-\bQ\ba=\ddot\bc+2\dot\bQ\bv+\ddot\bQ\bx .
\end{equation*}
Everything now depends on rewriting the last two terms.

\textbf{Step 2: eliminate $\bx$.} From \eqr{7.3},
$\bQ\bx=\bx^{+}-\bc$. Also, from $\bOm=\dot\bQ\bQ^{t}$ and
$\bQ^{t}\bQ=\I$,
\[
   \dot\bQ=\dot\bQ\,\I=\dot\bQ\bQ^{t}\bQ=\bOm\bQ ,
\]
which is the identity that will be used three times. Differentiating it,
\[
   \ddot\bQ=\big(\bOm\bQ\big)^{\textstyle\cdot}
   =\dot\bOm\bQ+\bOm\dot\bQ=\dot\bOm\bQ+\bOm\big(\bOm\bQ\big)
   =\big(\dot\bOm+\bOm^{2}\big)\bQ .
\]
Hence
\[
   \ddot\bQ\bx=\big(\dot\bOm+\bOm^{2}\big)\bQ\bx
   =\big(\dot\bOm+\bOm^{2}\big)\big(\bx^{+}-\bc\big).
\]

\textbf{Step 3: eliminate $\bv$.} Solve \eqr{7.7} for $\bQ\bv$:
\[
   \bQ\bv=\bv^{+}-\dot\bQ\bx-\dot\bc
   =\bv^{+}-\dot\bc-\bOm\big(\bx^{+}-\bc\big),
\]
using the result of Step 2 for $\dot\bQ\bx=\bOm\bQ\bx
=\bOm(\bx^{+}-\bc)$. This is \eqr{7.8} rearranged, as it must be.
Therefore
\[
\begin{aligned}
   2\dot\bQ\bv&=2\bOm\bQ\bv
    =2\bOm\Big[\bv^{+}-\dot\bc-\bOm\big(\bx^{+}-\bc\big)\Big]\\
   &=2\bOm\big(\bv^{+}-\dot\bc\big)-2\bOm^{2}\big(\bx^{+}-\bc\big).
\end{aligned}
\]

\textbf{Step 4: collect.} Substituting the results of Steps 2 and 3 into
the display at the end of Step 1,
\[
\begin{aligned}
   \ba^{+}-\bQ\ba&=\ddot\bc
     +2\bOm\big(\bv^{+}-\dot\bc\big)-2\bOm^{2}\big(\bx^{+}-\bc\big)
     +\big(\dot\bOm+\bOm^{2}\big)\big(\bx^{+}-\bc\big)\\
   &=\ddot\bc+2\bOm\big(\bv^{+}-\dot\bc\big)
     +\big(-2\bOm^{2}+\dot\bOm+\bOm^{2}\big)\big(\bx^{+}-\bc\big),
\end{aligned}
\]
the two brackets in $\bx^{+}-\bc$ having been collected by linearity.
Since $-2\bOm^{2}+\bOm^{2}=-\bOm^{2}$,
\[
   \boxed{\ \ba^{+}-\bQ\ba=\ddot\bc+2\bOm\big(\bv^{+}-\dot\bc\big)
   +\big(\dot\bOm-\bOm^{2}\big)\big(\bx^{+}-\bc\big),\ }
\]
which is \eqr{7.9}.

\emph{Two checks.} First, put $\dot\bQ=\Ob$ and $\ddot\bc=\bze$. Then
$\bOm=\Ob$, $\dot\bOm=\Ob$, and every term on the right vanishes, leaving
$\ba^{+}=\bQ\ba$: this is \eqr{7.10}$_{2}$, so the Galilean case is
recovered. Second, put $\bv^{+}=\dot\bc$ and $\bx^{+}=\bc$, meaning a
particle sitting at the origin of $O^{+}$'s frame and moving with it.
Then only $\ddot\bc$ survives, which is right: such a particle has, for
$O^{+}$, no acceleration at all, and the whole discrepancy is the
acceleration of the frame itself.

\emph{Where the minus sign comes from.} The centrifugal term appears as
$-\bOm^{2}$ and not $+\bOm^{2}$ because the two contributions to
$\bOm^{2}$ have opposite signs and different origins: $+\bOm^{2}$ comes
from $\ddot\bQ$, the turning of the frame's turning, while $-2\bOm^{2}$
comes from the Coriolis term acting on the part of $\bv$ that is itself
due to the frame's rotation. Their sum is $-\bOm^{2}$. With
$\bOm\bs=\bom\times\bs$ this is $-\bom\times(\bom\times\bs)$, which by
the expansion of the triple cross product is
$|\bom|^{2}\bs-\ip\bom\bs\,\bom$, the component of $|\bom|^{2}\bs$
perpendicular to the axis. It points away from the axis, as centrifugal
effects do.

The result is what \eqr{7.62} and \eqr{7.63} are built on: the entire
difference between the two observers' equations of motion is the boxed
expression multiplied by $\rho$.
\end{solution}

%-----------------------------------------------------------------------
\begin{problem}[The rate of the shifter in a spinning frame]
Show that $(\shift^{+})^{\textstyle\cdot}=\bOm\shift^{+}$, where
$\bOm=\dot\bQ\bQ^{t}$.
\end{problem}

\begin{solution}
Everything is in \eqr{7.18}, $\shift^{+}=\bQ(t)\shift\bK^{t}$, once it is
noticed which of the three factors can move.

The shifter $\shift$ of the observer $O$ is fixed, because the frame of
$O$ is taken to be inertial, as stated after \eqr{7.18}; so
$\dot\shift=\Ob$ in the sense that $\shift$ is a constant two-point
tensor. The tensor $\bK$ is fixed by \eqr{7.15}, every factor in it
having been evaluated at the single instant $t_{0}$; so $\dot\bK=\hOb$
and hence $\dot{\bK^{t}}=(\dot\bK)^{t}=\hOb$ by Problem 1(b) of
Chapter 4. Only $\bQ(t)$ moves. Differentiating \eqr{7.18} with the
product rule, two of the three terms therefore vanish:
\[
   \big(\shift^{+}\big)^{\textstyle\cdot}
   =\dot\bQ\shift\bK^{t}+\bQ\dot\shift\bK^{t}
     +\bQ\shift\dot{\bK^{t}}
   =\dot\bQ\shift\bK^{t}.
\]
Now insert $\I=\bQ^{t}\bQ$, which is legitimate by
Proposition~\ref{prop:orth}, and regroup so that \eqr{7.18} reappears on
the right:
\[
   \dot\bQ\shift\bK^{t}=\dot\bQ\,\I\,\shift\bK^{t}
   =\big(\dot\bQ\bQ^{t}\big)\big(\bQ\shift\bK^{t}\big)
   =\bOm\shift^{+},
\]
using \eqr{7.8}$_{2}$ for the first bracket and \eqr{7.18} for the
second. Hence
\[
   \boxed{\ \big(\shift^{+}\big)^{\textstyle\cdot}=\bOm\shift^{+},
   \qquad \bOm=\dot\bQ\bQ^{t}.\ }
\]

\emph{What the formula says.} Compare it with \eqr{4.3},
$\dot\bF=\bL\bF$. The two are the same statement with $\shift^{+}$ in
place of $\bF$ and $\bOm$ in place of $\bL$. That is exactly right, and
it is Remark~\ref{rem:ch7omega} in another guise: the frame of $O^{+}$
undergoes a rigid displacement relative to that of $O$, whose velocity
gradient is $\bOm$ by \eqr{4.29}, and $\shift^{+}$ is the two-point
tensor that gets carried along by it. The shifter of a spinning observer
turns with him, at the rate $\bOm$.

\emph{The consistency check that matters.} The relation
$\bH^{+}=\bF^{+}-\shift^{+}$ of \eqr{7.25}$_{2}$ would be worthless if
its two terms evolved by different rules. They do not:
$\dot\bF^{+}=\bL^{+}\bF^{+}$ by \eqr{4.3}, and for a frame at rest
relative to the material, so that $\bL^{+}$ reduces to the frame's own
contribution $\bOm$ by \eqr{7.33}, the two rates coincide. In particular,
if the body is at rest for $O$ then $\bF$ is constant, and
\[
   \dot\bF^{+}=\big(\bQ\bF\bK^{t}\big)^{\textstyle\cdot}
   =\dot\bQ\bF\bK^{t}=\bOm\bF^{+},
\]
by the same insertion of $\I=\bQ^{t}\bQ$, so
$\dot\bH^{+}=\bOm(\bF^{+}-\shift^{+})=\bOm\bH^{+}$ and the displacement
gradient simply rides round with the observer. Had the shifter been
replaced by $\hI$, its counterpart $\hI^{+}=\bK\hI\bK^{t}$ of \eqr{7.26}
would have satisfied $(\hI^{+})^{\textstyle\cdot}=\hOb^{+}$ instead: it
would have stood still while $\bF^{+}$ turned, and the difference
$\bF^{+}-\hI^{+}$ would have changed in time even though nothing whatever
was happening to the body. That is the dynamic version of the objection
raised after \eqr{7.26}, and it is the sharper of the two.

\emph{When it vanishes.} If $O^{+}$ is related to $O$ by a Galilean
transformation then $\dot\bQ=\Ob$, so $\bOm=\Ob$ and
$(\shift^{+})^{\textstyle\cdot}=\Ob$: the shifter is fixed in every
inertial frame, which is the claim made in
Remark~\ref{rem:ch7whyshiftass}.
\end{solution}

%-----------------------------------------------------------------------
\begin{problem}[An observer standing on a spinning planet]
Observer $O^{+}$ adopts a frame of reference attached to the surface of a
planet of radius $R$, with origin located at azimuth $\theta$ (degrees
longitude) and elevation $\varphi$ (degrees latitude). The planet spins
about its axis, the unit vector $\be_{3}$, at the constant rate $\omega$,
i.e. $\theta=\omega t+\theta_{0}$ with $\theta_{0}=\theta(0)$. What are
the equations of motion for this observer?
\end{problem}

\begin{solution}
Required: \eqr{7.62} with \eqr{7.63}, i.e. $\bc,\bOm,\dot\bOm$ from
$\bQ(t),\bc(t)$ of \eqr{7.3}. Geometry: Figures~\ref{fig:ch7planet}
and~\ref{fig:ch7planet3d}. Data:
\begin{gather*}
   \ip{\be_{i}}{\be_{j}}=\dd_{ij},\qquad
   \be_{1}\times\be_{2}=\be_{3},\qquad \be_{2}\times\be_{3}=\be_{1},\qquad
   \be_{3}\times\be_{1}=\be_{2},\\
   \theta=\omega t+\theta_{0},\qquad \dot\theta=\omega,\qquad
   \dot\varphi=0,\qquad \dot\be_{i}=\bze .
\end{gather*}
Equatorial plane (Figure~\ref{fig:ch7planet3d}): unit vector at angle
$\theta$ from $\be_{1}$, and the same at $\theta+\pi/2$:
\begin{align*}
   \bff_{\rho}&=\cos\theta\,\be_{1}+\sin\theta\,\be_{2},\\
   \bff_{E}&=\cos(\theta+\tfrac{\pi}{2})\,\be_{1}
            +\sin(\theta+\tfrac{\pi}{2})\,\be_{2}
           =-\sin\theta\,\be_{1}+\cos\theta\,\be_{2},\\
   \be_{3}\times\bff_{\rho}
     &=\cos\theta\,\be_{3}\times\be_{1}+\sin\theta\,\be_{3}\times\be_{2}
      =\cos\theta\,\be_{2}-\sin\theta\,\be_{1}=\bff_{E},\\
   |\bff_{\rho}|^{2}&=\cos^{2}\theta+\sin^{2}\theta=1,\qquad
   |\bff_{E}|^{2}=\sin^{2}\theta+\cos^{2}\theta=1,\\
   \ip{\bff_{\rho}}{\bff_{E}}
     &=-\cos\theta\sin\theta+\sin\theta\cos\theta=0,\qquad
   \ip{\bff_{\rho}}{\be_{3}}=\ip{\bff_{E}}{\be_{3}}=0,\\
   \bff_{\rho}\times\bff_{E}
     &=\cos^{2}\theta\,\be_{1}\times\be_{2}
       -\sin^{2}\theta\,\be_{2}\times\be_{1}
      =(\cos^{2}\theta+\sin^{2}\theta)\,\be_{3}=\be_{3},\\
   \bff_{E}\times\be_{3}
     &=-\sin\theta\,\be_{1}\times\be_{3}+\cos\theta\,\be_{2}\times\be_{3}
      =\sin\theta\,\be_{2}+\cos\theta\,\be_{1}=\bff_{\rho}.
\end{align*}
Meridian plane (Figure~\ref{fig:ch7planet}(a)), spanned by
$\bff_{\rho},\be_{3}$: unit vector at angle $\varphi$ from
$\bff_{\rho}$, and the same at $\varphi+\pi/2$:
\begin{align*}
   \bff_{U}&=\cos\varphi\,\bff_{\rho}+\sin\varphi\,\be_{3},
   \qquad \bo^{+}=R\,\bff_{U},\\
   \bff_{N}&=\cos(\varphi+\tfrac{\pi}{2})\,\bff_{\rho}
            +\sin(\varphi+\tfrac{\pi}{2})\,\be_{3}
           =-\sin\varphi\,\bff_{\rho}+\cos\varphi\,\be_{3},\\
   |\bff_{U}|^{2}&=\cos^{2}\varphi\,|\bff_{\rho}|^{2}
      +2\cos\varphi\sin\varphi\,\ip{\bff_{\rho}}{\be_{3}}
      +\sin^{2}\varphi\,|\be_{3}|^{2}\\
      &=\cos^{2}\varphi+0+\sin^{2}\varphi=1,\\
   |\bff_{N}|^{2}&=\sin^{2}\varphi\,|\bff_{\rho}|^{2}
      -2\sin\varphi\cos\varphi\,\ip{\bff_{\rho}}{\be_{3}}
      +\cos^{2}\varphi\,|\be_{3}|^{2}\\
      &=\sin^{2}\varphi+0+\cos^{2}\varphi=1,\\
   \ip{\bff_{N}}{\bff_{U}}
     &=-\sin\varphi\cos\varphi\,|\bff_{\rho}|^{2}
       +(\cos^{2}\varphi-\sin^{2}\varphi)\ip{\bff_{\rho}}{\be_{3}}
       +\cos\varphi\sin\varphi\,|\be_{3}|^{2}\\
     &=-\sin\varphi\cos\varphi+0+\cos\varphi\sin\varphi=0,\\
   \ip{\bff_{E}}{\bff_{U}}
     &=\cos\varphi\ip{\bff_{E}}{\bff_{\rho}}
       +\sin\varphi\ip{\bff_{E}}{\be_{3}}=0,\\
   \ip{\bff_{E}}{\bff_{N}}
     &=-\sin\varphi\ip{\bff_{E}}{\bff_{\rho}}
       +\cos\varphi\ip{\bff_{E}}{\be_{3}}=0,\\
   \bff_{E}\times\bff_{N}
     &=-\sin\varphi\,\bff_{E}\times\bff_{\rho}
       +\cos\varphi\,\bff_{E}\times\be_{3}
      =\sin\varphi\,\be_{3}+\cos\varphi\,\bff_{\rho}=\bff_{U},\\
   \tbox{\bff_{E}}{\bff_{N}}{\bff_{U}}
     &=\ip{\bff_{E}\times\bff_{N}}{\bff_{U}}
      =\ip{\bff_{U}}{\bff_{U}}=1 .
\end{align*}
Hence, for $\alpha,\beta\in\{E,N,U\}$,
$\ip{\bff_{\alpha}}{\bff_{\beta}}=\dd_{\alpha\beta}$ and
$\bff_{E}\times\bff_{N}=\bff_{U}$,
$\bff_{N}\times\bff_{U}=\bff_{E}$,
$\bff_{U}\times\bff_{E}=\bff_{N}$. Inverse of the last two
definitions: $\cos\varphi\times$first $-\sin\varphi\times$second, and
$\sin\varphi\times$first $+\cos\varphi\times$second,
\begin{align*}
   \cos\varphi\,\bff_{U}-\sin\varphi\,\bff_{N}
     &=(\cos^{2}\varphi+\sin^{2}\varphi)\,\bff_{\rho}
       +(\cos\varphi\sin\varphi-\sin\varphi\cos\varphi)\,\be_{3}
      =\bff_{\rho},\\
   \sin\varphi\,\bff_{U}+\cos\varphi\,\bff_{N}
     &=(\sin\varphi\cos\varphi-\cos\varphi\sin\varphi)\,\bff_{\rho}
       +(\sin^{2}\varphi+\cos^{2}\varphi)\,\be_{3}
      =\be_{3},
\end{align*}
and in the equatorial plane in the same way
\begin{align*}
   \cos\theta\,\bff_{\rho}-\sin\theta\,\bff_{E}
     &=(\cos^{2}\theta+\sin^{2}\theta)\,\be_{1}
       +(\cos\theta\sin\theta-\sin\theta\cos\theta)\,\be_{2}=\be_{1},\\
   \sin\theta\,\bff_{\rho}+\cos\theta\,\bff_{E}
     &=(\sin\theta\cos\theta-\cos\theta\sin\theta)\,\be_{1}
       +(\sin^{2}\theta+\cos^{2}\theta)\,\be_{2}=\be_{2}.
\end{align*}

$O^{+}$: fixed basis $\{\be^{+}_{E},\be^{+}_{N},\be^{+}_{U}\}$ of
$\mathbb{E}^{3+}$, $\ip{\be^{+}_{\alpha}}{\be^{+}_{\beta}}
=\dd_{\alpha\beta}$, $\be^{+}_{E}\times\be^{+}_{N}=\be^{+}_{U}$
cyclically. $\bQ$ of \eqr{7.3} with $\bQ\bff_{\alpha}=\be^{+}_{\alpha}$,
by Lemma~\ref{lem:dyadprod}:
\begin{align*}
   \bQ&=\be^{+}_{\alpha}\otimes\bff_{\alpha},\qquad
   \bQ^{t}=\bff_{\alpha}\otimes\be^{+}_{\alpha},\qquad
   \bQ\bff_{\beta}=\be^{+}_{\alpha}\ip{\bff_{\alpha}}{\bff_{\beta}}
     =\be^{+}_{\alpha}\dd_{\alpha\beta}=\be^{+}_{\beta},\\
   \bQ\bQ^{t}&=\big(\be^{+}_{\alpha}\otimes\bff_{\alpha}\big)
     \big(\bff_{\beta}\otimes\be^{+}_{\beta}\big)
   =\ip{\bff_{\alpha}}{\bff_{\beta}}\,
     \be^{+}_{\alpha}\otimes\be^{+}_{\beta}\\
   &=\dd_{\alpha\beta}\,\be^{+}_{\alpha}\otimes\be^{+}_{\beta}
   =\be^{+}_{\alpha}\otimes\be^{+}_{\alpha}=\I^{+},\\
   \bQ^{t}\bQ&=\big(\bff_{\alpha}\otimes\be^{+}_{\alpha}\big)
     \big(\be^{+}_{\beta}\otimes\bff_{\beta}\big)
   =\ip{\be^{+}_{\alpha}}{\be^{+}_{\beta}}\,
     \bff_{\alpha}\otimes\bff_{\beta}\\
   &=\dd_{\alpha\beta}\,\bff_{\alpha}\otimes\bff_{\beta}
   =\bff_{\alpha}\otimes\bff_{\alpha}=\I .
\end{align*}
Station at the origin of $\mathbb{E}^{3+}$:
\begin{equation*}
   \bze=\bQ\bo^{+}+\bc=R\,\bQ\bff_{U}+\bc=R\,\be^{+}_{U}+\bc
   \ \Longrightarrow\
   \bc=-R\,\be^{+}_{U},\qquad
   \dot\bc=\ddot\bc=\bze .
\end{equation*}
Rates, with $\dot\theta=\omega$, $\dot\varphi=0$, $\dot\be_{i}=\bze$,
and the inverse relation for $\bff_{\rho}$:
\begin{align*}
   \dot\bff_{\rho}&=\dot\theta\big(-\sin\theta\,\be_{1}
      +\cos\theta\,\be_{2}\big)=\omega\,\bff_{E},\\
   \dot\bff_{E}&=\dot\theta\big(-\cos\theta\,\be_{1}
      -\sin\theta\,\be_{2}\big)=-\omega\,\bff_{\rho}
      =-\omega\cos\varphi\,\bff_{U}+\omega\sin\varphi\,\bff_{N},\\
   \dot\bff_{N}&=-\sin\varphi\,\dot\bff_{\rho}+\cos\varphi\,\dot\be_{3}
      =-\omega\sin\varphi\,\bff_{E},\\
   \dot\bff_{U}&=\cos\varphi\,\dot\bff_{\rho}+\sin\varphi\,\dot\be_{3}
      =\omega\cos\varphi\,\bff_{E}.
\end{align*}
Spin, \eqr{7.8}$_{2}$ and Lemma~\ref{lem:dyadprod}:
\begin{align*}
   \dot\bQ&=\be^{+}_{\alpha}\otimes\dot\bff_{\alpha},\qquad
   \bOm=\dot\bQ\bQ^{t}
   =\big(\be^{+}_{\alpha}\otimes\dot\bff_{\alpha}\big)
    \big(\bff_{\beta}\otimes\be^{+}_{\beta}\big)\\
   &=\ip{\dot\bff_{\alpha}}{\bff_{\beta}}\,
     \be^{+}_{\alpha}\otimes\be^{+}_{\beta}
   =\Omega_{\alpha\beta}\,\be^{+}_{\alpha}\otimes\be^{+}_{\beta},\\
   \Omega_{\alpha\beta}&=\ip{\dot\bff_{\alpha}}{\bff_{\beta}}:
   \qquad
   \Omega_{EE}=\ip{\dot\bff_{E}}{\bff_{E}}=0,\qquad
   \Omega_{EN}=\ip{\dot\bff_{E}}{\bff_{N}}=\omega\sin\varphi,\\
   \Omega_{EU}&=\ip{\dot\bff_{E}}{\bff_{U}}=-\omega\cos\varphi,\qquad
   \Omega_{NE}=\ip{\dot\bff_{N}}{\bff_{E}}=-\omega\sin\varphi,\\
   \Omega_{UE}&=\ip{\dot\bff_{U}}{\bff_{E}}=\omega\cos\varphi,\qquad
   \Omega_{NN}=\Omega_{NU}=\Omega_{UN}=\Omega_{UU}=0,\\
   [\Omega_{\alpha\beta}]&=\omega
   \begin{pmatrix}
      0 & \sin\varphi & -\cos\varphi\\
      -\sin\varphi & 0 & 0\\
      \cos\varphi & 0 & 0
   \end{pmatrix}
   =-[\Omega_{\alpha\beta}]^{t},\\
   \Omega_{\alpha\beta}+\Omega_{\beta\alpha}
   &=\ip{\dot\bff_{\alpha}}{\bff_{\beta}}
    +\ip{\bff_{\alpha}}{\dot\bff_{\beta}}
   =\tfrac{\dl}{\dl t}\ip{\bff_{\alpha}}{\bff_{\beta}}
   =\tfrac{\dl}{\dl t}\dd_{\alpha\beta}=0 .
\end{align*}
Axial vector, Proposition~\ref{prop:axial} with $(E,N,U)$ as $(1,2,3)$:
\begin{align*}
   w_{\gamma}&=-\tfrac12\eps_{\gamma\alpha\beta}\Omega_{\alpha\beta},\qquad
   w_{E}=-\tfrac12(\Omega_{NU}-\Omega_{UN})=0,\\
   w_{N}&=-\tfrac12(\Omega_{UE}-\Omega_{EU})
        =-\tfrac12(\omega\cos\varphi+\omega\cos\varphi)=-\omega\cos\varphi,\\
   w_{U}&=-\tfrac12(\Omega_{EN}-\Omega_{NE})
        =-\tfrac12(\omega\sin\varphi+\omega\sin\varphi)=-\omega\sin\varphi,\\
   \bw&=-\omega\big(\cos\varphi\,\be^{+}_{N}+\sin\varphi\,\be^{+}_{U}\big)
      =:-\bom,\\
   \bOm\bs&=\bw\times\bs=-\bom\times\bs
      \qquad\forall\,\bs\in\mathbb{E}^{3+},\\
   \bQ\be_{3}&=\bQ\big(\sin\varphi\,\bff_{U}+\cos\varphi\,\bff_{N}\big)
      =\sin\varphi\,\be^{+}_{U}+\cos\varphi\,\be^{+}_{N}
      =\bom/\omega,\\
   \dot\bom&=\bze\ \Longrightarrow\ \dot\bOm=\Ob^{+}.
\end{align*}
Sign: \eqr{7.8} with $\bv=\bze$, $\dot\bc=\bze$ gives
$\bv^{+}=\bOm(\bx^{+}-\bc)=-\bom\times(\bx^{+}-\bc)$, a point at rest for
$O$ circles the axis backwards for $O^{+}$. $\theta_{0}$ enters $\bQ$
only, not $\bOm$.

\eqr{7.63} with $\ddot\bc=\dot\bc=\bze$, $\dot\bOm=\Ob^{+}$,
$\bs^{+}:=\bx^{+}-\bc=\bx^{+}+R\,\be^{+}_{U}$:
\begin{align*}
   \vek{i}^{+}&=2\bOm\bv^{+}-\bOm^{2}\bs^{+},\qquad
   \bOm\bv^{+}=-\bom\times\bv^{+},\\
   \bOm^{2}\bs^{+}&=\bOm(\bOm\bs^{+})=\bOm(-\bom\times\bs^{+})
   =-\bom\times(-\bom\times\bs^{+})=\bom\times(\bom\times\bs^{+}),
\end{align*}
\begin{equation*}
   \boxed{\
   \begin{aligned}
      &\dvg^{+}\bT^{+}+\rho^{+}\big(\bb^{+}+\vek{i}^{+}\big)
        =\rho^{+}\ba^{+},\\[0.6mm]
      &\vek{i}^{+}=-2\,\bom\times\bv^{+}
        -\bom\times\big(\bom\times\bs^{+}\big),
      \qquad
      \bom=\omega\big(\cos\varphi\,\be^{+}_{N}+\sin\varphi\,\be^{+}_{U}\big).
   \end{aligned}
   \ }
\end{equation*}
\eqr{7.59}, \eqr{7.60} and the jump conditions are unchanged.

Centrifugal term, by $\ba\times(\bb\times\bc)=\ip{\ba}{\bc}\bb
-\ip{\ba}{\bb}\bc$ and
$|\bom|^{2}=\omega^{2}(\cos^{2}\varphi+\sin^{2}\varphi)=\omega^{2}$:
\begin{equation*}
   -\bom\times(\bom\times\bs^{+})
   =-\ip{\bom}{\bs^{+}}\,\bom+\ip{\bom}{\bom}\,\bs^{+}
   =\omega^{2}\bs^{+}-\ip{\bom}{\bs^{+}}\,\bom ,
\end{equation*}
$\omega^{2}$ times the part of $\bs^{+}$ perpendicular to $\bom$. At the
station $\bx^{+}=\bze$, $\bs^{+}=R\,\be^{+}_{U}$,
$\ip{\bom}{\be^{+}_{U}}=\omega\sin\varphi$:
\begin{align*}
   -\bom\times(\bom\times R\,\be^{+}_{U})
   &=\omega^{2}R\,\be^{+}_{U}
     -R\,\omega\sin\varphi\cdot\omega\big(\cos\varphi\,\be^{+}_{N}
     +\sin\varphi\,\be^{+}_{U}\big)\\
   &=\omega^{2}R\,(1-\sin^{2}\varphi)\,\be^{+}_{U}
     -\omega^{2}R\sin\varphi\cos\varphi\,\be^{+}_{N}\\
   &=\omega^{2}R\cos\varphi
     \big(\cos\varphi\,\be^{+}_{U}-\sin\varphi\,\be^{+}_{N}\big),\\
   \big|\cdots\big|&=\omega^{2}R\cos\varphi
     \sqrt{\cos^{2}\varphi+\sin^{2}\varphi}=\omega^{2}R\cos\varphi,\\
   \cos\varphi\,\be^{+}_{U}-\sin\varphi\,\be^{+}_{N}
   &=\bQ\big(\cos\varphi\,\bff_{U}-\sin\varphi\,\bff_{N}\big)
    =\bQ\bff_{\rho},
\end{align*}
$\omega^{2}\times$ the distance $R\cos\varphi$ from the axis, along the
outward horizontal of the equatorial plane
(Figure~\ref{fig:ch7planet}(a)). With gravity $-g\,\be^{+}_{U}$,
\begin{align*}
   \bg_{\mathrm{eff}}
   &=-g\,\be^{+}_{U}+\omega^{2}R\cos\varphi
     \big(\cos\varphi\,\be^{+}_{U}-\sin\varphi\,\be^{+}_{N}\big)\\
   &=-\big(g-\omega^{2}R\cos^{2}\varphi\big)\be^{+}_{U}
     -\omega^{2}R\sin\varphi\cos\varphi\,\be^{+}_{N},\\
   \tan\delta&=\frac{|\text{horizontal part}|}{|\text{vertical part}|}
   =\frac{\omega^{2}R\sin\varphi\cos\varphi}
         {g-\omega^{2}R\cos^{2}\varphi},
\end{align*}
Figure~\ref{fig:ch7planet}(b),(c), Remark~\ref{rem:ch7inertial}: the
term is uniform and $\propto\rho^{+}$ like gravity, hence
indistinguishable from it locally.

Coriolis term, $\bv^{+}=v_{E}\be^{+}_{E}+v_{N}\be^{+}_{N}+v_{U}\be^{+}_{U}$:
\begin{align*}
   \be^{+}_{N}\times\bv^{+}
   &=v_{E}\,\be^{+}_{N}\times\be^{+}_{E}+v_{U}\,\be^{+}_{N}\times\be^{+}_{U}
    =-v_{E}\,\be^{+}_{U}+v_{U}\,\be^{+}_{E},\\
   \be^{+}_{U}\times\bv^{+}
   &=v_{E}\,\be^{+}_{U}\times\be^{+}_{E}+v_{N}\,\be^{+}_{U}\times\be^{+}_{N}
    =v_{E}\,\be^{+}_{N}-v_{N}\,\be^{+}_{E},\\
   -2\,\bom\times\bv^{+}
   &=-2\omega\cos\varphi\,\be^{+}_{N}\times\bv^{+}
     -2\omega\sin\varphi\,\be^{+}_{U}\times\bv^{+}\\
   &=-2\omega\cos\varphi\big(v_{U}\,\be^{+}_{E}-v_{E}\,\be^{+}_{U}\big)
     -2\omega\sin\varphi\big(-v_{N}\,\be^{+}_{E}+v_{E}\,\be^{+}_{N}\big)\\
   &=2\omega\Big[\big(v_{N}\sin\varphi-v_{U}\cos\varphi\big)\be^{+}_{E}
     -v_{E}\sin\varphi\,\be^{+}_{N}
     +v_{E}\cos\varphi\,\be^{+}_{U}\Big].
\end{align*}
$\sin\varphi>0$: $v_{N}>0\Rightarrow$ push along $+\be^{+}_{E}$,
$v_{E}>0\Rightarrow$ push along $-\be^{+}_{N}$; both to the right of the
motion, reversed for $\sin\varphi<0$.

Earth: $\omega=2\pi/86164\,\mathrm{s}=7.29\times10^{-5}\,\mathrm{s}^{-1}$,
$R=6.37\times10^{6}\,\mathrm{m}$, $g=9.81\,\mathrm{m\,s^{-2}}$:
\begin{align*}
   \omega^{2}R&=(7.29\times10^{-5})^{2}\times6.37\times10^{6}
     =5.32\times10^{-9}\times6.37\times10^{6}\\
     &=3.39\times10^{-2}\ \mathrm{m\,s^{-2}}=0.0035\,g,\\
   2\omega|\bv^{+}|&=1.46\times10^{-4}\ \mathrm{m\,s^{-2}}
     \quad\text{at }|\bv^{+}|=1\ \mathrm{m\,s^{-1}},\\
   \tan\delta&=\frac{\eps\sin\varphi\cos\varphi}{1-\eps\cos^{2}\varphi}
   =\frac{\tfrac12\eps\sin2\varphi}{1-\eps\cos^{2}\varphi},
   \qquad \eps=\frac{\omega^{2}R}{g}=0.00345,\\
   \frac{\dl}{\dl\varphi}\tan\delta&=0
   \ \Leftrightarrow\ \eps\cos2\varphi+O(\eps^{2})=0
   \ \Leftrightarrow\ \varphi=45^{\circ}+O(\eps),\\
   \tan\delta_{\max}&=\frac{\eps/2}{1-\eps/2}=1.73\times10^{-3},
   \qquad \delta_{\max}=0.099^{\circ}=5.9' .
\end{align*}
Negligible against $g$ in a laboratory; decisive in the ocean and the
atmosphere, where the comparison is with the residual accelerations
left after gravity is balanced by pressure.
\end{solution}

\begin{figure}[htbp]\centering
\renewcommand{\thefigure}{J}
% ---------------------------------------------------------------------
%  GENERATED by fig_ch7_planet.py -- do not edit; edit the script.
%  Figure J of Chapter 7: the geometry of Problem 3.
%  Panels (a),(b): latitude 35 deg; panel (b) draws omega^2 R / g = 0.55
%  (Earth: 0.0035) so that the tilt of the plumb line is visible.
% ---------------------------------------------------------------------
\begin{tikzpicture}[
   ar/.style={-{Stealth[length=2.4mm]},thick},
   sm/.style={-{Stealth[length=1.9mm]},semithick},
   guide/.style={gray!65,densely dashed,thin}]
\pgfplotsset{
   panel/.style={axis lines=none, x=1.20cm, y=1.20cm, clip=false,
                 every axis plot/.append style={no markers, line join=round}}}

% ======================= (a) the planet =======================
\begin{axis}[panel, name=pa, xmin=-2.550, xmax=3.550, ymin=-3.350, ymax=3.850]
% the planet in the meridional section through the station, and its equator
\addplot[thick,gray!55,fill=blue!4,domain=0:360,samples=181,smooth]
   ({2.000*cos(x)},{2.000*sin(x)});
\addplot[guide] coordinates {(-2.000,0) (2.000,0)};
% the axis  e_3  and the spin about it
\addplot[guide] coordinates {(0.000,-2.550) (0.000,2.400)};
\addplot[ar] coordinates {(0.000,2.400) (0.000,3.150)};
\addplot[{Stealth[length=1.9mm]}-,semithick,red!65!black,domain=-15:200,samples=40]
   ({0.620*cos(x)},{2.620+0.220*sin(x)});
\addplot[only marks,mark=*,mark size=1.4pt,black] coordinates {(0,0)};
% the equatorial radial direction at the station's longitude, and the latitude
\addplot[ar,gray!70] coordinates {(0.000,0.000) (1.300,0.000)};
\addplot[guide,domain=0:35.000,samples=20]
   ({0.720*cos(x)},{0.720*sin(x)});
% the station: distance R from the centre, R cos(phi) from the axis
\addplot[thick] coordinates {(0.000,0.000) (1.515,1.061)};
\addplot[guide] coordinates {(0.000,1.147) (1.488,1.147)};
% the local triad:  f_U outward, f_N along the meridian, f_E into the page
\addplot[ar,blue!55!black] coordinates {(1.761,1.233) (2.498,1.749)};
\addplot[ar,blue!55!black] coordinates {(1.552,1.270) (1.151,1.843)};
\addplot[semithick,gray!75,domain=0:360,samples=37] ({1.638+0.150*cos(x)},{1.147+0.150*sin(x)});
\addplot[semithick,gray!75] coordinates {(1.532,1.041) (1.744,1.253)};
\addplot[semithick,gray!75] coordinates {(1.532,1.253) (1.744,1.041)};
\node[font=\small, anchor=south west, inner sep=1pt] at (axis cs:0.100,3.200) {$\be_{3}$};
\node[font=\small, anchor=west, inner sep=1pt, text=red!65!black] at (axis cs:0.780,2.800) {$\omega$};
\node[font=\small, anchor=north, inner sep=1pt, text=gray!85] at (axis cs:0.650,-0.120) {$\bff_{\rho}$};
\node[font=\small, anchor=west, inner sep=1pt] at (axis cs:1.050,0.330) {$\varphi$};
\node[font=\small, anchor=south east, inner sep=1pt] at (axis cs:0.668,0.615) {$R$};
\node[font=\footnotesize, anchor=south, inner sep=1pt, text=gray!85] at (axis cs:0.530,1.237) {$R\cos\varphi$};
\node[font=\small, anchor=west, inner sep=1pt, text=blue!55!black] at (axis cs:2.564,1.795) {$\bff_{U}$};
\node[font=\small, anchor=south, inner sep=1pt, text=blue!55!black] at (axis cs:1.225,1.909) {$\bff_{N}$};
\node[font=\small, anchor=west, inner sep=1pt, text=gray!85] at (axis cs:2.081,1.069) {$\bff_{E}$};
\node[font=\small, anchor=north, inner sep=1pt] at (axis cs:0.500,-2.780) {(a) the local triad at the station};
\end{axis}

% ======================= (b) the station ======================
\begin{axis}[panel, name=pb, xmin=-2.500, xmax=2.600, ymin=-3.350, ymax=3.850, at={(pa.east)}, anchor=west, xshift=4mm]
% the horizon, with the ground hatched beneath it
\addplot[thick,gray!60] coordinates {(-2.250,0.000) (2.150,0.000)};
\pgfplotsinvokeforeach{-2.050,-1.640,-1.230,-0.820,-0.410,0.000,0.410,0.820,1.230,1.640,2.050}{
   \addplot[gray!45,thin] coordinates {(#1,0) (#1-0.16,-0.20)};}
% the observer's fixed triad:  north to the left, up, east into the page
\addplot[ar] coordinates {(-0.150,0.000) (-1.650,0.000)};
\addplot[ar] coordinates {(-0.000,0.150) (-0.000,1.650)};
\addplot[semithick,gray!75,domain=0:360,samples=37] ({0.000+0.150*cos(x)},{0.000+0.150*sin(x)});
\addplot[semithick,gray!75] coordinates {(-0.106,-0.106) (0.106,0.106)};
\addplot[semithick,gray!75] coordinates {(-0.106,0.106) (0.106,-0.106)};
% the angular velocity, at elevation phi above the northern horizon
\addplot[ar,red!65!black] coordinates {(-0.123,0.086) (-1.679,1.176)};
\addplot[guide,domain=145.000:180,samples=20]
   ({0.900*cos(x)},{0.900*sin(x)});
% gravity, the centrifugal term (perpendicular to the axis) and their sum
\addplot[ar,gray!70] coordinates {(-0.000,-0.150) (-0.000,-1.550)};
\addplot[sm,red!65!black] coordinates {(-0.000,-1.550) (0.401,-0.978)};
\addplot[ar,blue!55!black] coordinates {(0.057,-0.139) (0.401,-0.978)};
\addplot[guide,domain=-90:-67.728,samples=12]
   ({0.700*cos(x)},{0.700*sin(x)});
\node[font=\small, anchor=north, inner sep=1pt] at (axis cs:-1.620,-0.100) {$\be^{+}_{N}$};
\node[font=\small, anchor=west, inner sep=1pt] at (axis cs:0.100,1.600) {$\be^{+}_{U}$};
\node[font=\small, anchor=west, inner sep=1pt, text=gray!85] at (axis cs:0.220,0.340) {$\be^{+}_{E}$};
\node[font=\small, anchor=south east, inner sep=1pt, text=red!65!black] at (axis cs:-1.828,1.295) {$\bom$};
\node[font=\small, anchor=east, inner sep=1pt] at (axis cs:-1.049,0.331) {$\varphi$};
\node[font=\small, anchor=east, inner sep=1pt] at (axis cs:-0.100,-1.250) {$-g\,\be^{+}_{U}$};
\node[font=\footnotesize, anchor=west, inner sep=1pt, text=red!65!black] at (axis cs:0.550,-1.820) {$\omega^{2}R\cos\varphi$};
\node[font=\small, anchor=west, inner sep=1pt, text=blue!55!black] at (axis cs:0.550,-1.160) {$\bg_{\mathrm{eff}}$};
\node[font=\small, anchor=west, inner sep=1pt, text=blue!55!black] at (axis cs:0.450,-0.536) {$\delta$};
\node[font=\small, anchor=north, inner sep=1pt] at (axis cs:-0.050,-2.780) {(b) as $O^{+}$ draws it};
\end{axis}

% ============ (c) the centrifugal term against latitude ===========
%  Earth: omega^2 R = 33.9 mm/s^2; the plumb line tilts most, by
%  5.9 arcmin, at phi = 45.0 deg  (quoted in the caption).
\begin{axis}[name=pc, at={(pa.south west)}, anchor=outer north west,
   xshift=0.2mm, yshift=-3mm,
   width=13.8cm, height=4.9cm,
   xmin=0, xmax=90, ymin=0, ymax=1.05,
   xtick={0,15,...,90}, ytick={0,0.25,0.5,0.75,1},
   xticklabel={$\pgfmathprintnumber{\tick}^{\circ}$},
   yticklabel={$\pgfmathprintnumber{\tick}$},
   tick label style={font=\footnotesize},
   xlabel={(c) the centrifugal term against the latitude $\varphi$},
   ylabel={in units of $\omega^{2}R$},
   label style={font=\small},
   axis lines=left, axis line style={thin,gray!70},
   grid=major, grid style={gray!20, thin},
   legend style={font=\footnotesize, draw=none, fill=none,
                 at={(0.5,1.02)}, anchor=south, legend columns=-1,
                 /tikz/every even column/.append style={column sep=0.45cm}},
   legend cell align=left,
   every axis plot/.append style={no markers, line join=round, thick}]
\addplot[red!65!black]  table {
0.0 1.00000
0.5 0.99996
1.0 0.99985
1.5 0.99966
2.0 0.99939
2.5 0.99905
3.0 0.99863
3.5 0.99813
4.0 0.99756
4.5 0.99692
5.0 0.99619
5.5 0.99540
6.0 0.99452
6.5 0.99357
7.0 0.99255
7.5 0.99144
8.0 0.99027
8.5 0.98902
9.0 0.98769
9.5 0.98629
10.0 0.98481
10.5 0.98325
11.0 0.98163
11.5 0.97992
12.0 0.97815
12.5 0.97630
13.0 0.97437
13.5 0.97237
14.0 0.97030
14.5 0.96815
15.0 0.96593
15.5 0.96363
16.0 0.96126
16.5 0.95882
17.0 0.95630
17.5 0.95372
18.0 0.95106
18.5 0.94832
19.0 0.94552
19.5 0.94264
20.0 0.93969
20.5 0.93667
21.0 0.93358
21.5 0.93042
22.0 0.92718
22.5 0.92388
23.0 0.92050
23.5 0.91706
24.0 0.91355
24.5 0.90996
25.0 0.90631
25.5 0.90259
26.0 0.89879
26.5 0.89493
27.0 0.89101
27.5 0.88701
28.0 0.88295
28.5 0.87882
29.0 0.87462
29.5 0.87036
30.0 0.86603
30.5 0.86163
31.0 0.85717
31.5 0.85264
32.0 0.84805
32.5 0.84339
33.0 0.83867
33.5 0.83389
34.0 0.82904
34.5 0.82413
35.0 0.81915
35.5 0.81412
36.0 0.80902
36.5 0.80386
37.0 0.79864
37.5 0.79335
38.0 0.78801
38.5 0.78261
39.0 0.77715
39.5 0.77162
40.0 0.76604
40.5 0.76041
41.0 0.75471
41.5 0.74896
42.0 0.74314
42.5 0.73728
43.0 0.73135
43.5 0.72537
44.0 0.71934
44.5 0.71325
45.0 0.70711
45.5 0.70091
46.0 0.69466
46.5 0.68835
47.0 0.68200
47.5 0.67559
48.0 0.66913
48.5 0.66262
49.0 0.65606
49.5 0.64945
50.0 0.64279
50.5 0.63608
51.0 0.62932
51.5 0.62251
52.0 0.61566
52.5 0.60876
53.0 0.60182
53.5 0.59482
54.0 0.58779
54.5 0.58070
55.0 0.57358
55.5 0.56641
56.0 0.55919
56.5 0.55194
57.0 0.54464
57.5 0.53730
58.0 0.52992
58.5 0.52250
59.0 0.51504
59.5 0.50754
60.0 0.50000
60.5 0.49242
61.0 0.48481
61.5 0.47716
62.0 0.46947
62.5 0.46175
63.0 0.45399
63.5 0.44620
64.0 0.43837
64.5 0.43051
65.0 0.42262
65.5 0.41469
66.0 0.40674
66.5 0.39875
67.0 0.39073
67.5 0.38268
68.0 0.37461
68.5 0.36650
69.0 0.35837
69.5 0.35021
70.0 0.34202
70.5 0.33381
71.0 0.32557
71.5 0.31730
72.0 0.30902
72.5 0.30071
73.0 0.29237
73.5 0.28402
74.0 0.27564
74.5 0.26724
75.0 0.25882
75.5 0.25038
76.0 0.24192
76.5 0.23345
77.0 0.22495
77.5 0.21644
78.0 0.20791
78.5 0.19937
79.0 0.19081
79.5 0.18224
80.0 0.17365
80.5 0.16505
81.0 0.15643
81.5 0.14781
82.0 0.13917
82.5 0.13053
83.0 0.12187
83.5 0.11320
84.0 0.10453
84.5 0.09585
85.0 0.08716
85.5 0.07846
86.0 0.06976
86.5 0.06105
87.0 0.05234
87.5 0.04362
88.0 0.03490
88.5 0.02618
89.0 0.01745
89.5 0.00873
90.0 0.00000
};
\addlegendentry{$\cos\varphi$, the magnitude}
\addplot[blue!55!black] table {
0.0 1.00000
0.5 0.99992
1.0 0.99970
1.5 0.99931
2.0 0.99878
2.5 0.99810
3.0 0.99726
3.5 0.99627
4.0 0.99513
4.5 0.99384
5.0 0.99240
5.5 0.99081
6.0 0.98907
6.5 0.98719
7.0 0.98515
7.5 0.98296
8.0 0.98063
8.5 0.97815
9.0 0.97553
9.5 0.97276
10.0 0.96985
10.5 0.96679
11.0 0.96359
11.5 0.96025
12.0 0.95677
12.5 0.95315
13.0 0.94940
13.5 0.94550
14.0 0.94147
14.5 0.93731
15.0 0.93301
15.5 0.92858
16.0 0.92402
16.5 0.91934
17.0 0.91452
17.5 0.90958
18.0 0.90451
18.5 0.89932
19.0 0.89401
19.5 0.88857
20.0 0.88302
20.5 0.87735
21.0 0.87157
21.5 0.86568
22.0 0.85967
22.5 0.85355
23.0 0.84733
23.5 0.84100
24.0 0.83457
24.5 0.82803
25.0 0.82139
25.5 0.81466
26.0 0.80783
26.5 0.80091
27.0 0.79389
27.5 0.78679
28.0 0.77960
28.5 0.77232
29.0 0.76496
29.5 0.75752
30.0 0.75000
30.5 0.74240
31.0 0.73474
31.5 0.72700
32.0 0.71919
32.5 0.71131
33.0 0.70337
33.5 0.69537
34.0 0.68730
34.5 0.67918
35.0 0.67101
35.5 0.66278
36.0 0.65451
36.5 0.64619
37.0 0.63782
37.5 0.62941
38.0 0.62096
38.5 0.61248
39.0 0.60396
39.5 0.59540
40.0 0.58682
40.5 0.57822
41.0 0.56959
41.5 0.56093
42.0 0.55226
42.5 0.54358
43.0 0.53488
43.5 0.52617
44.0 0.51745
44.5 0.50873
45.0 0.50000
45.5 0.49127
46.0 0.48255
46.5 0.47383
47.0 0.46512
47.5 0.45642
48.0 0.44774
48.5 0.43907
49.0 0.43041
49.5 0.42178
50.0 0.41318
50.5 0.40460
51.0 0.39604
51.5 0.38752
52.0 0.37904
52.5 0.37059
53.0 0.36218
53.5 0.35381
54.0 0.34549
54.5 0.33722
55.0 0.32899
55.5 0.32082
56.0 0.31270
56.5 0.30463
57.0 0.29663
57.5 0.28869
58.0 0.28081
58.5 0.27300
59.0 0.26526
59.5 0.25760
60.0 0.25000
60.5 0.24248
61.0 0.23504
61.5 0.22768
62.0 0.22040
62.5 0.21321
63.0 0.20611
63.5 0.19909
64.0 0.19217
64.5 0.18534
65.0 0.17861
65.5 0.17197
66.0 0.16543
66.5 0.15900
67.0 0.15267
67.5 0.14645
68.0 0.14033
68.5 0.13432
69.0 0.12843
69.5 0.12265
70.0 0.11698
70.5 0.11143
71.0 0.10599
71.5 0.10068
72.0 0.09549
72.5 0.09042
73.0 0.08548
73.5 0.08066
74.0 0.07598
74.5 0.07142
75.0 0.06699
75.5 0.06269
76.0 0.05853
76.5 0.05450
77.0 0.05060
77.5 0.04685
78.0 0.04323
78.5 0.03975
79.0 0.03641
79.5 0.03321
80.0 0.03015
80.5 0.02724
81.0 0.02447
81.5 0.02185
82.0 0.01937
82.5 0.01704
83.0 0.01485
83.5 0.01281
84.0 0.01093
84.5 0.00919
85.0 0.00760
85.5 0.00616
86.0 0.00487
86.5 0.00373
87.0 0.00274
87.5 0.00190
88.0 0.00122
88.5 0.00069
89.0 0.00030
89.5 0.00008
90.0 0.00000
};
\addlegendentry{$\cos^{2}\!\varphi$, along $\be^{+}_{U}$}
\addplot[blue!55!black,densely dashed] table {
0.0 0.00000
0.5 0.00873
1.0 0.01745
1.5 0.02617
2.0 0.03488
2.5 0.04358
3.0 0.05226
3.5 0.06093
4.0 0.06959
4.5 0.07822
5.0 0.08682
5.5 0.09540
6.0 0.10396
6.5 0.11248
7.0 0.12096
7.5 0.12941
8.0 0.13782
8.5 0.14619
9.0 0.15451
9.5 0.16278
10.0 0.17101
10.5 0.17918
11.0 0.18730
11.5 0.19537
12.0 0.20337
12.5 0.21131
13.0 0.21919
13.5 0.22700
14.0 0.23474
14.5 0.24240
15.0 0.25000
15.5 0.25752
16.0 0.26496
16.5 0.27232
17.0 0.27960
17.5 0.28679
18.0 0.29389
18.5 0.30091
19.0 0.30783
19.5 0.31466
20.0 0.32139
20.5 0.32803
21.0 0.33457
21.5 0.34100
22.0 0.34733
22.5 0.35355
23.0 0.35967
23.5 0.36568
24.0 0.37157
24.5 0.37735
25.0 0.38302
25.5 0.38857
26.0 0.39401
26.5 0.39932
27.0 0.40451
27.5 0.40958
28.0 0.41452
28.5 0.41934
29.0 0.42402
29.5 0.42858
30.0 0.43301
30.5 0.43731
31.0 0.44147
31.5 0.44550
32.0 0.44940
32.5 0.45315
33.0 0.45677
33.5 0.46025
34.0 0.46359
34.5 0.46679
35.0 0.46985
35.5 0.47276
36.0 0.47553
36.5 0.47815
37.0 0.48063
37.5 0.48296
38.0 0.48515
38.5 0.48719
39.0 0.48907
39.5 0.49081
40.0 0.49240
40.5 0.49384
41.0 0.49513
41.5 0.49627
42.0 0.49726
42.5 0.49810
43.0 0.49878
43.5 0.49931
44.0 0.49970
44.5 0.49992
45.0 0.50000
45.5 0.49992
46.0 0.49970
46.5 0.49931
47.0 0.49878
47.5 0.49810
48.0 0.49726
48.5 0.49627
49.0 0.49513
49.5 0.49384
50.0 0.49240
50.5 0.49081
51.0 0.48907
51.5 0.48719
52.0 0.48515
52.5 0.48296
53.0 0.48063
53.5 0.47815
54.0 0.47553
54.5 0.47276
55.0 0.46985
55.5 0.46679
56.0 0.46359
56.5 0.46025
57.0 0.45677
57.5 0.45315
58.0 0.44940
58.5 0.44550
59.0 0.44147
59.5 0.43731
60.0 0.43301
60.5 0.42858
61.0 0.42402
61.5 0.41934
62.0 0.41452
62.5 0.40958
63.0 0.40451
63.5 0.39932
64.0 0.39401
64.5 0.38857
65.0 0.38302
65.5 0.37735
66.0 0.37157
66.5 0.36568
67.0 0.35967
67.5 0.35355
68.0 0.34733
68.5 0.34100
69.0 0.33457
69.5 0.32803
70.0 0.32139
70.5 0.31466
71.0 0.30783
71.5 0.30091
72.0 0.29389
72.5 0.28679
73.0 0.27960
73.5 0.27232
74.0 0.26496
74.5 0.25752
75.0 0.25000
75.5 0.24240
76.0 0.23474
76.5 0.22700
77.0 0.21919
77.5 0.21131
78.0 0.20337
78.5 0.19537
79.0 0.18730
79.5 0.17918
80.0 0.17101
80.5 0.16278
81.0 0.15451
81.5 0.14619
82.0 0.13782
82.5 0.12941
83.0 0.12096
83.5 0.11248
84.0 0.10396
84.5 0.09540
85.0 0.08682
85.5 0.07822
86.0 0.06959
86.5 0.06093
87.0 0.05226
87.5 0.04358
88.0 0.03488
88.5 0.02617
89.0 0.01745
89.5 0.00873
90.0 0.00000
};
\addlegendentry{$\sin\varphi\cos\varphi$, along $-\be^{+}_{N}$}
\addplot[guide,forget plot] coordinates {(45,0) (45,0.5)};
\addplot[only marks,mark=*,mark size=1.6pt,blue!55!black,forget plot]
   coordinates {(45,0.5)};
\end{axis}
\end{tikzpicture}

\caption{The geometry of Problem 3 at latitude $\varphi=35^{\circ}$.
(a) The meridional section through the station: the planet turns at
$\omega$ about $\be_{3}$, $\bff_{\rho}$ is the equatorial radial
direction at the station's longitude, and
$\bff_{U}=\cos\varphi\,\bff_{\rho}+\sin\varphi\,\be_{3}$,
$\bff_{N}=-\sin\varphi\,\bff_{\rho}+\cos\varphi\,\be_{3}$, with
$\bff_{E}$ into the page, is the local triad; the station is $R$
from the centre and $R\cos\varphi$ from the axis, and that second
distance is the whole of the centrifugal term. (b) The same picture as
$O^{+}$ draws it, $\be^{+}_{U}$ vertical, north to the left, east into
the page: the angular velocity $\bom$ points at the celestial pole, at elevation
$\varphi$ above the northern horizon, and the shifter triad of
Remark~\ref{rem:ch7whatinertial} turns backwards about it once a
sidereal day; the centrifugal term, perpendicular to the axis,
tilts the effective gravity $\bg_{\mathrm{eff}}$ from $-g\be^{+}_{U}$
towards the equator by $\delta$, so a plumb line does not point at the
centre of the planet ($\omega^{2}R/g$ is drawn as $0.55$; for the Earth
it is $0.0035$). (c) Magnitude and components of the centrifugal term
against latitude.}
\label{fig:ch7planet}
\end{figure}
\begin{figure}[htbp]\centering
\renewcommand{\thefigure}{K}
% ---------------------------------------------------------------------
%  GENERATED by fig_ch7_planet.py -- do not edit; edit the script.
%  Figure K of Chapter 7: the geometry of Problem 3 in three dimensions,
%  longitude 40 deg, latitude 35 deg, R = 1.  Explicit unit vectors
%  (azimuth 100, elevation 22); curves on the far side are dotted.
% ---------------------------------------------------------------------
\begin{tikzpicture}[
   ar/.style={-{Stealth[length=2.4mm]},thick},
   sm/.style={-{Stealth[length=1.9mm]},semithick},
   guide/.style={gray!65,densely dashed,thin},
   back/.style={gray!45,dotted,thin}]
\begin{axis}[axis lines=none, clip=false,
   x={(-0.729cm,-1.549cm)}, y={(4.136cm,-0.273cm)}, z={(0cm,3.894cm)},
   xmin=-1.6, xmax=1.6, ymin=-1.6, ymax=1.6, zmin=-1.2, zmax=1.6,
   every axis plot/.append style={no markers, line join=round}]
% the sphere, as a translucent surface
\addplot3[surf, shader=flat, draw=gray!35, line width=0.15pt, fill=blue!25,
   opacity=0.13, z buffer=sort, samples=37, samples y=19,
   domain=0:360, y domain=-90:90]
   ({cos(x)*cos(y)},{sin(x)*cos(y)},{sin(y)});

% the equator, the meridian of the station and its parallel
\addplot3[thin,gray!70] coordinates {(1.000,0.000,0.000) (0.999,0.035,0.000) (0.998,0.070,0.000) (0.995,0.105,0.000) (0.990,0.139,0.000) (0.985,0.174,0.000) (0.978,0.208,0.000) (0.970,0.242,0.000) (0.961,0.276,0.000) (0.951,0.309,0.000) (0.940,0.342,0.000) (0.927,0.375,0.000) (0.914,0.407,0.000) (0.899,0.438,0.000) (0.883,0.469,0.000) (0.866,0.500,0.000) (0.848,0.530,0.000) (0.829,0.559,0.000) (0.809,0.588,0.000) (0.788,0.616,0.000) (0.766,0.643,0.000) (0.743,0.669,0.000) (0.719,0.695,0.000) (0.695,0.719,0.000) (0.669,0.743,0.000) (0.643,0.766,0.000) (0.616,0.788,0.000) (0.588,0.809,0.000) (0.559,0.829,0.000) (0.530,0.848,0.000) (0.500,0.866,0.000) (0.469,0.883,0.000) (0.438,0.899,0.000) (0.407,0.914,0.000) (0.375,0.927,0.000) (0.342,0.940,0.000) (0.309,0.951,0.000) (0.276,0.961,0.000) (0.242,0.970,0.000) (0.208,0.978,0.000) (0.174,0.985,0.000) (0.139,0.990,0.000) (0.105,0.995,0.000) (0.070,0.998,0.000) (0.035,0.999,0.000) (0.000,1.000,0.000) (-0.035,0.999,0.000) (-0.070,0.998,0.000) (-0.105,0.995,0.000) (-0.139,0.990,0.000) (-0.174,0.985,0.000) (-0.208,0.978,0.000)};
\addplot3[thin,gray!70] coordinates {(0.208,-0.978,0.000) (0.242,-0.970,0.000) (0.276,-0.961,0.000) (0.309,-0.951,0.000) (0.342,-0.940,0.000) (0.375,-0.927,0.000) (0.407,-0.914,0.000) (0.438,-0.899,0.000) (0.469,-0.883,0.000) (0.500,-0.866,0.000) (0.530,-0.848,0.000) (0.559,-0.829,0.000) (0.588,-0.809,0.000) (0.616,-0.788,0.000) (0.643,-0.766,0.000) (0.669,-0.743,0.000) (0.695,-0.719,0.000) (0.719,-0.695,0.000) (0.743,-0.669,0.000) (0.766,-0.643,0.000) (0.788,-0.616,0.000) (0.809,-0.588,0.000) (0.829,-0.559,0.000) (0.848,-0.530,0.000) (0.866,-0.500,0.000) (0.883,-0.469,0.000) (0.899,-0.438,0.000) (0.914,-0.407,0.000) (0.927,-0.375,0.000) (0.940,-0.342,0.000) (0.951,-0.309,0.000) (0.961,-0.276,0.000) (0.970,-0.242,0.000) (0.978,-0.208,0.000) (0.985,-0.174,0.000) (0.990,-0.139,0.000) (0.995,-0.105,0.000) (0.998,-0.070,0.000) (0.999,-0.035,0.000) (1.000,-0.000,0.000)};
\addplot3[thin,gray!70] coordinates {(0.766,0.643,0.000) (0.766,0.642,0.035) (0.764,0.641,0.070) (0.762,0.639,0.105) (0.759,0.637,0.139) (0.754,0.633,0.174) (0.749,0.629,0.208) (0.743,0.624,0.242) (0.736,0.618,0.276) (0.729,0.611,0.309) (0.720,0.604,0.342) (0.710,0.596,0.375) (0.700,0.587,0.407) (0.689,0.578,0.438) (0.676,0.568,0.469) (0.663,0.557,0.500) (0.650,0.545,0.530) (0.635,0.533,0.559) (0.620,0.520,0.588) (0.604,0.507,0.616) (0.587,0.492,0.643) (0.569,0.478,0.669) (0.551,0.462,0.695) (0.532,0.447,0.719) (0.513,0.430,0.743) (0.492,0.413,0.766) (0.472,0.396,0.788) (0.450,0.378,0.809) (0.428,0.359,0.829) (0.406,0.341,0.848) (0.383,0.321,0.866) (0.360,0.302,0.883) (0.336,0.282,0.899) (0.312,0.261,0.914) (0.287,0.241,0.927) (0.262,0.220,0.940) (0.237,0.199,0.951) (0.211,0.177,0.961) (0.185,0.156,0.970) (0.159,0.134,0.978) (0.133,0.112,0.985) (0.107,0.089,0.990) (0.080,0.067,0.995) (0.053,0.045,0.998) (0.027,0.022,0.999) (0.000,0.000,1.000) (-0.027,-0.022,0.999) (-0.053,-0.045,0.998) (-0.080,-0.067,0.995) (-0.107,-0.089,0.990) (-0.133,-0.112,0.985) (-0.159,-0.134,0.978) (-0.185,-0.156,0.970) (-0.211,-0.177,0.961) (-0.237,-0.199,0.951) (-0.262,-0.220,0.940) (-0.287,-0.241,0.927) (-0.312,-0.261,0.914) (-0.336,-0.282,0.899)};
\addplot3[thin,gray!70] coordinates {(0.336,0.282,-0.899) (0.360,0.302,-0.883) (0.383,0.321,-0.866) (0.406,0.341,-0.848) (0.428,0.359,-0.829) (0.450,0.378,-0.809) (0.472,0.396,-0.788) (0.492,0.413,-0.766) (0.513,0.430,-0.743) (0.532,0.447,-0.719) (0.551,0.462,-0.695) (0.569,0.478,-0.669) (0.587,0.492,-0.643) (0.604,0.507,-0.616) (0.620,0.520,-0.588) (0.635,0.533,-0.559) (0.650,0.545,-0.530) (0.663,0.557,-0.500) (0.676,0.568,-0.469) (0.689,0.578,-0.438) (0.700,0.587,-0.407) (0.710,0.596,-0.375) (0.720,0.604,-0.342) (0.729,0.611,-0.309) (0.736,0.618,-0.276) (0.743,0.624,-0.242) (0.749,0.629,-0.208) (0.754,0.633,-0.174) (0.759,0.637,-0.139) (0.762,0.639,-0.105) (0.764,0.641,-0.070) (0.766,0.642,-0.035) (0.766,0.643,-0.000)};
\addplot3[thin,gray!70] coordinates {(0.819,0.000,0.574) (0.819,0.029,0.574) (0.817,0.057,0.574) (0.815,0.086,0.574) (0.811,0.114,0.574) (0.807,0.142,0.574) (0.801,0.170,0.574) (0.795,0.198,0.574) (0.787,0.226,0.574) (0.779,0.253,0.574) (0.770,0.280,0.574) (0.760,0.307,0.574) (0.748,0.333,0.574) (0.736,0.359,0.574) (0.723,0.385,0.574) (0.709,0.410,0.574) (0.695,0.434,0.574) (0.679,0.458,0.574) (0.663,0.481,0.574) (0.646,0.504,0.574) (0.628,0.527,0.574) (0.609,0.548,0.574) (0.589,0.569,0.574) (0.569,0.589,0.574) (0.548,0.609,0.574) (0.527,0.628,0.574) (0.504,0.646,0.574) (0.481,0.663,0.574) (0.458,0.679,0.574) (0.434,0.695,0.574) (0.410,0.709,0.574) (0.385,0.723,0.574) (0.359,0.736,0.574) (0.333,0.748,0.574) (0.307,0.760,0.574) (0.280,0.770,0.574) (0.253,0.779,0.574) (0.226,0.787,0.574) (0.198,0.795,0.574) (0.170,0.801,0.574) (0.142,0.807,0.574) (0.114,0.811,0.574) (0.086,0.815,0.574) (0.057,0.817,0.574) (0.029,0.819,0.574) (0.000,0.819,0.574) (-0.029,0.819,0.574) (-0.057,0.817,0.574) (-0.086,0.815,0.574) (-0.114,0.811,0.574) (-0.142,0.807,0.574) (-0.170,0.801,0.574) (-0.198,0.795,0.574) (-0.226,0.787,0.574) (-0.253,0.779,0.574) (-0.280,0.770,0.574) (-0.307,0.760,0.574) (-0.333,0.748,0.574) (-0.359,0.736,0.574) (-0.385,0.723,0.574)};
\addplot3[thin,gray!70] coordinates {(-0.086,-0.815,0.574) (-0.057,-0.817,0.574) (-0.029,-0.819,0.574) (-0.000,-0.819,0.574) (0.029,-0.819,0.574) (0.057,-0.817,0.574) (0.086,-0.815,0.574) (0.114,-0.811,0.574) (0.142,-0.807,0.574) (0.170,-0.801,0.574) (0.198,-0.795,0.574) (0.226,-0.787,0.574) (0.253,-0.779,0.574) (0.280,-0.770,0.574) (0.307,-0.760,0.574) (0.333,-0.748,0.574) (0.359,-0.736,0.574) (0.385,-0.723,0.574) (0.410,-0.709,0.574) (0.434,-0.695,0.574) (0.458,-0.679,0.574) (0.481,-0.663,0.574) (0.504,-0.646,0.574) (0.527,-0.628,0.574) (0.548,-0.609,0.574) (0.569,-0.589,0.574) (0.589,-0.569,0.574) (0.609,-0.548,0.574) (0.628,-0.527,0.574) (0.646,-0.504,0.574) (0.663,-0.481,0.574) (0.679,-0.458,0.574) (0.695,-0.434,0.574) (0.709,-0.410,0.574) (0.723,-0.385,0.574) (0.736,-0.359,0.574) (0.748,-0.333,0.574) (0.760,-0.307,0.574) (0.770,-0.280,0.574) (0.779,-0.253,0.574) (0.787,-0.226,0.574) (0.795,-0.198,0.574) (0.801,-0.170,0.574) (0.807,-0.142,0.574) (0.811,-0.114,0.574) (0.815,-0.086,0.574) (0.817,-0.057,0.574) (0.819,-0.029,0.574) (0.819,-0.000,0.574)};
% the fixed triad and the axis, the radial direction, the angles
\addplot3[guide] coordinates {(-0.000,-0.000,-1.000) (0.000,0.000,0.000)};
\addplot3[ar] coordinates {(0.000,0.000,0.000) (1.450,0.000,0.000)};
\addplot3[ar] coordinates {(0.000,0.000,0.000) (0.000,1.450,0.000)};
\addplot3[ar] coordinates {(0.000,0.000,0.000) (0.000,0.000,1.500)};
\addplot3[ar,gray!70] coordinates {(0.000,0.000,0.000) (0.766,0.643,0.000)};
\addplot3[guide] coordinates {(0.420,0.000,0.000) (0.420,0.012,0.000) (0.419,0.024,0.000) (0.418,0.037,0.000) (0.417,0.049,0.000) (0.416,0.061,0.000) (0.414,0.073,0.000) (0.411,0.085,0.000) (0.409,0.097,0.000) (0.406,0.109,0.000) (0.402,0.120,0.000) (0.399,0.132,0.000) (0.395,0.144,0.000) (0.390,0.155,0.000) (0.386,0.166,0.000) (0.381,0.177,0.000) (0.375,0.188,0.000) (0.370,0.199,0.000) (0.364,0.210,0.000) (0.357,0.220,0.000) (0.351,0.231,0.000) (0.344,0.241,0.000) (0.337,0.251,0.000) (0.329,0.260,0.000) (0.322,0.270,0.000)};
\addplot3[guide] coordinates {(0.245,0.206,0.000) (0.245,0.206,0.008) (0.245,0.205,0.016) (0.244,0.205,0.024) (0.244,0.205,0.033) (0.243,0.204,0.041) (0.242,0.203,0.049) (0.241,0.202,0.057) (0.240,0.201,0.065) (0.239,0.200,0.073) (0.237,0.199,0.081) (0.236,0.198,0.088) (0.234,0.196,0.096) (0.232,0.195,0.104) (0.230,0.193,0.112) (0.227,0.191,0.119) (0.225,0.189,0.127) (0.223,0.187,0.134) (0.220,0.184,0.142) (0.217,0.182,0.149) (0.214,0.180,0.156) (0.211,0.177,0.163) (0.208,0.174,0.170) (0.204,0.171,0.177) (0.201,0.168,0.184)};
% the station, its distances from the centre and from the axis
\addplot3[thick] coordinates {(0.000,0.000,0.000) (0.628,0.527,0.574)};
\addplot3[guide] coordinates {(0.000,0.000,0.574) (0.628,0.527,0.574)};
\addplot3[only marks,mark=*,mark size=1.6pt,black] coordinates {(0.000,0.000,0.000) (0.628,0.527,0.574)};
% the moving triad at the station
\addplot3[ar,blue!55!black] coordinates {(0.628,0.527,0.574) (0.338,0.871,0.574)};
\addplot3[ar,blue!55!black] coordinates {(0.628,0.527,0.574) (0.430,0.361,0.942)};
\addplot3[ar,blue!55!black] coordinates {(0.628,0.527,0.574) (0.910,0.763,0.832)};
% the spin about the axis
\addplot3[sm,red!65!black] coordinates {(-0.167,0.199,1.220) (-0.192,0.175,1.220) (-0.214,0.148,1.220) (-0.232,0.118,1.220) (-0.245,0.086,1.220) (-0.255,0.052,1.220) (-0.259,0.017,1.220) (-0.259,-0.017,1.220) (-0.255,-0.052,1.220) (-0.245,-0.086,1.220) (-0.232,-0.118,1.220) (-0.214,-0.148,1.220) (-0.192,-0.175,1.220) (-0.167,-0.199,1.220) (-0.139,-0.220,1.220) (-0.108,-0.236,1.220) (-0.076,-0.249,1.220) (-0.042,-0.257,1.220) (-0.007,-0.260,1.220) (0.028,-0.259,1.220) (0.062,-0.252,1.220) (0.095,-0.242,1.220) (0.127,-0.227,1.220) (0.156,-0.208,1.220) (0.183,-0.185,1.220) (0.206,-0.159,1.220) (0.225,-0.130,1.220) (0.241,-0.099,1.220) (0.252,-0.066,1.220) (0.258,-0.031,1.220) (0.260,0.003,1.220) (0.257,0.038,1.220) (0.250,0.072,1.220) (0.238,0.105,1.220) (0.222,0.136,1.220) (0.201,0.164,1.220) (0.178,0.190,1.220) (0.151,0.212,1.220) (0.121,0.230,1.220) (0.089,0.244,1.220)};
\node[font=\small, anchor=north east, inner sep=1pt, shift={(-0.13cm,-0.27cm)}] at (axis cs:1.450,0.000,0.000) {$\be_{1}$};
\node[font=\small, anchor=west, inner sep=1pt, shift={(0.30cm,-0.02cm)}] at (axis cs:0.000,1.450,0.000) {$\be_{2}$};
\node[font=\small, anchor=south, inner sep=1pt, shift={(0.00cm,0.30cm)}] at (axis cs:0.000,0.000,1.500) {$\be_{3}$};
\node[font=\small, anchor=north west, inner sep=1pt, text=gray!85, shift={(0.25cm,-0.16cm)}] at (axis cs:0.766,0.643,0.000) {$\bff_{\rho}$};
\node[font=\small, anchor=west, inner sep=1pt, text=blue!55!black, shift={(0.29cm,0.06cm)}] at (axis cs:0.338,0.871,0.574) {$\bff_{E}$};
\node[font=\small, anchor=west, inner sep=1pt, text=blue!55!black, shift={(0.45cm,0.15cm)}] at (axis cs:0.430,0.361,0.942) {$\bff_{N}$};
\node[font=\small, anchor=south, inner sep=1pt, text=blue!55!black, shift={(-0.15cm,0.35cm)}] at (axis cs:0.910,0.763,0.832) {$\bff_{U}$};
\node[font=\footnotesize, anchor=center, inner sep=1pt, shift={(0.00cm,0.00cm)}] at (axis cs:0.639,0.233,0.000) {$\theta$};
\node[font=\footnotesize, anchor=center, inner sep=1pt, shift={(1.64cm,-0.44cm)}] at (axis cs:0.000,0.000,0.000) {$\varphi$};
\node[font=\footnotesize, anchor=north west, inner sep=1pt, shift={(1.32cm,0.70cm)}] at (axis cs:0.000,0.000,0.000) {$R$};
\node[font=\footnotesize, anchor=east, inner sep=1pt, text=gray!85, shift={(-0.12cm,0.00cm)}] at (axis cs:0.000,0.000,0.574) {$R\cos\varphi$};
\node[font=\small, anchor=west, inner sep=1pt, text=red!65!black, shift={(1.30cm,0.00cm)}] at (axis cs:0.000,0.000,1.220) {$\omega$};
\node[font=\small, anchor=east, inner sep=1pt, shift={(-0.20cm,0.10cm)}] at (axis cs:0.000,0.000,0.000) {$O$};
\node[font=\footnotesize, anchor=north, inner sep=1pt, text=gray!85, shift={(0.90cm,-0.12cm)}] at (axis cs:0.985,0.174,0.000) {equator};
\node[font=\footnotesize, anchor=north east, inner sep=1pt, text=gray!85, shift={(-0.10cm,-0.15cm)}] at (axis cs:0.579,-0.579,0.574) {parallel $\varphi$};
\node[font=\footnotesize, anchor=north, inner sep=1pt, text=gray!85, shift={(0.30cm,-0.28cm)}] at (axis cs:0.336,0.282,-0.899) {meridian $\theta$};
\end{axis}
\end{tikzpicture}

\caption{Figure~\ref{fig:ch7planet}(a) in three dimensions: what is
where, for a station at longitude $\theta=40^{\circ}$ and latitude
$\varphi=35^{\circ}$. $O$ is the inertial observer at the centre with
the fixed triad $\{\be_{i}\}$, $\be_{3}$ the axis of the spin $\omega$.
The station lies on the meridian of longitude $\theta=\omega t+\theta_{0}$,
measured in the equatorial plane from $\be_{1}$, and on the parallel of
latitude $\varphi$, measured from the equator along that meridian;
$\bff_{\rho}$ is the radial direction in the equatorial plane at that
longitude. The station is $R$ from $O$ and $R\cos\varphi$ from the axis.
At the station, $\bff_{U}$ points straight up along $O$--station,
$\bff_{N}$ points north along the meridian and $\bff_{E}$ points east
along the parallel; only the near side of the three circles is drawn.}
\label{fig:ch7planet3d}
\end{figure}
